\input{preamble}

% OK, commençons ici.
%
\begin{document}

\title{Espaces de Quot et de Hilbert}

\maketitle

\phantomsection
\label{section-phantom}

\tableofcontents




\section{Introduction}
\label{section-introduction}

\noindent
Tel qu'il avait été conçu initialement, ce chapitre devait traiter des
foncteurs de Quot et de Hilbert et démontrer que ce sont des espaces
algébriques sous réserve que certaines conditions techniques soient
satisfaites. Dans le cadre des schémas, ce sujet est traité dans les exposés
de Grothendieck au séminaire Bourbaki ; voir
\cite{Gr-I},
\cite{Gr-II},
\cite{Gr-III},
\cite{Gr-IV},
\cite{Gr-V} et
\cite{Gr-VI}. Pour les schémas projectifs, les schémas de Quot et de Hilbert
se plongent dans des grassmanniennes d'espaces de sections de faisceaux
inversibles très amples convenables, ce qui fournit une méthode de
construction de ces schémas. Notre approche est différente : nous utilisons
les axiomes d'Artin pour démontrer que Quot et Hilb sont des espaces
algébriques.

\medskip\noindent
Après plus ample réflexion, il s'est révélé plus commode, pour le
développement de la théorie dans le projet Stacks, de commencer par le champ
$\Cohstack_{X/B}$ des faisceaux cohérents (à support propre sur la base),
introduit dans \cite{lieblich_remarks}. Pour nous, $f : X \to B$ est un
morphisme d'espaces algébriques satisfaisant des conditions techniques
convenables, bien que ceci puisse être généralisé (voir ci-dessous). Étant
donnés des modules $\mathcal{F}$ et $\mathcal{G}$ sur $X$, sous des
hypothèses convenables, le foncteur
$T/B \mapsto \Hom_{X_T}(\mathcal{F}_T, \mathcal{G}_T)$ est un espace
algébrique $\mathit{Hom}(\mathcal{F}, \mathcal{G})$ sur $B$. Voir la
section \ref{section-hom}. On montre à la section \ref{section-isom} que le
sous-foncteur $\mathit{Isom}(\mathcal{F}, \mathcal{G})$ des isomorphismes est
un espace algébrique. Ceci est utilisé dans les sections suivantes pour
montrer que la diagonale du champ $\Cohstack_{X/B}$ est représentable. Nous
démontrons que $\Cohstack_{X/B}$ est un champ algébrique dans la section
\ref{section-stack-coherent-sheaves} lorsque $X \to B$ est plat, et dans la
section \ref{section-not-flat} en général. On trouvera dans l'introduction de
cette section des indications bibliographiques.

\medskip\noindent
Cela établi, il est assez immédiat de démontrer que
$\Quotfunctor_{\mathcal{F}/X/B}$, $\Hilbfunctor_{X/B}$ et
$\Picardfunctor_{X/B}$ sont des espaces algébriques et que
$\Picardstack_{X/B}$ est un champ algébrique. Voir les sections
\ref{section-quot}, \ref{section-hilb},
\ref{section-picard-functor} et \ref{section-picard-stack}.

\medskip\noindent
De la manière habituelle, nous en déduisons à la section
\ref{section-relative-morphisms} que le foncteur
$\mathit{Mor}_B(Z, X)$ des morphismes relatifs est un espace algébrique (sous
des hypothèses convenables).

\medskip\noindent
Dans la section \ref{section-stack-of-spaces}, nous démontrons que le champ en
groupoïdes
$$
\Spacesstack'_{fp, flat, proper}
$$
qui paramètre les familles plates d'espaces algébriques propres satisfait à
tous les axiomes d'Artin (y compris l'ouverture de la versalité), sauf à
l'effectivité formelle. Nous avons choisi à dessein cette notation fort
malcommode pour ce champ, car le lecteur doit user de ses propriétés avec
prudence.

\medskip\noindent
Dans la section \ref{section-polarized}, nous démontrons que le champ
$\Polarizedstack$ qui paramètre les familles plates d'espaces algébriques
propres polarisés est un champ algébrique. Grâce à notre étude des familles
plates d'espaces algébriques propres, il suffit pour cela de démontrer
l'effectivité formelle pour les schémas polarisés, résultat souvent connu
sous le nom de théorème d'algébrisation de Grothendieck.

\medskip\noindent
Dans la section \ref{section-curves}, nous démontrons que le champ
$\Curvesstack$ qui paramètre les familles de courbes est algébrique.

\medskip\noindent
Dans la section \ref{section-moduli-complexes}, nous étudions les modules de
complexes sur un morphisme propre et obtenons un champ algébrique
$\Complexesstack_{X/B}$. L'idée de l'énoncé et celle de la démonstration sont
empruntées à \cite{lieblich-complexes}.

\medskip\noindent
Que ne trouve-t-on pas dans ce chapitre ? Il n'y est presque pas question des
propriétés que possèdent les espaces et champs de modules ainsi obtenus (au-delà
de leur algébricité) ; à ce sujet, nous renvoyons à Champs de modules,
section \ref{moduli-section-introduction}. Dans la plupart des résultats
examinés, on peut généraliser les constructions en considérant un morphisme
$\mathcal{X} \to \mathcal{B}$ de champs algébriques au lieu d'un morphisme
$X \to B$ d'espaces algébriques. Nous en traiterons (insérer ici une référence
future). Dans le cas des espaces de Hilbert, il existe une notion plus générale
de « champs de Hilbert », dont nous traiterons dans un chapitre séparé ; voir
(insérer ici une référence future).




\section{Conventions}
\label{section-conventions}

\noindent
Nous avons placé à dessein ce chapitre, ainsi que les chapitres « Exemples de
champs », « Faisceaux sur les champs algébriques », « Critères de
représentabilité » et « Axiomes d'Artin », avant le développement général de
la théorie des champs algébriques. La raison en est qu'à partir du chapitre
suivant (voir Propriétés des champs, section
\ref{stacks-properties-section-conventions}), nous ne distinguerons plus un
schéma du champ algébrique auquel il donne naissance. Notre langage deviendra
donc plus souple et plus facile à lire pour un être humain, mais aussi moins
précis. Ces premiers chapitres, y compris le chapitre initial « Champs
algébriques », posent les fondements qui nous permettront plus tard de négliger
certaines distinctions très techniques entre diverses façons d'envisager les
champs algébriques. Cependant, notamment dans les chapitres « Axiomes d'Artin »
et « Critères de représentabilité », nous devons préciser exactement les objets
avec lesquels nous travaillons, puisque nous cherchons à démontrer que
certaines constructions produisent des champs algébriques ou des espaces
algébriques.

\medskip\noindent
Malheureusement, il en résulte que certaines notations, conventions et certains
termes sont malcommodes et peuvent paraître contraires aux usages au lecteur
plus expérimenté. Nous espérons qu'il nous le pardonnera !

\medskip\noindent
Nous supposerons constamment que tous les schémas sont contenus dans un grand
site fppf $\Sch_{fppf}$. En outre, tout anneau $A$ considéré possède la
propriété que $\Spec(A)$ est (isomorphe à) un objet de ce grand site.

\medskip\noindent
Soit $S$ un schéma et soit $X$ un espace algébrique sur $S$. Dans ce chapitre
et dans le suivant, nous écrirons $X \times_S X$ pour le produit de $X$ par
lui-même (dans la catégorie des espaces algébriques sur $S$), plutôt que
$X \times X$.














\section{Le foncteur Hom}
\label{section-hom}

\noindent
Dans cette section, nous étudions le foncteur des homomorphismes défini
ci-dessous.

\begin{situation}
\label{situation-hom}
Soit $S$ un schéma. Soit $f : X \to B$ un morphisme d'espaces algébriques sur
$S$. Soient $\mathcal{F}$ et $\mathcal{G}$ des $\mathcal{O}_X$-modules
quasi-cohérents. Pour tout schéma $T$ sur $B$, nous noterons $\mathcal{F}_T$
et $\mathcal{G}_T$ les changements de base de $\mathcal{F}$ et
$\mathcal{G}$ à $T$, autrement dit leurs images inverses par le morphisme de
projection $X_T = X \times_B T \to X$. Nous considérons le foncteur
\begin{equation}
\label{equation-hom}
\mathit{Hom}(\mathcal{F}, \mathcal{G}) :
(\Sch/B)^{opp}
\longrightarrow
\textit{Sets},\quad
T
\longrightarrow
\Hom_{\mathcal{O}_{X_T}}(\mathcal{F}_T, \mathcal{G}_T)
\end{equation}
\end{situation}

\noindent
Dans la situation \ref{situation-hom}, nous regarderons parfois
$\mathit{Hom}(\mathcal{F}, \mathcal{G})$ comme un foncteur
$(\Sch/S)^{opp} \to \textit{Sets}$ muni d'un morphisme
$\mathit{Hom}(\mathcal{F}, \mathcal{G}) \to B$. En effet, si $T$ est un
schéma sur $S$, un élément de
$\mathit{Hom}(\mathcal{F}, \mathcal{G})(T)$ consiste en un couple $(h, u)$,
où $h$ est un morphisme $h : T \to B$ et
$u : \mathcal{F}_T \to \mathcal{G}_T$ est un homomorphisme de
$\mathcal{O}_{X_T}$-modules, avec $X_T = T \times_{h, B} X$ et
$\mathcal{F}_T$, $\mathcal{G}_T$ les images inverses sur $X_T$. En
particulier, lorsque nous disons que
$\mathit{Hom}(\mathcal{F}, \mathcal{G})$ est un espace algébrique, nous
entendons que le foncteur correspondant
$(\Sch/S)^{opp} \to \textit{Sets}$ est un espace algébrique.

\begin{lemma}
\label{lemma-hom-sheaf}
Dans la situation \ref{situation-hom}, le foncteur
$\mathit{Hom}(\mathcal{F}, \mathcal{G})$
satisfait à la propriété de faisceau pour la topologie fpqc.
\end{lemma}

\begin{proof}
Soit $\{T_i \to T\}_{i \in I}$ un recouvrement fpqc de schémas sur $B$.
Posons $X_i = X_{T_i} = X \times_S T_i$ et $\mathcal{F}_i = u_{T_i}$,
ainsi que $\mathcal{G}_i = \mathcal{G}_{T_i}$. Remarquons que
$\{X_i \to X_T\}_{i \in I}$ est un recouvrement fpqc de $X_T$ ; voir
Topologies sur les espaces, lemme
\ref{spaces-topologies-lemma-fpqc}. Ainsi, une famille d'homomorphismes
$u_i : \mathcal{F}_i \to \mathcal{G}_i$ tels que $u_i$ et $u_j$
induisent le même homomorphisme sur $X_{T_i \times_T T_j}$ provient par
descente d'un unique homomorphisme
$u : \mathcal{F}_T \to \mathcal{G}_T$
(Descente sur les espaces, proposition
\ref{spaces-descent-proposition-fpqc-descent-quasi-coherent}).
\end{proof}

\begin{lemma}
\label{lemma-extend-hom-to-spaces}
Dans la situation \ref{situation-hom}. Soit $T$ un espace algébrique sur $S$.
Nous avons
$$
\Mor_{\Sh((\Sch/S)_{fppf})}(T, \mathit{Hom}(\mathcal{F}, \mathcal{G})) =
\{(h, u) \mid h : T \to B, u : \mathcal{F}_T \to \mathcal{G}_T\}
$$
où $\mathcal{F}_T, \mathcal{G}_T$ désignent les images inverses de
$\mathcal{F}$ et $\mathcal{G}$ sur l'espace algébrique
$X \times_{B, h} T$.
\end{lemma}

\begin{proof}
Choisissons un schéma $U$ et un morphisme étale surjectif $p : U \to T$.
Soit $R = U \times_T U$, de projections $t, s : R \to U$.

\medskip\noindent
Soit $v : T \to \mathit{Hom}(\mathcal{F}, \mathcal{G})$ une transformation
naturelle. Alors $v(p)$ correspond à un couple $(h_U, u_U)$ sur $U$. Comme
$v$ est une transformation de foncteurs, les images inverses de
$(h_U, u_U)$ par $s$ et $t$ coïncident. Puisque $T = U/R$ (Espaces, lemme
\ref{spaces-lemma-space-presentation}), nous obtenons un morphisme
$h : T \to B$ tel que $h_U = h \circ p$. Alors $\mathcal{F}_U$ est l'image
inverse de $\mathcal{F}_T$ sur $X_U$, et il en est de même pour
$\mathcal{G}_U$. Par conséquent, $u_U$ descend en un homomorphisme de
$\mathcal{O}_{X_T}$-modules $u : \mathcal{F}_T \to \mathcal{G}_T$, d'après
Descente sur les espaces, proposition
\ref{spaces-descent-proposition-fpqc-descent-quasi-coherent}.

\medskip\noindent
Réciproquement, soit $(h, u)$ un couple sur $T$. Nous obtenons alors une
transformation naturelle $v : T \to \mathit{Hom}(\mathcal{F}, \mathcal{G})$
en envoyant un morphisme $a : T' \to T$, où $T'$ est un schéma, sur
$(h \circ a, a^*u)$. Nous omettons de vérifier que cette construction et celle
du paragraphe précédent sont inverses l'une de l'autre.
\end{proof}

\begin{remark}
\label{remark-hom-base-change}
Dans la situation \ref{situation-hom}, soit $B' \to B$ un morphisme d'espaces
algébriques sur $S$. Posons $X' = X \times_B B'$ et désignons par
$\mathcal{F}'$, $\mathcal{G}'$ les images inverses de $\mathcal{F}$,
$\mathcal{G}$ sur $X'$. Nous obtenons alors un foncteur
$\mathit{Hom}(\mathcal{F}', \mathcal{G}') : (\Sch/B')^{opp} \to \textit{Sets}$
associé au changement de base $f' : X' \to B'$. Pour un schéma $T$ sur $B'$,
il est clair que
$$
\mathit{Hom}(\mathcal{F}', \mathcal{G}')(T) =
\mathit{Hom}(\mathcal{F}, \mathcal{G})(T)
$$
où, dans le membre de droite, on regarde $T$ comme un schéma sur $B$ au moyen
du composé $T \to B' \to B$. Cette remarque élémentaire sera parfois utile
pour changer l'espace algébrique de base.
\end{remark}

\begin{lemma}
\label{lemma-hom-sheaf-in-X}
Dans la situation \ref{situation-hom}, soit $\{X_i \to X\}_{i \in I}$ un
recouvrement fppf et, pour chaque $i, j \in I$, soit
$\{X_{ijk} \to X_i \times_X X_j\}$ un recouvrement fppf. Désignons par
$\mathcal{F}_i$, resp.\ $\mathcal{F}_{ijk}$, l'image inverse de
$\mathcal{F}$ sur $X_i$, resp.\ sur $X_{ijk}$. Définissons de même
$\mathcal{G}_i$ et $\mathcal{G}_{ijk}$. Pour tout schéma $T$ sur $B$, le
diagramme
$$
\xymatrix{
\mathit{Hom}(\mathcal{F}, \mathcal{G})(T) \ar[r] &
\prod\nolimits_i
\mathit{Hom}(\mathcal{F}_i, \mathcal{G}_i)(T)
\ar@<1ex>[r]^-{\text{pr}_0^*} \ar@<-1ex>[r]_-{\text{pr}_1^*}
&
\prod\nolimits_{i, j, k}
\mathit{Hom}(\mathcal{F}_{ijk}, \mathcal{G}_{ijk})(T)
}
$$
présente la première flèche comme l'égalisateur des deux autres.
\end{lemma}

\begin{proof}
Soit $u_i : \mathcal{F}_{i, T} \to \mathcal{G}_{i, T}$ un élément de
l'égalisateur de $\text{pr}_0^*$ et $\text{pr}_1^*$. Comme le changement de
base d'un recouvrement fppf est un recouvrement fppf (Topologies sur les
espaces, lemme \ref{spaces-topologies-lemma-fppf}), nous voyons que
$\{X_{i, T} \to X_T\}_{i \in I}$ et
$\{X_{ijk, T} \to X_{i, T} \times_{X_T} X_{j, T}\}$ sont des recouvrements
fppf. En appliquant Descente sur les espaces, proposition
\ref{spaces-descent-proposition-fpqc-descent-quasi-coherent}, nous concluons
d'abord que $u_i$ et $u_j$ induisent le même morphisme sur
$X_{i, T} \times_{X_T} X_{j, T}$ ; une seconde application montre alors qu'il
existe un unique morphisme $u : \mathcal{F}_T \to \mathcal{G}_T$ dont la
restriction est $u_i$ pour tout $i$. Ceci achève la démonstration.
\end{proof}

\begin{lemma}
\label{lemma-hom-limits}
Dans la situation \ref{situation-hom}. Si $\mathcal{F}$ est de présentation
finie et si $f$ est quasi-compact et quasi-séparé, alors
$\mathit{Hom}(\mathcal{F}, \mathcal{G})$ commute aux limites.
\end{lemma}

\begin{proof}
Soit $T = \lim_{i \in I} T_i$ une limite filtrante de $B$-schémas affines.
Nous devons montrer que
$$
\mathit{Hom}(\mathcal{F}, \mathcal{G})(T) =
\colim \mathit{Hom}(\mathcal{F}, \mathcal{G})(T_i)
$$
Choisissons $0 \in I$. Nous pouvons remplacer $B$ par $T_0$, $X$ par
$X_{T_0}$, $\mathcal{F}$ par $\mathcal{F}_{T_0}$, $\mathcal{G}$ par
$\mathcal{G}_{T_0}$ et $I$ par $\{i \in I \mid i \geq 0\}$. Voir la
remarque \ref{remark-hom-base-change}. Nous pouvons donc supposer que
$B = \Spec(R)$ est affine.

\medskip\noindent
Lorsque $B$ est affine, $X$ est quasi-compact et quasi-séparé. Choisissons un
morphisme étale surjectif $U \to X$, où $U$ est un schéma affine (Propriétés
des espaces, lemme
\ref{spaces-properties-lemma-quasi-compact-affine-cover}). Comme $X$ est
quasi-séparé, le schéma $U \times_X U$ est quasi-compact, et nous pouvons
choisir un morphisme étale surjectif $V \to U \times_X U$, où $V$ est un
schéma affine. En appliquant le lemme \ref{lemma-hom-sheaf-in-X}, nous voyons
que $\mathit{Hom}(\mathcal{F}, \mathcal{G})$ est l'égalisateur de deux
applications entre
$$
\mathit{Hom}(\mathcal{F}|_U, \mathcal{G}|_U)
\quad\text{et}\quad
\mathit{Hom}(\mathcal{F}|_V, \mathcal{G}|_V)
$$
Nous sommes ainsi ramenés au cas où $X$ est affine.

\medskip\noindent
Dans le cas affine, l'énoncé du lemme se ramène au problème suivant : étant
donnés un homomorphisme d'anneaux $R \to A$, deux $A$-modules $M$, $N$ et un
système inductif de $R$-algèbres $C = \colim C_i$. Quand l'application
$$
\colim \Hom_{A \otimes_R C_i}(M \otimes_R C_i, N \otimes_R C_i)
\longrightarrow
\Hom_{A \otimes_R C}(M \otimes_R C, N \otimes_R C)
$$
est-elle bijective ? D'après Algèbre, lemme
\ref{algebra-lemma-module-map-property-in-colimit}, elle l'est si
$M \otimes_R C$ est de présentation finie sur $A \otimes_R C$, c'est-à-dire
si $M$ est de présentation finie sur $A$.
\end{proof}

\begin{lemma}
\label{lemma-hom-closed}
Soit $S$ un schéma. Soit $B$ un espace algébrique sur $S$. Soit
$i : X' \to X$ une immersion fermée d'espaces algébriques sur $B$. Soit
$\mathcal{F}$ un $\mathcal{O}_X$-module quasi-cohérent et soit
$\mathcal{G}'$ un $\mathcal{O}_{X'}$-module quasi-cohérent. Alors
$$
\mathit{Hom}(\mathcal{F}, i_*\mathcal{G}') =
\mathit{Hom}(i^*\mathcal{F}, \mathcal{G}')
$$
comme foncteurs sur $(\Sch/B)$.
\end{lemma}

\begin{proof}
Soit $g : T \to B$ un morphisme, où $T$ est un schéma. Désignons par
$i_T : X'_T \to X_T$ le changement de base de $i$. Désignons par
$h : X_T \to X$ et $h' : X'_T \to X'$ les projections. Observons que
$(h')^*i^*\mathcal{F} = i_T^*h^*\mathcal{F}$. Comme une immersion fermée est
affine (Morphismes d'espaces, lemme
\ref{spaces-morphisms-lemma-closed-immersion-affine}), nous avons
$h^*i_*\mathcal{G} = i_{T, *}(h')^*\mathcal{G}$ d'après Cohomologie des
espaces, lemme \ref{spaces-cohomology-lemma-affine-base-change}. Ainsi,
\begin{align*}
\mathit{Hom}(\mathcal{F}, i_*\mathcal{G}')(T)
& =
\Hom_{\mathcal{O}_{X_T}}(h^*\mathcal{F}, h^*i_*\mathcal{G}') \\
& =
\Hom_{\mathcal{O}_{X_T}}(h^*\mathcal{F}, i_{T, *}(h')^*\mathcal{G}) \\
& =
\Hom_{\mathcal{O}_{X'_T}}(i_T^*h^*\mathcal{F}, (h')^*\mathcal{G}) \\
& =
\Hom_{\mathcal{O}_{X'_T}}((h')^*i^*\mathcal{F}, (h')^*\mathcal{G}) \\
& =
\mathit{Hom}(i^*\mathcal{F}, \mathcal{G}')(T),
\end{align*}
comme voulu. L'égalité du milieu résulte de l'adjonction des foncteurs
$i_{T, *}$ et $i_T^*$.
\end{proof}

\begin{lemma}
\label{lemma-cohomology-perfect-complex}
Soit $S$ un schéma. Soit $B$ un espace algébrique sur $S$. Soit $K$ un objet
pseudo-cohérent de $D(\mathcal{O}_B)$.
\begin{enumerate}
\item Si, pour tout $g : T \to B$ dans $(\Sch/B)$, le faisceau de cohomologie
$H^{-1}(Lg^*K)$ est nul, alors le foncteur
$$
(\Sch/B)^{opp} \longrightarrow \textit{Sets},\quad
(g : T \to B) \longmapsto H^0(T, H^0(Lg^*K))
$$
est un espace algébrique affine et de présentation finie sur $B$.
\item Si, pour tout $g : T \to B$ dans $(\Sch/B)$, les faisceaux de
cohomologie $H^i(Lg^*K)$ sont nuls pour $i < 0$, alors $K$ est parfait,
$K$ est localement d'amplitude de Tor contenue dans $[0, b]$, et le foncteur
$$
(\Sch/B)^{opp} \longrightarrow \textit{Sets},\quad
(g : T \to B) \longmapsto H^0(T, Lg^*K)
$$
est un espace algébrique affine et de présentation finie sur $B$.
\end{enumerate}
\end{lemma}

\begin{proof}
Sous les hypothèses de (2), nous avons
$H^0(T, Lg^*K) = H^0(T, H^0(Lg^*K))$. Montrons que la règle
$T \mapsto H^0(T, H^0(Lg^*K))$ satisfait à la propriété de faisceau pour la
topologie fppf. Supposons pour cela que nous disposions d'un recouvrement fppf
$\{h_i : T_i \to T\}$ d'un schéma $g : T \to B$ sur $B$. Posons
$g_i = g \circ h_i$. Remarquons que, puisque $h_i$ est plat, nous avons
$Lh_i^* = h_i^*$ et $h_i^*$ commute à la prise de la cohomologie. Par suite,
$$
H^0(T_i, H^0(Lg_i^*K)) =
H^0(T_i, H^0(h_i^*Lg^*K)) =
H^0(T, h_i^*H^0(Lg^*K))
$$
Il en est de même pour l'image inverse sur $T_i \times_T T_j$. Comme
$Lg^*K$ est un complexe pseudo-cohérent sur $T$ (Cohomologie sur les sites,
lemme \ref{sites-cohomology-lemma-pseudo-coherent-pullback}), le faisceau de
cohomologie $\mathcal{F} = H^0(Lg^*K)$ est quasi-cohérent (Catégories dérivées
des espaces, lemme \ref{spaces-perfect-lemma-pseudo-coherent}). Ainsi, d'après
Descente sur les espaces, proposition
\ref{spaces-descent-proposition-fpqc-descent-quasi-coherent}, nous avons
$$
H^0(T, \mathcal{F}) = \Ker(
\prod H^0(T_i, h_i^*\mathcal{F}) \to
\prod H^0(T_i \times_T T_j, (T_i \times_T T_j \to T)^*\mathcal{F}))
$$
Nous voyons ainsi que les règles de (1) et (2) satisfont à la propriété de
faisceau pour les recouvrements fppf. Nous pouvons donc appliquer Bootstrap,
lemme \ref{bootstrap-lemma-locally-algebraic-space-finite-type}, pour voir
qu'il suffit de démontrer la représentabilité localement pour la topologie
étale sur $B$. De plus, on peut vérifier localement pour la topologie étale
sur $B$ si le résultat final est affine et de présentation finie ; voir
Morphismes d'espaces, lemmes \ref{spaces-morphisms-lemma-affine-local} et
\ref{spaces-morphisms-lemma-finite-presentation-local}. Nous pouvons donc
supposer que $B$ est un schéma affine.

\medskip\noindent
Supposons que $B = \Spec(A)$ soit un schéma affine. D'après les résultats de
Catégories dérivées des espaces, lemmes
\ref{spaces-perfect-lemma-pseudo-coherent},
\ref{spaces-perfect-lemma-derived-quasi-coherent-small-etale-site} et
\ref{spaces-perfect-lemma-descend-pseudo-coherent}, nous pouvons, dans le reste
de la démonstration, considérer $K$ comme un objet parfait de la catégorie
dérivée des complexes de modules sur $B$ pour la topologie de Zariski.
D'après Catégories dérivées des schémas, lemmes
\ref{perfect-lemma-pseudo-coherent},
\ref{perfect-lemma-affine-compare-bounded} et
\ref{perfect-lemma-pseudo-coherent-affine}, nous pouvons trouver un complexe
pseudo-cohérent $M^\bullet$ de $A$-modules tel que $K$ soit l'objet
correspondant de $D(\mathcal{O}_B)$. Notre hypothèse sur les images inverses
implique que $M^\bullet \otimes^\mathbf{L}_A \kappa(\mathfrak p)$ a son
$H^{-1}$ nul pour tout idéal premier $\mathfrak p \subset A$. D'après Plus
d'algèbre, lemme \ref{more-algebra-lemma-cut-complex-in-two}, nous pouvons
écrire
$$
M^\bullet =
\tau_{\geq 0}M^\bullet \oplus \tau_{\leq - 1}M^\bullet
$$
où $\tau_{\geq 0}M^\bullet$ est parfait d'amplitude de Tor contenue dans
$[0, b]$ pour un certain $b \geq 0$ (nous avons aussi utilisé ici Plus
d'algèbre, lemmes \ref{more-algebra-lemma-glue-perfect} et
\ref{more-algebra-lemma-glue-tor-amplitude}). Remarquons que, dans le cas (2),
nous voyons aussi que $\tau_{\leq - 1}M^\bullet = 0$ dans $D(A)$, de sorte que
$M^\bullet$ et $K$ sont parfaits d'amplitude de Tor contenue dans $[0, b]$.
Pour tout $B$-schéma $g : T \to B$, nous avons
$$
H^0(T, H^0(Lg^*K)) = H^0(T, H^0(Lg^*\tau_{\geq 0}K))
$$
(par le dual de Catégories dérivées, lemme
\ref{derived-lemma-negative-vanishing}) ; nous pouvons donc remplacer $K$ par
$\tau_{\geq 0}K$ et, corrélativement, $M^\bullet$ par
$\tau_{\geq 0}M^\bullet$. Autrement dit, nous pouvons supposer que
$M^\bullet$ est d'amplitude de Tor contenue dans $[0, b]$.

\medskip\noindent
Supposons $M^\bullet$ d'amplitude de Tor contenue dans $[0, b]$. Nous pouvons
supposer que $M^\bullet$ est un complexe borné supérieurement de $A$-modules
libres de type fini (d'après notre définition des complexes pseudo-cohérents ;
voir Plus d'algèbre, définition
\ref{more-algebra-definition-pseudo-coherent}, et la discussion qui suit la
définition). D'après Plus d'algèbre, lemme
\ref{more-algebra-lemma-last-one-flat}, nous voyons que
$M = \Coker(M^{- 1} \to M^0)$ est plat. D'après Algèbre, lemme
\ref{algebra-lemma-finite-projective}, nous voyons que $M$ est localement
libre de type fini. Par conséquent, $M^\bullet$ est quasi-isomorphe à
$$
M \to M^1 \to M^2 \to \ldots \to M^d \to 0 \ldots
$$
Remarquons qu'il s'agit d'un complexe K-plat (Cohomologie, lemme
\ref{cohomology-lemma-bounded-flat-K-flat}) ; l'image inverse dérivée de $K$
par un morphisme $T \to B$ est donc calculée par le complexe
$$
g^*\widetilde{M} \to g^*\widetilde{M^1} \to \ldots
$$
Il suffit ainsi de montrer que le foncteur
$$
(g : T \to B) \longmapsto
\Ker(
\Gamma(T,g^*\widetilde{M})
\to
\Gamma(T, g^*(\widetilde{M^1})
)
$$
est représentable par un schéma affine de présentation finie sur $B$.

\medskip\noindent
Nous pouvons encore remplacer $B$ par les membres d'un recouvrement ouvert
affine pour démontrer ce dernier énoncé. Nous pouvons donc supposer que $M$
est libre de type fini (rappelons que $M^1$ est libre de type fini dès le
départ). Écrivons $M = A^{\oplus n}$ et $M^1 = A^{\oplus m}$. Supposons que
l'application $M \to M^1$ soit donnée par la matrice $m \times n$
$(a_{ij})$ à coefficients dans $A$. Alors
$\widetilde{M} = \mathcal{O}_B^{\oplus n}$ et
$\widetilde{M^1} = \mathcal{O}_B^{\oplus m}$. Le foncteur ci-dessus est donc
égal au foncteur
$$
(g : T \to B) \longmapsto
\{(f_1, \ldots, f_n) \in \Gamma(T, \mathcal{O}_T) \mid
\sum g^\sharp(a_{ij})f_i = 0,\ j = 1, \ldots, m\}
$$
Il est clairement représentable par le schéma affine
$$
\Spec\left(A[x_1, \ldots, x_n]/(\sum a_{ij}x_i; j = 1, \ldots, m)\right)
$$
ce qui démontre le lemme.
\end{proof}

\noindent
Le foncteur $\mathit{Hom}(\mathcal{F}, \mathcal{G})$ est représentable dans
plusieurs situations. Tous nos résultats reposeront sur le cas fondamental
suivant. La démonstration ci-dessous de ce lemme est, en un certain sens, la
généralisation naturelle de celle de \cite[III, Cor 7.7.8]{EGA}.

\begin{lemma}
\label{lemma-noetherian-hom}
Dans la situation \ref{situation-hom}, supposons que
\begin{enumerate}
\item $B$ soit un espace algébrique noethérien,
\item $f$ soit localement de type fini et quasi-séparé,
\item $\mathcal{F}$ soit un $\mathcal{O}_X$-module de type fini, et
\item $\mathcal{G}$ soit un $\mathcal{O}_X$-module de type fini, plat sur
$B$, à support propre sur $B$.
\end{enumerate}
Alors le foncteur $\mathit{Hom}(\mathcal{F}, \mathcal{G})$ est un espace
algébrique affine et de présentation finie sur $B$.
\end{lemma}

\begin{proof}
Nous pouvons remplacer $X$ par un voisinage ouvert quasi-compact du support de
$\mathcal{G}$ ; nous pouvons donc supposer $X$ noethérien. Dans ce cas, $X$ et
$f$ sont quasi-compacts et quasi-séparés. Choisissons une approximation
$P \to \mathcal{F}$ par un complexe parfait $P$ du triplet
$(X, \mathcal{F}, -1)$ ; voir Catégories dérivées des espaces, définition
\ref{spaces-perfect-definition-approximation-holds} et théorème
\ref{spaces-perfect-theorem-approximation}). Alors l'application induite
$$
\Hom_{\mathcal{O}_X}(\mathcal{F}, \mathcal{G})
\longrightarrow
\Hom_{D(\mathcal{O}_X)}(P, \mathcal{G})
$$
est un isomorphisme, parce que $P \to \mathcal{F}$ induit un isomorphisme
$H^0(P) \to \mathcal{F}$ et que $H^i(P) = 0$ pour $i > 0$. De plus, pour tout
morphisme $g : T \to B$, désignons par $h : X_T = T \times_B X \to X$ la
projection et posons $P_T = Lh^*P$. Il est alors encore vrai que
$$
\Hom_{\mathcal{O}_{X_T}}(\mathcal{F}_T, \mathcal{G}_T)
\longrightarrow
\Hom_{D(\mathcal{O}_{X_T})}(P_T, \mathcal{G}_T)
$$
est un isomorphisme, puisque
$P_T = Lh^*P \to Lh^*\mathcal{F} \to \mathcal{F}_T$ induit un isomorphisme
$H^0(P_T) \to \mathcal{F}_T$ (car $h^*$ est exact à droite et
$H^i(P) = 0$ pour $i > 0$). Il suffit donc de démontrer le résultat pour le
foncteur
$$
T \longmapsto \Hom_{D(\mathcal{O}_{X_T})}(P_T, \mathcal{G}_T).
$$
D'après la suite spectrale de Leray (voir Cohomologie sur les sites, remarque
\ref{sites-cohomology-remark-before-Leray}), nous avons
$$
\Hom_{D(\mathcal{O}_{X_T})}(P_T, \mathcal{G}_T) =
H^0(X_T, R\SheafHom(P_T, \mathcal{G}_T)) =
H^0(T, Rf_{T, *}R\SheafHom(P_T, \mathcal{G}_T))
$$
où $f_T : X_T \to T$ est le changement de base de $f$. D'après Catégories
dérivées des espaces, lemme
\ref{spaces-perfect-lemma-base-change-RHom}, nous avons
$$
Rf_{T, *}R\SheafHom(P_T, \mathcal{G}_T) = Lg^*Rf_*R\SheafHom(P, \mathcal{G}).
$$
D'après Catégories dérivées des espaces, lemme
\ref{spaces-perfect-lemma-ext-perfect}, l'objet
$K = Rf_*R\SheafHom(P, \mathcal{G})$ de $D(\mathcal{O}_B)$ est parfait. Nous
pouvons donc appliquer le lemme \ref{lemma-cohomology-perfect-complex}, pourvu
que nous démontrions que le faisceau de cohomologie $H^i(Lg^*K)$ est $0$ pour
tout $i < 0$ et tout $g : T \to B$ comme ci-dessus. Ceci résulte clairement de
la dernière formule affichée, puisque les faisceaux de cohomologie de
$Rf_{T, *}R\SheafHom(P_T, \mathcal{G}_T)$ sont nuls en degrés négatifs : en
effet, les faisceaux de cohomologie de
$R\SheafHom(P_T, \mathcal{G}_T)$ sont nuls en degrés négatifs, car $P_T$ est
parfait et ses faisceaux de cohomologie sont nuls en degrés positifs.
\end{proof}

\noindent
Voici une conséquence facile du lemme \ref{lemma-noetherian-hom}.

\begin{proposition}
\label{proposition-hom}
Dans la situation \ref{situation-hom}, supposons que
\begin{enumerate}
\item $f$ soit de présentation finie, et
\item $\mathcal{G}$ soit un $\mathcal{O}_X$-module de présentation finie,
plat sur $B$, à support propre sur $B$.
\end{enumerate}
Alors le foncteur $\mathit{Hom}(\mathcal{F}, \mathcal{G})$ est un espace
algébrique affine sur $B$. Si $\mathcal{F}$ est de présentation finie, alors
$\mathit{Hom}(\mathcal{F}, \mathcal{G})$ est de présentation finie sur $B$.
\end{proposition}

\begin{proof}
D'après le lemme \ref{lemma-hom-sheaf}, le foncteur
$\mathit{Hom}(\mathcal{F}, \mathcal{G})$ satisfait à la propriété de faisceau
pour les recouvrements fppf. Cela signifie que nous pouvons\footnote{Nous
omettons de vérifier la condition ensembliste (3) du lemme cité.} appliquer
Bootstrap, lemme \ref{bootstrap-lemma-locally-algebraic-space}, pour vérifier
la représentabilité localement pour la topologie étale sur $B$. De plus, nous
pouvons vérifier localement pour la topologie étale sur $B$ si le résultat
final est affine ou de présentation finie ; voir Morphismes d'espaces, lemmes
\ref{spaces-morphisms-lemma-affine-local} et
\ref{spaces-morphisms-lemma-finite-presentation-local}. Nous pouvons donc
supposer que $B$ est un schéma affine.

\medskip\noindent
Supposons que $B$ soit un schéma affine. Comme $f$ est de présentation finie,
il en résulte que $X$ est quasi-compact et quasi-séparé. Nous pouvons donc
écrire $\mathcal{F} = \colim \mathcal{F}_i$ comme une colimite filtrante de
$\mathcal{O}_X$-modules de présentation finie (Limites d'espaces, lemme
\ref{spaces-limits-lemma-colimit-finitely-presented}). Il est clair que
$$
\mathit{Hom}(\mathcal{F}, \mathcal{G}) =
\lim \mathit{Hom}(\mathcal{F}_i, \mathcal{G})
$$
Si nous pouvons donc montrer que chaque
$\mathit{Hom}(\mathcal{F}_i, \mathcal{G})$ est représentable par un schéma
affine, nous verrons qu'il en va de même de
$\mathit{Hom}(\mathcal{F}, \mathcal{G})$. Utiliser les résultats de Limites,
section \ref{limits-section-limits}, et de Limites d'espaces, section
\ref{spaces-limits-section-limits}. Nous pouvons donc supposer que
$\mathcal{F}$ est de présentation finie.

\medskip\noindent
Écrivons $B = \Spec(R)$. Écrivons $R = \colim R_i$, chaque $R_i$ étant une
$\mathbf{Z}$-algèbre de type fini. Posons $B_i = \Spec(R_i)$. D'après les
résultats de Limites d'espaces, lemmes
\ref{spaces-limits-lemma-descend-finite-presentation} et
\ref{spaces-limits-lemma-descend-modules-finite-presentation}, nous pouvons
trouver un $i$, un morphisme d'espaces algébriques $X_i \to B_i$, et des
$\mathcal{O}_{X_i}$-modules de présentation finie $\mathcal{F}_i$ et
$\mathcal{G}_i$, tels que le changement de base de
$(X_i, \mathcal{F}_i, \mathcal{G}_i)$ à $B$ redonne
$(X, \mathcal{F}, \mathcal{G})$. D'après Limites d'espaces, lemme
\ref{spaces-limits-lemma-descend-flat}, nous pouvons, quitte à augmenter $i$,
supposer que $\mathcal{G}_i$ est plat sur $B_i$. D'après Limites d'espaces,
lemme \ref{spaces-limits-lemma-eventually-proper-support}, nous pouvons de
même supposer que le support schématique de $\mathcal{G}_i$ est propre sur
$B_i$. Nous pouvons alors appliquer le lemme \ref{lemma-noetherian-hom} pour
voir que $H_i = \mathit{Hom}(\mathcal{F}_i, \mathcal{G}_i)$ est un espace
algébrique affine de présentation finie sur $B_i$. En effectuant l'image
inverse sur $B$ (à l'aide de la remarque \ref{remark-hom-base-change}), nous
voyons que $H_i \times_{B_i} B = \mathit{Hom}(\mathcal{F}, \mathcal{G})$,
et nous avons gagné.
\end{proof}

\noindent
Vérification de cohérence : le faisceau $\mathit{Hom}$ joue le même rôle parmi
les espaces algébriques sur $S$.

\section{Le foncteur Isom}
\label{section-isom}

\noindent
Dans la situation \ref{situation-hom}, nous pouvons considérer le
sous-foncteur
$$
\mathit{Isom}(\mathcal{F}, \mathcal{G}) \subset
\mathit{Hom}(\mathcal{F}, \mathcal{G})
$$
dont la valeur sur un schéma $T$ sur $B$ est l'ensemble des homomorphismes
{\it inversibles} de $\mathcal{O}_{X_T}$-modules
$u : \mathcal{F}_T \to \mathcal{G}_T$.

\medskip\noindent
Nous regardons parfois $\mathit{Isom}(\mathcal{F}, \mathcal{G})$ comme un
foncteur $(\Sch/S)^{opp} \to \textit{Sets}$ muni d'un morphisme
$\mathit{Isom}(\mathcal{F}, \mathcal{G}) \to B$. En effet, si $T$ est un
schéma sur $S$, un élément de
$\mathit{Isom}(\mathcal{F}, \mathcal{G})(T)$ consiste en un couple $(h, u)$,
où $h$ est un morphisme $h : T \to B$ et
$u : \mathcal{F}_T \to \mathcal{G}_T$ est un isomorphisme de
$\mathcal{O}_{X_T}$-modules, avec $X_T = T \times_{h, B} X$ et
$\mathcal{F}_T$, $\mathcal{G}_T$ les images inverses sur $X_T$. En
particulier, lorsque nous disons que
$\mathit{Isom}(\mathcal{F}, \mathcal{G})$ est un espace algébrique, nous
entendons que le foncteur correspondant
$(\Sch/S)^{opp} \to \textit{Sets}$ est un espace algébrique.

\begin{lemma}
\label{lemma-isom-sheaf}
Dans la situation \ref{situation-hom}, le foncteur
$\mathit{Isom}(\mathcal{F}, \mathcal{G})$
satisfait à la propriété de faisceau pour la topologie fpqc.
\end{lemma}

\begin{proof}
Nous avons déjà vu que $\mathit{Hom}(\mathcal{F}, \mathcal{G})$ satisfait à
la propriété de faisceau. Il reste donc à montrer ceci : si
$\{T_i \to T\}_{i \in I}$ est un recouvrement fpqc de schémas sur $B$ et si
$u : \mathcal{F}_T \to \mathcal{G}_T$ est une application
$\mathcal{O}_{X_T}$-linéaire telle que $u_{T_i}$ soit un isomorphisme pour tout
$i$, alors $u$ est un isomorphisme. Puisque
$\{X_i \to X_T\}_{i \in I}$ est un recouvrement fpqc de $X_T$ ; voir
Topologies sur les espaces, lemme \ref{spaces-topologies-lemma-fpqc}, ceci
résulte de Descente sur les espaces, proposition
\ref{spaces-descent-proposition-fpqc-descent-quasi-coherent}.
\end{proof}

\noindent
Vérification de cohérence : le faisceau $\mathit{Isom}$ joue le même rôle parmi
les espaces algébriques sur $S$.

\begin{lemma}
\label{lemma-extend-isom-to-spaces}
Dans la situation \ref{situation-hom}. Soit $T$ un espace algébrique sur $S$.
Nous avons
$$
\Mor_{\Sh((\Sch/S)_{fppf})}(T, \mathit{Isom}(\mathcal{F}, \mathcal{G})) =
\{(h, u) \mid
h : T \to B, u : \mathcal{F}_T \to \mathcal{G}_T\text{ isomorphisme}\}
$$
où $\mathcal{F}_T, \mathcal{G}_T$ désignent les images inverses de
$\mathcal{F}$ et $\mathcal{G}$ sur l'espace algébrique
$X \times_{B, h} T$.
\end{lemma}

\begin{proof}
Observons que les membres de gauche et de droite de l'égalité sont des
sous-ensembles des membres de gauche et de droite de l'égalité du lemme
\ref{lemma-extend-hom-to-spaces}. Nous omettons de vérifier que ces
sous-ensembles se correspondent par l'identification donnée dans la
démonstration de ce lemme.
\end{proof}

\begin{proposition}
\label{proposition-isom}
Dans la situation \ref{situation-hom}, supposons que
\begin{enumerate}
\item $f$ soit de présentation finie, et
\item $\mathcal{F}$ et $\mathcal{G}$ soient des $\mathcal{O}_X$-modules de
présentation finie, plats sur $B$, à support propre sur $B$.
\end{enumerate}
Alors le foncteur $\mathit{Isom}(\mathcal{F}, \mathcal{G})$ est un espace
algébrique affine de présentation finie sur $B$.
\end{proposition}

\begin{proof}
Nous emploierons les abréviations
$H = \mathit{Hom}(\mathcal{F}, \mathcal{G})$,
$I = \mathit{Hom}(\mathcal{F}, \mathcal{F})$,
$H' = \mathit{Hom}(\mathcal{G}, \mathcal{F})$ et
$I' = \mathit{Hom}(\mathcal{G}, \mathcal{G})$.
D'après la proposition \ref{proposition-hom}, les foncteurs
$H$, $I$, $H'$, $I'$ sont des espaces algébriques et les morphismes
$H \to B$, $I \to B$, $H' \to B$ et $I' \to B$ sont affines et de
présentation finie. La composition des applications donne un morphisme
$$
c : H' \times_B H \longrightarrow I \times_B I',\quad
(u', u) \longmapsto (u \circ u', u' \circ u)
$$
d'espaces algébriques sur $B$. Comme $I \times_B I' \to B$ est séparé, la
section $\sigma : B \to I \times_B I'$ correspondant à
$(\text{id}_\mathcal{F}, \text{id}_\mathcal{G})$ est une immersion fermée
(Morphismes d'espaces, lemme
\ref{spaces-morphisms-lemma-section-immersion}). De plus, $\sigma$ est de
présentation finie (Morphismes d'espaces, lemme
\ref{spaces-morphisms-lemma-finite-presentation-permanence}). Par conséquent,
$$
\mathit{Isom}(\mathcal{F}, \mathcal{G}) =
(H' \times_B H) \times_{c, I \times_B I', \sigma} B
$$
est lui aussi un espace algébrique affine de présentation finie sur $B$.
Nous omettons quelques détails.
\end{proof}






\section{Le champ des faisceaux cohérents}
\label{section-stack-coherent-sheaves}

\noindent
Dans cette section, nous démontrons que le champ des faisceaux cohérents sur
$X/B$ est algébrique sous des hypothèses convenables. C'est un cas particulier
de \cite[Theorem 2.1.1]{lieblich_remarks}, qui traite du champ des faisceaux
cohérents sur un champ d'Artin au-dessus d'une base.

\begin{situation}
\label{situation-coherent}
Soit $S$ un schéma. Soit $f : X \to B$ un morphisme d'espaces algébriques sur
$S$. Supposons que $f$ soit de présentation finie. Nous désignons par
$\Cohstack_{X/B}$ la catégorie dont les objets sont les triplets
$(T, g, \mathcal{F})$ où
\begin{enumerate}
\item $T$ est un schéma sur $S$,
\item $g : T \to B$ est un morphisme sur $S$, et, en posant
$X_T = T \times_{g, B} X$,
\item $\mathcal{F}$ est un $\mathcal{O}_{X_T}$-module quasi-cohérent de
présentation finie, plat sur $T$, à support propre sur $T$.
\end{enumerate}
Un morphisme $(T, g, \mathcal{F}) \to (T', g', \mathcal{F}')$ est donné par
un couple $(h, \varphi)$ où
\begin{enumerate}
\item $h : T \to T'$ est un morphisme de schémas sur $B$
(c'est-à-dire que $g' \circ h = g$), et
\item $\varphi : (h')^*\mathcal{F}' \to \mathcal{F}$ est un isomorphisme de
$\mathcal{O}_{X_T}$-modules, où $h' : X_T \to X_{T'}$ est le changement de
base de $h$.
\end{enumerate}
\end{situation}

\noindent
Ainsi, $\Cohstack_{X/B}$ est une catégorie et la règle
$$
p : \Cohstack_{X/B} \longrightarrow (\Sch/S)_{fppf},
\quad
(T, g, \mathcal{F}) \longmapsto T
$$
est un foncteur. Pour un schéma $T$ sur $S$, nous désignons par
$\Cohstack_{X/B, T}$ la catégorie fibre de $p$ en $T$. Ces catégories fibres
sont des groupoïdes.

\begin{lemma}
\label{lemma-coherent-fibred-in-groupoids}
Dans la situation \ref{situation-coherent}, le foncteur
$p : \Cohstack_{X/B} \longrightarrow (\Sch/S)_{fppf}$ est fibré en
groupoïdes.
\end{lemma}

\begin{proof}
Nous montrons que $p$ est fibré en groupoïdes en vérifiant les conditions (1)
et (2) de Catégories, définition
\ref{categories-definition-fibred-groupoids}. Étant donnés un objet
$(T', g', \mathcal{F}')$ de $\Cohstack_{X/B}$ et un morphisme
$h : T \to T'$ de schémas sur $S$, nous pouvons poser $g = h \circ g'$ et
$\mathcal{F} = (h')^*\mathcal{F}'$, où $h' : X_T \to X_{T'}$ est le
changement de base de $h$. Il est alors clair que nous obtenons un morphisme
$(T, g, \mathcal{F}) \to (T', g', \mathcal{F}')$ de $\Cohstack_{X/B}$
au-dessus de $h$. Cela démontre (1). Pour (2), supposons donnés des morphismes
$$
(h_1, \varphi_1) : (T_1, g_1, \mathcal{F}_1) \to (T, g, \mathcal{F})
\quad\text{et}\quad
(h_2, \varphi_2) : (T_2, g_2, \mathcal{F}_2) \to (T, g, \mathcal{F})
$$
de $\Cohstack_{X/B}$ et un morphisme $h : T_1 \to T_2$ tel que
$h_2 \circ h = h_1$. Nous pouvons alors prendre pour $\varphi$ le composé
$$
(h')^*\mathcal{F}_2
\xrightarrow{(h')^*\varphi_2^{-1}}
(h')^*(h_2)^*\mathcal{F} = (h_1)^*\mathcal{F}
\xrightarrow{\varphi_1}
\mathcal{F}_1
$$
et obtenir ainsi le morphisme
$(h, \varphi) : (T_1, g_1, \mathcal{F}_1) \to (T_2, g_2, \mathcal{F}_2)$
qui atteste la condition (2).
\end{proof}

\begin{lemma}
\label{lemma-coherent-diagonal}
Dans la situation \ref{situation-coherent}. Posons
$\mathcal{X} = \Cohstack_{X/B}$. Alors
$\Delta : \mathcal{X} \to \mathcal{X} \times \mathcal{X}$ est représentable
par des espaces algébriques.
\end{lemma}

\begin{proof}
Considérons deux objets $x = (T, g, \mathcal{F})$ et
$y = (T, h, \mathcal{G})$ de $\mathcal{X}$ sur un schéma $T$. Nous devons
montrer que $\mathit{Isom}_\mathcal{X}(x, y)$ est un espace algébrique sur
$T$ ; voir Champs algébriques, lemme
\ref{algebraic-lemma-representable-diagonal}. Si, pour $a : T' \to T$, les
restrictions $x|_{T'}$ et $y|_{T'}$ sont isomorphes dans la catégorie fibre
$\mathcal{X}_{T'}$, alors $g \circ a = h \circ a$. Il existe donc une
transformation de préfaisceaux
$$
\mathit{Isom}_\mathcal{X}(x, y) \longrightarrow \text{Equalizer}(g, h)
$$
Comme la diagonale de $B$ est représentable (par des schémas), cet égalisateur
est un schéma. Nous pouvons donc remplacer $T$ par cet égalisateur et les
faisceaux $\mathcal{F}$ et $\mathcal{G}$ par leurs images inverses. Nous
pouvons ainsi supposer $g = h$. Dans ce cas,
$\mathit{Isom}_\mathcal{X}(x, y) = \mathit{Isom}(\mathcal{F}, \mathcal{G})$,
et le résultat découle de la proposition \ref{proposition-isom}.
\end{proof}

\begin{lemma}
\label{lemma-coherent-stack}
Dans la situation \ref{situation-coherent}, le foncteur
$p : \Cohstack_{X/B} \longrightarrow (\Sch/S)_{fppf}$ est un champ en
groupoïdes.
\end{lemma}

\begin{proof}
Pour démontrer que $\Cohstack_{X/B}$ est un champ en groupoïdes, il faut
montrer que les préfaisceaux $\mathit{Isom}$ sont des faisceaux et que les
données de descente sont effectives. L'assertion relative à $\mathit{Isom}$
résulte du lemme \ref{lemma-coherent-diagonal} ; voir Champs algébriques,
lemme \ref{algebraic-lemma-representable-diagonal}. Démontrons l'assertion
relative aux données de descente. Supposons que
$\{a_i : T_i \to T\}$ soit un recouvrement fppf de schémas sur $S$. Soit
$(\xi_i, \varphi_{ij})$ une donnée de descente pour $\{T_i \to T\}$ à valeurs
dans $\Cohstack_{X/B}$. Pour chaque $i$, écrivons
$\xi_i = (T_i, g_i, \mathcal{F}_i)$. Désignons par
$\text{pr}_0 : T_i \times_T T_j \to T_i$ et
$\text{pr}_1 : T_i \times_T T_j \to T_j$ les projections. La condition
$\xi_i|_{T_i \times_T T_j} = \xi_j|_{T_i \times_T T_j}$ implique notamment
que $g_i \circ \text{pr}_0 = g_j \circ \text{pr}_1$. Il existe donc un unique
morphisme $g : T \to B$ tel que $g_i = g \circ a_i$ ; voir Descente sur les
espaces, lemme
\ref{spaces-descent-lemma-fpqc-universal-effective-epimorphisms}.
Désignons par $X_T = T \times_{g, B} X$. Posons
$X_i = X_{T_i} = T_i \times_{g_i, B} X = T_i \times_{a_i, T} X_T$ et
$$
X_{ij} = X_{T_i} \times_{X_T} X_{T_j} = X_i \times_{X_T} X_j
$$
avec les projections $\text{pr}_i$ et $\text{pr}_j$ sur $X_i$ et $X_j$.
Observons que l'image inverse de $(T_i, g_i, \mathcal{F}_i)$ par
$\text{pr}_0 : T_i \times_T T_j \to T_i$ est donnée par
$(T_i \times_T T_j, g_i \circ \text{pr}_0, \text{pr}_i^*\mathcal{F}_i)$.
Ainsi, une donnée de descente pour $\{T_i \to T\}$ dans
$\Cohstack_{X/B}$ est donnée par les objets
$(T_i, g \circ a_i, \mathcal{F}_i)$ et, pour chaque couple $i, j$, par un
isomorphisme de $\mathcal{O}_{X_{ij}}$-modules
$$
\varphi_{ij} :
\text{pr}_i^*\mathcal{F}_i \longrightarrow \text{pr}_j^*\mathcal{F}_j
$$
satisfaisant à la condition de cocycle sur (l'image inverse de $X$ à)
$T_i \times_T T_j \times_T T_k$. Bien ; il suffit maintenant d'utiliser le
fait que $\{X_i \to X_T\}$ est un recouvrement fppf pour regarder
$(\mathcal{F}_i, \varphi_{ij})$ comme une donnée de descente pour ce
recouvrement. D'après Descente sur les espaces, proposition
\ref{spaces-descent-proposition-fpqc-descent-quasi-coherent}, cette donnée de
descente est effective et nous obtenons un faisceau quasi-cohérent
$\mathcal{F}$ sur $X_T$, dont la restriction à $X_i$ est $\mathcal{F}_i$.
D'après Morphismes d'espaces, lemme
\ref{spaces-morphisms-lemma-flat-permanence}, nous voyons que $\mathcal{F}$ est
plat sur $T$, et Descente sur les espaces, lemme
\ref{spaces-descent-lemma-finite-presentation-descends}, garantit que
$\mathcal{F}$ est de présentation finie comme $\mathcal{O}_{X_T}$-module.
Enfin, d'après Descente sur les espaces, lemme
\ref{spaces-descent-lemma-descending-property-proper}, nous voyons que le
support schématique de $\mathcal{F}$ est propre sur $T$, puisque nous avons
supposé le support schématique de $\mathcal{F}_i$ propre sur $T_i$
(remarquons que la formation du support schématique commute au changement de
base plat d'après Morphismes d'espaces, lemme
\ref{spaces-morphisms-lemma-flat-pullback-support}). Nous obtenons ainsi
l'objet recherché sur $T$.
\end{proof}

\begin{remark}
\label{remark-coherent-base-change}
Dans la situation \ref{situation-coherent}, la règle
$(T, g, \mathcal{F}) \mapsto (T, g)$ définit un $1$-morphisme
$$
\Cohstack_{X/B} \longrightarrow \mathcal{S}_B
$$
de champs en groupoïdes (voir le lemme \ref{lemma-coherent-stack}, Champs
algébriques, section \ref{algebraic-section-split}, et Exemples de champs,
section \ref{examples-stacks-section-stack-associated-to-sheaf}). Soit
$B' \to B$ un morphisme d'espaces algébriques sur $S$. Soit
$\mathcal{S}_{B'} \to \mathcal{S}_B$ le $1$-morphisme associé de champs
fibrés en ensembles. Posons $X' = X \times_B B'$. Nous obtenons un champ en
groupoïdes $\Cohstack_{X'/B'} \to (\Sch/S)_{fppf}$ associé au changement de
base $f' : X' \to B'$. Dans cette situation, le diagramme
$$
\vcenter{
\xymatrix{
\Cohstack_{X'/B'} \ar[r] \ar[d] & \Cohstack_{X/B} \ar[d] \\
\mathcal{S}_{B'} \ar[r] & \mathcal{S}_B
}
}
\quad
\begin{matrix}
\text{ou, dans} \\
\text{une autre} \\
\text{notation}
\end{matrix}
\quad
\vcenter{
\xymatrix{
\Cohstack_{X'/B'} \ar[r] \ar[d] & \Cohstack_{X/B} \ar[d] \\
\Sch/B' \ar[r] & \Sch/B
}
}
$$
est un carré $2$-cartésien. Cette remarque élémentaire sera parfois utile pour
changer l'espace algébrique de base.
\end{remark}

\begin{lemma}
\label{lemma-coherent-limits}
Dans la situation \ref{situation-coherent}, supposons que $B \to S$ soit
localement de présentation finie. Alors
$p : \Cohstack_{X/B} \to (\Sch/S)_{fppf}$ commute aux limites (Axiomes
d'Artin, définition \ref{artin-definition-limit-preserving}).
\end{lemma}

\begin{proof}
Notons $B(T)$ la catégorie discrète dont les objets sont les $S$-morphismes
$T \to B$. Soit $T = \lim T_i$ une limite filtrante de schémas affines sur
$S$. Associer à un objet $(T, h, \mathcal{F})$ de
$\Cohstack_{X/B, T}$ l'objet $h$ de $B(T)$ donne un diagramme commutatif de
catégories fibres
$$
\xymatrix{
\colim \Cohstack_{X/B, T_i} \ar[r] \ar[d] &
\Cohstack_{X/B, T} \ar[d] \\
\colim B(T_i) \ar[r] & B(T)
}
$$
Nous devons montrer que la flèche horizontale supérieure est une équivalence.
Puisque nous avons supposé que $B$ est localement de présentation finie sur
$S$, il résulte de Limites d'espaces, remarque
\ref{spaces-limits-remark-limit-preserving}, que la flèche horizontale
inférieure est une équivalence. Cela signifie que nous pouvons supposer que
$T = \lim T_i$ soit une limite filtrante de schémas affines sur $B$.
Désignons par $g_i : T_i \to B$ et $g : T \to B$ les morphismes
correspondants. Posons $X_i = T_i \times_{g_i, B} X$ et
$X_T = T \times_{g, B} X$. Observons que $X_T = \colim X_i$ et que les
espaces algébriques $X_i$ et $X_T$ sont quasi-séparés et quasi-compacts
(puisqu'ils sont de présentation finie sur les schémas affines $T_i$ et $T$).
D'après Limites d'espaces, lemme
\ref{spaces-limits-lemma-descend-modules-finite-presentation}, nous avons
$$
\colim \textit{FP}(X_i) = \textit{FP}(X_T).
$$
où $\textit{FP}(W)$ abrège la catégorie des $\mathcal{O}_W$-modules de
présentation finie. Les résultats de Limites d'espaces, lemmes
\ref{spaces-limits-lemma-descend-flat} et
\ref{spaces-limits-lemma-eventually-proper-support}, nous disent que la même
assertion demeure vraie si nous remplaçons $\textit{FP}(X_i)$ et
$\textit{FP}(X_T)$ par la sous-catégorie pleine des objets plats sur $T_i$ et
$T$, dont le support schématique est propre sur $T_i$ et $T$. Cela démontre le
lemme.
\end{proof}

\begin{lemma}
\label{lemma-coherent-RS-star}
Dans la situation \ref{situation-coherent}. Soit
$$
\xymatrix{
Z \ar[r] \ar[d] & Z' \ar[d] \\
Y \ar[r] & Y'
}
$$
une somme amalgamée dans la catégorie des schémas sur $S$, où
$Z \to Z'$ est un épaississement et $Z \to Y$ est affine ; voir Plus sur les
morphismes, lemme
\ref{more-morphisms-lemma-pushout-along-thickening}. Alors le foncteur entre
catégories fibres
$$
\Cohstack_{X/B, Y'}
\longrightarrow
\Cohstack_{X/B, Y} \times_{\Cohstack_{X/B, Z}} \Cohstack_{X/B, Z'}
$$
est une équivalence.
\end{lemma}

\begin{proof}
Observons que l'application correspondante
$$
B(Y') \longrightarrow B(Y) \times_{B(Z)} B(Z')
$$
est une bijection ; voir Sommes amalgamées d'espaces, lemme
\ref{spaces-pushouts-lemma-pushout-along-thickening-schemes}. Le diagramme
commutatif
$$
\xymatrix{
\Cohstack_{X/B, Y'} \ar[r] \ar[d] &
\Cohstack_{X/B, Y} \times_{\Cohstack_{X/B, Z}} \Cohstack_{X/B, Z'}
\ar[d] \\
B(Y') \ar[r] & B(Y) \times_{B(Z)} B(Z')
}
$$
montre donc que nous pouvons supposer que $Y'$ est un schéma sur $B'$.
D'après la remarque \ref{remark-coherent-base-change}, nous pouvons remplacer
$B$ par $Y'$ et $X$ par $X \times_B Y'$. Nous pouvons ainsi supposer
$B = Y'$. Dans ce cas, l'assertion résulte de Sommes amalgamées d'espaces,
lemme \ref{spaces-pushouts-lemma-space-over-pushout-flat-modules}.
\end{proof}

\begin{lemma}
\label{lemma-coherent-over-first-order-thickening}
Soit
$$
\xymatrix{
X \ar[d] \ar[r]_i & X' \ar[d] \\
T \ar[r] & T'
}
$$
un carré cartésien d'espaces algébriques, où $T \to T'$ est un
épaississement du premier ordre. Soit $\mathcal{F}'$ un
$\mathcal{O}_{X'}$-module plat sur $T'$. Posons
$\mathcal{F} = i^*\mathcal{F}'$. Les conditions suivantes sont équivalentes :
\begin{enumerate}
\item $\mathcal{F}'$ est un $\mathcal{O}_{X'}$-module quasi-cohérent de
présentation finie,
\item $\mathcal{F}'$ est un $\mathcal{O}_{X'}$-module de présentation finie,
\item $\mathcal{F}$ est un $\mathcal{O}_X$-module quasi-cohérent de
présentation finie,
\item $\mathcal{F}$ est un $\mathcal{O}_X$-module de présentation finie.
\end{enumerate}
\end{lemma}

\begin{proof}
Rappelons qu'un module de présentation finie est quasi-cohérent ; cela donne
les équivalences entre (1) et (2), et entre (3) et (4). L'équivalence de (2)
et (4) est un cas particulier de Théorie des déformations, lemme
\ref{defos-lemma-deform-fp-module-ringed-topoi}.
\end{proof}

\begin{lemma}
\label{lemma-coherent-tangent-space}
Dans la situation \ref{situation-coherent}, supposons que $S$ soit un schéma
localement noethérien et que $B \to S$ soit localement de présentation finie.
Soit $k$ un corps de type fini sur $S$ et soit
$x_0 = (\Spec(k), g_0, \mathcal{G}_0)$ un objet de
$\mathcal{X} = \Cohstack_{X/B}$ sur $k$. Alors les espaces
$T\mathcal{F}_{\mathcal{X}, k, x_0}$ et
$\text{Inf}(\mathcal{F}_{\mathcal{X}, k, x_0})$ (Axiomes d'Artin, section
\ref{artin-section-tangent-spaces}) sont de dimension finie.
\end{lemma}

\begin{proof}
Observons que, d'après le lemme \ref{lemma-coherent-RS-star}, notre champ en
groupoïdes $\mathcal{X}$ satisfait à la propriété (RS*) définie dans Axiomes
d'Artin, section \ref{artin-section-inf}. En particulier, $\mathcal{X}$
satisfait à (RS). Toutes les catégories de prédéformations associées sont donc
des catégories de déformations (Axiomes d'Artin, lemme
\ref{artin-lemma-deformation-category}), et l'énoncé a un sens.

\medskip\noindent
Nous montrons dans ce paragraphe que nous pouvons nous ramener au cas
$B = \Spec(k)$. Posons $X_0 = \Spec(k) \times_{g_0, B} X$ et désignons par
$\mathcal{X}_0 = \Cohstack_{X_0/k}$. Dans la remarque
\ref{remark-coherent-base-change}, nous avons vu que $\mathcal{X}_0$ est le
$2$-produit fibré de $\mathcal{X}$ et $\Spec(k)$ sur $B$, comme catégories
fibrées en groupoïdes sur $(\Sch/S)_{fppf}$. Ainsi, d'après Axiomes d'Artin,
lemme \ref{artin-lemma-fibre-product-tangent-spaces}, nous sommes ramenés à
démontrer que $B$, $\Spec(k)$ et $\mathcal{X}_0$ possèdent des espaces
tangents et des espaces d'automorphismes infinitésimaux de dimension finie.
Les espaces tangents de $B$ et de $\Spec(k)$ sont de dimension finie d'après
Axiomes d'Artin, lemme \ref{artin-lemma-finite-dimension}, et leur
$\text{Inf}$ est évidemment nul. Il suffit donc de traiter $\mathcal{X}_0$.

\medskip\noindent
Soit $k[\epsilon]$ l'algèbre des nombres duaux sur $k$. Soit
$\Spec(k[\epsilon]) \to B$ le composé de $g_0 : \Spec(k) \to B$ et du
morphisme $\Spec(k[\epsilon]) \to \Spec(k)$ provenant de l'inclusion
$k \to k[\epsilon]$. Posons $X_0 = \Spec(k) \times_B X$ et
$X_\epsilon = \Spec(k[\epsilon]) \times_B X$. Observons que $X_\epsilon$ est
un épaississement du premier ordre de $X_0$, plat sur l'épaississement du
premier ordre $\Spec(k) \to \Spec(k[\epsilon])$. En développant les
définitions et en utilisant le lemme
\ref{lemma-coherent-over-first-order-thickening}, nous voyons que
$T\mathcal{F}_{\mathcal{X}_0, k, x_0}$ est l'ensemble des relèvements de
$\mathcal{G}_0$ en un module plat sur $X_\epsilon$. D'après Théorie des
déformations, lemme \ref{defos-lemma-flat-ringed-topoi}, nous concluons que
$$
T\mathcal{F}_{\mathcal{X}_0, k, x_0} =
\Ext^1_{\mathcal{O}_{X_0}}(\mathcal{G}_0, \mathcal{G}_0)
$$
Nous avons utilisé ici l'identification $\epsilon k[\epsilon] \cong k$ de
$k[\epsilon]$-modules. En utilisant une fois encore Théorie des déformations,
lemme \ref{defos-lemma-flat-ringed-topoi}, nous voyons que
$$
\text{Inf}(\mathcal{F}_{\mathcal{X}, k, x_0}) =
\Ext^0_{\mathcal{O}_{X_0}}(\mathcal{G}_0, \mathcal{G}_0)
$$
Ces espaces sont de dimension finie sur $k$, car $\mathcal{G}_0$ est à support
propre sur $\Spec(k)$. En effet, $X_0$ est de présentation finie sur
$\Spec(k)$, donc noethérien. Comme $\mathcal{G}_0$ est de présentation finie,
c'est un $\mathcal{O}_{X_0}$-module cohérent. Nous pouvons donc appliquer
Catégories dérivées des espaces, lemme
\ref{spaces-perfect-lemma-ext-finite}, pour conclure à la finitude voulue.
\end{proof}

\begin{lemma}
\label{lemma-coherent-existence}
Dans la situation \ref{situation-coherent}, supposons que $S$ soit un schéma
localement noethérien et que $f : X \to B$ soit séparé. Soit
$\mathcal{X} = \Cohstack_{X/B}$. Alors le foncteur d'Axiomes d'Artin,
équation (\ref{artin-equation-approximation}), est une équivalence.
\end{lemma}

\begin{proof}
Soit $A$ une $S$-algèbre qui est un anneau local noethérien complet, d'idéal
maximal $\mathfrak m$, et dont le corps résiduel $k$ est de type fini sur $S$.
Nous devons montrer que la catégorie des objets sur $A$ est équivalente à la
catégorie des objets formels sur $A$. Comme nous savons que ceci vaut pour la
catégorie $\mathcal{S}_B$ fibrée en ensembles associée à $B$, d'après Axiomes
d'Artin, lemme \ref{artin-lemma-effective}, il suffit de le démontrer pour les
objets situés au-dessus d'un morphisme donné $\Spec(A) \to B$.

\medskip\noindent
Posons $X_A = \Spec(A) \times_B X$ et
$X_n = \Spec(A/\mathfrak m^n) \times_B X$. D'après le théorème d'existence de
Grothendieck (Plus sur les morphismes d'espaces, théorème
\ref{spaces-more-morphisms-theorem-grothendieck-existence}), la catégorie des
modules cohérents $\mathcal{F}$ sur $X_A$ à support propre sur $\Spec(A)$ est
équivalente à la catégorie des systèmes $(\mathcal{F}_n)$ de modules cohérents
$\mathcal{F}_n$ sur $X_n$ à support propre sur
$\Spec(A/\mathfrak m^n)$. L'équivalence envoie $\mathcal{F}$ sur le système
$(\mathcal{F} \otimes_A A/\mathfrak m^n)$. Voir la discussion dans Plus sur
les morphismes d'espaces, remarque
\ref{spaces-more-morphisms-remark-reformulate-existence-theorem}. Pour achever
la démonstration du lemme, il suffit de montrer que $\mathcal{F}$ est plat sur
$A$ si et seulement si tous les $\mathcal{F} \otimes_A A/\mathfrak m^n$ sont
plats sur $A/\mathfrak m^n$. Cela résulte de Plus sur les morphismes d'espaces,
lemme \ref{spaces-more-morphisms-lemma-flatness-over-Noetherian-ring}.
\end{proof}

\begin{lemma}
\label{lemma-coherent-defo-thy}
Dans la situation \ref{situation-coherent}, supposons que $S$ soit un schéma
localement noethérien, que $S = B$ et que $f : X \to B$ soit plat. Soit
$\mathcal{X} = \Cohstack_{X/B}$. Alors la versalité est ouverte pour
$\mathcal{X}$ (voir Axiomes d'Artin, définition
\ref{artin-definition-openness-versality}).
\end{lemma}

\begin{proof}[Première démonstration]
Cette démonstration repose sur le critère d'Axiomes d'Artin, lemme
\ref{artin-lemma-dual-openness}. Soit $U \to S$ un morphisme de schémas de
type fini, soit $x$ un objet de $\mathcal{X}$ sur $U$ et soit $u_0 \in U$ un
point de type fini tel que $x$ soit versel en $u_0$. En rétrécissant $U$, nous
pouvons supposer que $u_0$ est un point fermé (Morphismes, lemme
\ref{morphisms-lemma-point-finite-type}) et que $U = \Spec(A)$, le morphisme
$U \to S$ étant à valeurs dans un ouvert affine $\Spec(\Lambda)$ de $S$.
Soit $\mathcal{F}$ le module cohérent sur
$X_A = \Spec(A) \times_S X$, plat sur $A$, qui correspond à l'objet donné
$x$.

\medskip\noindent
D'après Théorie des déformations, lemme
\ref{defos-lemma-flat-ringed-topoi}, nous avons un isomorphisme de foncteurs
$$
T_x(M) = \Ext^1_{X_A}(\mathcal{F}, \mathcal{F} \otimes_A M)
$$
et, étant donnée une surjection $A' \to A$ de $\Lambda$-algèbres, de noyau de
carré nul $I$, nous avons une classe d'obstruction
$$
\xi_{A'} \in \Ext^2_{X_A}(\mathcal{F}, \mathcal{F} \otimes_A I)
$$
Ceci utilise le fait que, pour tout $A' \to A$ comme ci-dessus, le changement
de base $X_{A'} = \Spec(A') \times_B X$ est plat sur $A'$. De plus, la
construction de la classe d'obstruction est fonctorielle en la surjection
$A' \to A$ (pour $A$ fixé), d'après Théorie des déformations, lemme
\ref{defos-lemma-functorial-ringed-topoi}. Appliquons Catégories dérivées des
espaces, lemme \ref{spaces-perfect-lemma-compute-ext}, au calcul des groupes
Ext $\Ext^i_{X_A}(\mathcal{F}, \mathcal{F} \otimes_A M)$ pour $i \leq m$,
avec $m = 2$. Nous trouvons un objet parfait $K \in D(A)$ et des
isomorphismes fonctoriels
$$
H^i(K \otimes_A^\mathbf{L} M)
\longrightarrow
\Ext^i_{X_A}(\mathcal{F}, \mathcal{F} \otimes_A M)
$$
pour $i \leq m$, compatibles aux morphismes de bord. Cet objet $K$, avec les
identifications affichées ci-dessus, nous fournit une donnée comme dans
Axiomes d'Artin, situation \ref{artin-situation-dual}. Enfin, la condition (iv)
d'Axiomes d'Artin, lemme \ref{artin-lemma-dual-obstruction}, est satisfaite
d'après Théorie des déformations, lemme
\ref{defos-lemma-verify-iv-ringed-topoi}. Axiomes d'Artin, lemme
\ref{artin-lemma-dual-openness}, s'applique donc bien, ce qui démontre le
lemme.
\end{proof}

\begin{proof}[Deuxième démonstration]
Cette démonstration repose sur Axiomes d'Artin, lemme
\ref{artin-lemma-get-openness-obstruction-theory}. Les conditions (1), (2) et
(3) de ce lemme correspondent aux lemmes \ref{lemma-coherent-diagonal},
\ref{lemma-coherent-RS-star} et \ref{lemma-coherent-limits}.

\medskip\noindent
Nous avons construit une théorie de l'obstruction dans le chapitre sur la
théorie des déformations. Plus précisément, étant données une $S$-algèbre $A$
et un objet $x$ de $\Cohstack_{X/B}$ sur $\Spec(A)$, donné par
$\mathcal{F}$ sur $X_A$, nous posons
$\mathcal{O}_x(M) = \Ext^2_{X_A}(\mathcal{F}, \mathcal{F} \otimes_A M)$ ; si
$A' \to A$ est une surjection de noyau $I$, nous prenons pour élément
d'obstruction l'élément
$$
o_x(A') = o(\mathcal{F}, \mathcal{F} \otimes_A I, 1) \in
\mathcal{O}_x(I) = \Ext^2_{X_A}(\mathcal{F}, \mathcal{F} \otimes_A I)
$$
de Théorie des déformations, lemme \ref{defos-lemma-flat-ringed-topoi}.
Toutes les propriétés d'une théorie de l'obstruction telle qu'elle est définie
dans Axiomes d'Artin, définition
\ref{artin-definition-obstruction-theory}, résultent de ce lemme, sauf la
fonctorialité des classes d'obstruction formulée dans la condition (ii) de la
définition. Mais, comme l'indique la note de bas de page jointe à l'hypothèse
(4) d'Axiomes d'Artin, lemme
\ref{artin-lemma-get-openness-obstruction-theory}, il suffit de vérifier la
fonctorialité des classes d'obstruction pour $A$ fixé, ce qui résulte de
Théorie des déformations, lemme
\ref{defos-lemma-functorial-ringed-topoi}. Théorie des déformations, lemme
\ref{defos-lemma-flat-ringed-topoi}, nous dit aussi que
$T_x(M) = \Ext^1_{X_A}(\mathcal{F}, \mathcal{F} \otimes_A M)$ pour tout
$A$-module $M$.

\medskip\noindent
Pour achever la démonstration, il suffit de montrer que
$T_x(\prod M_n) = \prod T_x(M_n)$ et
$\mathcal{O}_x(\prod M_n) = \prod \mathcal{O}_x(M)$. Appliquons Catégories
dérivées des espaces, lemme \ref{spaces-perfect-lemma-compute-ext}, au calcul
des groupes Ext $\Ext^i_{X_A}(\mathcal{F}, \mathcal{F} \otimes_A M)$ pour
$i \leq m$, avec $m = 2$. Nous trouvons un objet parfait $K \in D(A)$ et des
isomorphismes fonctoriels
$$
H^i(K \otimes_A^\mathbf{L} M)
\longrightarrow
\Ext^i_{X_A}(\mathcal{F}, \mathcal{F} \otimes_A M)
$$
pour $i = 1, 2$. Un argument immédiat montre que
$$
H^i(K \otimes_A^\mathbf{L} \prod M_n) =
\prod H^i(K \otimes_A^\mathbf{L} M_n)
$$
dès que $K$ est un objet pseudo-cohérent de $D(A)$. En fait, cette propriété
(pour tout $i$) caractérise les complexes pseudo-cohérents ; voir Plus
d'algèbre, lemme \ref{more-algebra-lemma-pseudo-coherent-tensor}.
\end{proof}

\begin{theorem}[Algébricité du champ des faisceaux cohérents ; cas plat]
\label{theorem-coherent-algebraic}
Soit $S$ un schéma. Soit $f : X \to B$ un morphisme d'espaces algébriques sur
$S$. Supposons que $f$ soit de présentation finie, séparé et
plat\footnote{Cette hypothèse n'est pas nécessaire. Voir la section
\ref{section-not-flat}.}. Alors $\Cohstack_{X/B}$ est un champ algébrique sur
$S$.
\end{theorem}

\begin{proof}
Posons $\mathcal{X} = \Cohstack_{X/B}$. Nous avons vu que $\mathcal{X}$ est un
champ en groupoïdes sur $(\Sch/S)_{fppf}$, à diagonale représentable par des
espaces algébriques (lemmes \ref{lemma-coherent-stack} et
\ref{lemma-coherent-diagonal}). Il suffit donc de trouver un schéma $W$ et un
morphisme lisse surjectif $W \to \mathcal{X}$.

\medskip\noindent
Soit $B'$ un schéma et soit $B' \to B$ un morphisme étale surjectif. Posons
$X' = B' \times_B X$ et désignons par $f' : X' \to B'$ la projection. Alors
$\mathcal{X}' = \Cohstack_{X'/B'}$ est égal au $2$-produit fibré de
$\mathcal{X}$ et de la catégorie fibrée en ensembles associée à $B'$, sur la
catégorie fibrée en ensembles associée à $B$ (remarque
\ref{remark-coherent-base-change}). D'après les résultats de Champs
algébriques, section \ref{algebraic-section-representable-properties}, le
morphisme $\mathcal{X}' \to \mathcal{X}$ est étale et surjectif. Il suffit donc
de démontrer le résultat pour $\mathcal{X}'$. Autrement dit, nous pouvons
supposer que $B$ est un schéma.

\medskip\noindent
Supposons que $B$ soit un schéma. Dans ce cas, nous pouvons remplacer $S$ par
$B$ ; voir Champs algébriques, section
\ref{algebraic-section-change-base-scheme}. Nous pouvons donc supposer
$S = B$.

\medskip\noindent
Supposons $S = B$. Choisissons un recouvrement ouvert affine
$S = \bigcup U_i$. Désignons par $\mathcal{X}_i$ la restriction de
$\mathcal{X}$ à $(\Sch/U_i)_{fppf}$. Si nous pouvons trouver des schémas $W_i$
sur $U_i$ et des morphismes lisses surjectifs
$W_i \to \mathcal{X}_i$, alors, en posant $W = \coprod W_i$, nous obtenons un
morphisme lisse surjectif $W \to \mathcal{X}$. Nous pouvons donc supposer que
$S = B$ est affine.

\medskip\noindent
Supposons $S = B$ affine, disons $S = \Spec(\Lambda)$. Écrivons
$\Lambda = \colim \Lambda_i$ comme une colimite filtrante, chaque $\Lambda_i$
étant de type fini sur $\mathbf{Z}$. Pour un certain $i$, nous pouvons trouver
un morphisme d'espaces algébriques $X_i \to \Spec(\Lambda_i)$ qui soit de
présentation finie, séparé et plat, et dont le changement de base à $\Lambda$
soit $X$. Voir Limites d'espaces, lemmes
\ref{spaces-limits-lemma-descend-finite-presentation},
\ref{spaces-limits-lemma-descend-separated-morphism} et
\ref{spaces-limits-lemma-descend-flat}. Si nous montrons que
$\Cohstack_{X_i/\Spec(\Lambda_i)}$ est un champ algébrique, il résultera par
changement de base (remarque \ref{remark-coherent-base-change} et Champs
algébriques, section \ref{algebraic-section-change-base-scheme}) que
$\mathcal{X}$ est un champ algébrique. Nous pouvons donc supposer que
$\Lambda$ est une $\mathbf{Z}$-algèbre de type fini.

\medskip\noindent
Supposons $S = B = \Spec(\Lambda)$ affine de type fini sur $\mathbf{Z}$. Dans
ce cas, nous vérifierons les conditions (1), (2), (3), (4) et (5) d'Axiomes
d'Artin, lemme \ref{artin-lemma-diagonal-representable}, pour conclure que
$\mathcal{X}$ est un champ algébrique. Remarquons que $\Lambda$ est un
anneau G ; voir Plus d'algèbre, proposition
\ref{more-algebra-proposition-ubiquity-G-ring}. Ainsi, tous les anneaux locaux
de $S$ sont des anneaux G. La condition (5) est donc satisfaite. D'après le
lemme \ref{lemma-coherent-defo-thy}, la versalité est ouverte pour
$\mathcal{X}$, donc (4) est satisfaite. Pour vérifier (2), nous devons vérifier
les axiomes [-1], [0], [1], [2] et [3] d'Axiomes d'Artin, section
\ref{artin-section-axioms}. Nous omettons la vérification de [-1], et les
axiomes [0], [1], [2], [3] correspondent respectivement aux lemmes
\ref{lemma-coherent-stack}, \ref{lemma-coherent-limits},
\ref{lemma-coherent-RS-star}, \ref{lemma-coherent-tangent-space}. La condition
(3) résulte du lemme \ref{lemma-coherent-existence}. Enfin, la condition (1)
est le lemme \ref{lemma-coherent-diagonal}. Ceci achève la démonstration du
théorème.
\end{proof}









\section{Le champ des faisceaux cohérents dans le cas non plat}
\label{section-not-flat}

\noindent
Dans le théorème \ref{theorem-coherent-algebraic}, l'hypothèse que
$f : X \to B$ est plat n'est pas nécessaire. Dans cette section, nous donnons
une démonstration différente qui évite l'hypothèse de platitude et la
vérification de l'ouverture de la versalité, en utilisant les résultats de
Platitude sur les espaces, section \ref{spaces-flat-section-existence}, et
d'Axiomes d'Artin, section
\ref{artin-section-strong-formal-effectiveness}.

\medskip\noindent
Pour une autre approche de ce problème, le lecteur pourra consulter
\cite{ArtinI} et suivre la méthode exposée dans les articles
\cite{olsson-starr}, \cite{lieblich_remarks}, \cite{olsson_proper},
\cite{Hall-Rydh}, \cite{Hall-Rydh-Hilbert}, \cite{rydh_representability}.
Certains de ces articles traitent du cas plus général du champ des faisceaux
cohérents sur un champ algébrique au-dessus d'un champ algébrique, et d'autres
traitent de problèmes analogues pour les champs de Hilbert ou les foncteurs de
Quot. Notre stratégie consistera à démontrer l'algébricité de certains champs
de Hilbert et foncteurs de Quot comme conséquence de l'algébricité du champ des
faisceaux cohérents.

\begin{theorem}[Algébricité du champ des faisceaux cohérents ; cas général]
\label{theorem-coherent-algebraic-general}
Soit $S$ un schéma. Soit $f : X \to B$ un morphisme d'espaces algébriques sur
$S$. Supposons que $f$ soit de présentation finie et séparé. Alors
$\Cohstack_{X/B}$ est un champ algébrique sur $S$.
\end{theorem}

\begin{proof}
Seule la dernière étape diffère de la démonstration du cas plat, mais nous
répétons ici tous les arguments afin de nous assurer que tout fonctionne.

\medskip\noindent
Posons $\mathcal{X} = \Cohstack_{X/B}$. Nous avons vu que $\mathcal{X}$ est un
champ en groupoïdes sur $(\Sch/S)_{fppf}$, à diagonale représentable par des
espaces algébriques (lemmes \ref{lemma-coherent-stack} et
\ref{lemma-coherent-diagonal}). Il suffit donc de trouver un schéma $W$ et un
morphisme lisse surjectif $W \to \mathcal{X}$.

\medskip\noindent
Soit $B'$ un schéma et soit $B' \to B$ un morphisme étale surjectif. Posons
$X' = B' \times_B X$ et désignons par $f' : X' \to B'$ la projection. Alors
$\mathcal{X}' = \Cohstack_{X'/B'}$ est égal au $2$-produit fibré de
$\mathcal{X}$ et de la catégorie fibrée en ensembles associée à $B'$, sur la
catégorie fibrée en ensembles associée à $B$ (remarque
\ref{remark-coherent-base-change}). D'après les résultats de Champs
algébriques, section \ref{algebraic-section-representable-properties}, le
morphisme $\mathcal{X}' \to \mathcal{X}$ est étale et surjectif. Il suffit donc
de démontrer le résultat pour $\mathcal{X}'$. Autrement dit, nous pouvons
supposer que $B$ est un schéma.

\medskip\noindent
Supposons que $B$ soit un schéma. Dans ce cas, nous pouvons remplacer $S$ par
$B$ ; voir Champs algébriques, section
\ref{algebraic-section-change-base-scheme}. Nous pouvons donc supposer
$S = B$.

\medskip\noindent
Supposons $S = B$. Choisissons un recouvrement ouvert affine
$S = \bigcup U_i$. Désignons par $\mathcal{X}_i$ la restriction de
$\mathcal{X}$ à $(\Sch/U_i)_{fppf}$. Si nous pouvons trouver des schémas $W_i$
sur $U_i$ et des morphismes lisses surjectifs
$W_i \to \mathcal{X}_i$, alors, en posant $W = \coprod W_i$, nous obtenons un
morphisme lisse surjectif $W \to \mathcal{X}$. Nous pouvons donc supposer que
$S = B$ est affine.

\medskip\noindent
Supposons $S = B$ affine, disons $S = \Spec(\Lambda)$. Écrivons
$\Lambda = \colim \Lambda_i$ comme une colimite filtrante, chaque $\Lambda_i$
étant de type fini sur $\mathbf{Z}$. Pour un certain $i$, nous pouvons trouver
un morphisme d'espaces algébriques $X_i \to \Spec(\Lambda_i)$ qui soit séparé
et de présentation finie et dont le changement de base à $\Lambda$ soit $X$.
Voir Limites d'espaces, lemmes
\ref{spaces-limits-lemma-descend-finite-presentation} et
\ref{spaces-limits-lemma-descend-separated-morphism}. Si nous montrons que
$\Cohstack_{X_i/\Spec(\Lambda_i)}$ est un champ algébrique, il résultera par
changement de base (remarque \ref{remark-coherent-base-change} et Champs
algébriques, section \ref{algebraic-section-change-base-scheme}) que
$\mathcal{X}$ est un champ algébrique. Nous pouvons donc supposer que
$\Lambda$ est une $\mathbf{Z}$-algèbre de type fini.

\medskip\noindent
Supposons $S = B = \Spec(\Lambda)$ affine de type fini sur $\mathbf{Z}$. Dans
ce cas, nous vérifierons les conditions (1), (2), (3), (4) et (5) d'Axiomes
d'Artin, lemme \ref{artin-lemma-diagonal-representable}, pour conclure que
$\mathcal{X}$ est un champ algébrique. Remarquons que $\Lambda$ est un anneau
G ; voir Plus d'algèbre, proposition
\ref{more-algebra-proposition-ubiquity-G-ring}. Ainsi, tous les anneaux locaux
de $S$ sont des anneaux G. La condition (5) est donc satisfaite. Pour vérifier
(2), nous devons vérifier les axiomes [-1], [0], [1], [2] et [3] d'Axiomes
d'Artin, section \ref{artin-section-axioms}. Nous omettons la vérification de
[-1], et les axiomes [0], [1], [2], [3] correspondent respectivement aux
lemmes \ref{lemma-coherent-stack}, \ref{lemma-coherent-limits},
\ref{lemma-coherent-RS-star}, \ref{lemma-coherent-tangent-space}. La condition
(3) est le lemme \ref{lemma-coherent-existence}. La condition (1) est le lemme
\ref{lemma-coherent-diagonal}.

\medskip\noindent
Il reste à démontrer la condition (4), qui est l'ouverture de la versalité.
Pour cela, nous utiliserons Axiomes d'Artin, lemme
\ref{artin-lemma-SGE-implies-openness-versality}. Nous avons déjà vu que
$\mathcal{X}$ a sa diagonale représentable par des espaces algébriques,
possède (RS*) et commute aux limites (voir les lemmes utilisés ci-dessus). Il
nous reste donc seulement à voir que $\mathcal{X}$ satisfait à l'effectivité
formelle forte formulée dans Axiomes d'Artin, lemme
\ref{artin-lemma-SGE-implies-openness-versality}. C'est le théorème de
Platitude sur les espaces \ref{spaces-flat-theorem-existence}, et la
démonstration est achevée.
\end{proof}

\section{Le foncteur des quotients}
\label{section-functor-quotients}

\noindent
Dans cette section, nous donnons quelques généralités sur le foncteur
$Q_{\mathcal{F}/X/B}$ défini ci-dessous. La notation
$\Quotfunctor_{\mathcal{F}/X/B}$ est réservée à un sous-foncteur de
$\text{Q}_{\mathcal{F}/X/B}$. Nous recommandons au lecteur d'omettre cette
section lors d'une première lecture.

\begin{situation}
\label{situation-q}
Soit $S$ un schéma. Soit $f : X \to B$ un morphisme d'espaces algébriques sur
$S$. Soit $\mathcal{F}$ un $\mathcal{O}_X$-module quasi-cohérent. Pour tout
schéma $T$ sur $B$, nous désignerons par $X_T$ le changement de base de $X$ à
$T$ et par $\mathcal{F}_T$ l'image inverse de $\mathcal{F}$ par le morphisme
de projection $X_T = X \times_B T \to X$. Pour un tel $T$, posons
$$
\text{Q}_{\mathcal{F}/X/B}(T) =
\left\{
\begin{matrix}
\text{quotients }\mathcal{F}_T \to \mathcal{Q}\text{ où }
\mathcal{Q}\text{ est un}\\
\mathcal{O}_{X_T}\text{-module quasi-cohérent plat sur }T
\end{matrix}
\right\}
$$
Nous identifions deux quotients s'ils ont le même noyau. Supposons que
$T' \to T$ soit un morphisme de schémas sur $B$ et que
$\mathcal{F}_T \to \mathcal{Q}$ soit un élément de
$\text{Q}_{\mathcal{F}/X/B}(T)$. Alors l'image inverse
$\mathcal{Q}' = (X_{T'} \to X_T)^*\mathcal{Q}$ est un
$\mathcal{O}_{X_{T'}}$-module quasi-cohérent plat sur $T'$, d'après Morphismes
d'espaces, lemme \ref{spaces-morphisms-lemma-base-change-module-flat}. Nous
obtenons donc un foncteur
\begin{equation}
\label{equation-q}
\text{Q}_{\mathcal{F}/X/B} : (\Sch/B)^{opp} \longrightarrow \textit{Sets}
\end{equation}
C'est le {\it foncteur des quotients de $\mathcal{F}/X/B$}. Nous définissons
un sous-foncteur
\begin{equation}
\label{equation-q-fp}
\text{Q}^{fp}_{\mathcal{F}/X/B} : (\Sch/B)^{opp} \longrightarrow \textit{Sets}
\end{equation}
qui associe à $T$ le sous-ensemble de $\text{Q}_{\mathcal{F}/X/B}(T)$ formé
des quotients $\mathcal{F}_T \to \mathcal{Q}$ tels que $\mathcal{Q}$ soit de
présentation finie comme $\mathcal{O}_{X_T}$-module. C'est un sous-foncteur
d'après Propriétés des espaces, section
\ref{spaces-properties-section-properties-modules}.
\end{situation}

\noindent
Dans la situation \ref{situation-q}, nous regardons parfois
$\text{Q}_{\mathcal{F}/X/B}$ comme un foncteur
$(\Sch/S)^{opp} \to \textit{Sets}$ muni d'un morphisme
$\text{Q}_{\mathcal{F}/X/S} \to B$. En effet, si $T$ est un schéma sur $S$,
un élément de $\text{Q}_{\mathcal{F}/X/B}(T)$ est un couple
$(h, \mathcal{Q})$, où $h$ un morphisme $h : T \to B$ et où $\mathcal{Q}$
est un quotient $\mathcal{F}_T \to \mathcal{Q}$ plat sur $T$, de présentation
finie sur $X_T = X \times_{B, h} T$. En particulier, lorsque nous disons que
$\text{Q}_{\mathcal{F}/X/S}$ est un espace algébrique, nous entendons que le
foncteur correspondant $(\Sch/S)^{opp} \to \textit{Sets}$ est un espace
algébrique. Des remarques analogues valent pour
$\text{Q}^{fp}_{\mathcal{F}/X/B}$.

\begin{remark}
\label{remark-q-base-change}
Dans la situation \ref{situation-q}, soit $B' \to B$ un morphisme d'espaces
algébriques sur $S$. Posons $X' = X \times_B B'$ et désignons par
$\mathcal{F}'$ l'image inverse de $\mathcal{F}$ sur $X'$. Nous avons donc le
foncteur $Q_{\mathcal{F}'/X'/B'}$ sur la catégorie des schémas sur $B'$. Pour
un schéma $T$ sur $B'$, il est clair que
$$
Q_{\mathcal{F}'/X'/B'}(T) = Q_{\mathcal{F}/X/B}(T)
$$
où, dans le membre de droite, nous regardons $T$ comme un schéma sur $B$ au
moyen du composé $T \to B' \to B$. Des remarques analogues valent pour
$\text{Q}^{fp}_{\mathcal{F}/X/B}$. Ces remarques élémentaires seront parfois
utiles pour changer l'espace algébrique de base.
\end{remark}

\begin{remark}
\label{remark-q-sheaf}
Soient $S$ un schéma, $X$ un espace algébrique sur $S$ et $\mathcal{F}$ un
$\mathcal{O}_X$-module quasi-cohérent. Supposons que
$\{f_i : X_i \to X\}_{i \in I}$ soit un recouvrement fpqc et que, pour chaque
$i, j \in I$, nous soit donné un recouvrement fpqc
$\{X_{ijk} \to X_i \times_X X_j\}$. Dans cette situation, nous avons une
bijection
$$
\left\{
\begin{matrix}
\text{quotients }\mathcal{F} \to \mathcal{Q}\text{ où } \\
\mathcal{Q}\text{ est quasi-cohérent}\\
\end{matrix}
\right\}
\longrightarrow
\left\{
\begin{matrix}
\text{familles de quotients }f_i^*\mathcal{F} \to \mathcal{Q}_i
\text{ où } \\
\mathcal{Q}_i\text{ est quasi-cohérent et où }
\mathcal{Q}_i\text{ et }\mathcal{Q}_j\\
\text{induisent le même quotient sur }X_{ijk}
\end{matrix}
\right\}
$$
En effet, soit $(f_i^*\mathcal{F} \to \mathcal{Q}_i)_{i \in I}$ un élément du
membre de droite. Comme $\{X_{ijk} \to X_i \times_X X_j\}$ est un
recouvrement fpqc, nous voyons que les images inverses de $\mathcal{Q}_i$ et
$\mathcal{Q}_j$ induisent le même quotient de l'image inverse de
$\mathcal{F}$ sur $X_i \times_X X_j$ (par pleine fidélité dans Descente sur
les espaces, proposition
\ref{spaces-descent-proposition-fpqc-descent-quasi-coherent}). Nous obtenons
donc une donnée de descente pour les modules quasi-cohérents relativement à
$\{X_i \to X\}_{i \in I}$. D'après Descente sur les espaces, proposition
\ref{spaces-descent-proposition-fpqc-descent-quasi-coherent}, nous trouvons un
homomorphisme de $\mathcal{O}_X$-modules quasi-cohérents
$\mathcal{F} \to \mathcal{Q}$ dont la restriction à $X_i$ redonne les
homomorphismes donnés $f_i^*\mathcal{F} \to \mathcal{Q}_i$. Comme la famille
de morphismes $\{X_i \to X\}$ est conjointement surjective et plate, pour
tout point $x \in |X|$, il existe un $i$ et un point $x_i \in |X_i|$
s'envoyant sur $x$. Remarquons que l'homomorphisme induit sur les anneaux
locaux
$\mathcal{O}_{X, \overline{x}} \to \mathcal{O}_{X_i, \overline{x_i}}$
est fidèlement plat ; voir Morphismes d'espaces, section
\ref{spaces-morphisms-section-flat}. Nous voyons donc que
$\mathcal{F} \to \mathcal{Q}$ est surjectif.
\end{remark}

\begin{lemma}
\label{lemma-q-sheaf}
Dans la situation \ref{situation-q}. Les foncteurs
$\text{Q}_{\mathcal{F}/X/B}$ et $\text{Q}^{fp}_{\mathcal{F}/X/B}$ satisfont à
la propriété de faisceau pour la topologie fpqc.
\end{lemma}

\begin{proof}
Soit $\{T_i \to T\}_{i \in I}$ un recouvrement fpqc de schémas sur $S$.
Posons $X_i = X_{T_i} = X \times_S T_i$ et
$\mathcal{F}_i = \mathcal{F}_{T_i}$. Remarquons que
$\{X_i \to X_T\}_{i \in I}$ est un recouvrement fpqc de $X_T$ (Topologies
sur les espaces, lemme \ref{spaces-topologies-lemma-fpqc}) et que
$X_{T_i \times_T T_{i'}} = X_i \times_{X_T} X_{i'}$. Supposons que
$\mathcal{F}_i \to \mathcal{Q}_i$ soit une famille d'éléments de
$\text{Q}_{\mathcal{F}/X/B}(T_i)$ telle que $\mathcal{Q}_i$ et
$\mathcal{Q}_{i'}$ induisent le même élément de
$\text{Q}_{\mathcal{F}/X/B}(T_i \times_T T_{i'})$. D'après la remarque
\ref{remark-q-sheaf}, nous obtenons une application surjective de
$\mathcal{O}_{X_T}$-modules quasi-cohérents
$\mathcal{F}_T \to \mathcal{Q}$ dont la restriction à $X_i$ redonne les
quotients donnés. D'après Morphismes d'espaces, lemme
\ref{spaces-morphisms-lemma-flat-permanence}, nous voyons que $\mathcal{Q}$ est
plat sur $T$. Enfin, dans le cas de $\text{Q}^{fp}_{\mathcal{F}/X/B}$,
c'est-à-dire si les $\mathcal{Q}_i$ sont de présentation finie, Descente sur
les espaces, lemme \ref{spaces-descent-lemma-finite-presentation-descends},
garantit que $\mathcal{Q}$ est de présentation finie comme
$\mathcal{O}_{X_T}$-module.
\end{proof}

\noindent
Vérification de cohérence : $\text{Q}_{\mathcal{F}/X/B}$ et
$\text{Q}^{fp}_{\mathcal{F}/X/B}$ jouent le même rôle parmi les espaces
algébriques sur $S$.

\begin{lemma}
\label{lemma-extend-q-to-spaces}
Dans la situation \ref{situation-q}. Soit $T$ un espace algébrique sur $S$.
Nous avons
$$
\Mor_{\Sh((\Sch/S)_{fppf})}(T,  \text{Q}_{\mathcal{F}/X/B}) =
\left\{
\begin{matrix}
(h, \mathcal{F}_T \to \mathcal{Q}) \text{ où }
h : T \to B \text{ et où}\\
\mathcal{Q}\text{ est quasi-cohérent et plat sur }T
\end{matrix}
\right\}
$$
où $\mathcal{F}_T$ désigne l'image inverse de $\mathcal{F}$ sur l'espace
algébrique $X \times_{B, h} T$. De même, nous avons
$$
\Mor_{\Sh((\Sch/S)_{fppf})}(T,  \text{Q}^{fp}_{\mathcal{F}/X/B}) =
\left\{
\begin{matrix}
(h, \mathcal{F}_T \to \mathcal{Q}) \text{ où }
h : T \to B \text{ et où}\\
\mathcal{Q}\text{ est de présentation finie et plat sur }T
\end{matrix}
\right\}
$$
\end{lemma}

\begin{proof}
Choisissons un schéma $U$ et un morphisme étale surjectif $p : U \to T$.
Soit $R = U \times_T U$, de projections $t, s : R \to U$.

\medskip\noindent
Soit $v : T \to \text{Q}_{\mathcal{F}/X/B}$ une transformation naturelle.
Alors $v(p)$ correspond à un couple
$(h_U, \mathcal{F}_U \to \mathcal{Q}_U)$ sur $U$. Comme $v$ est une
transformation de foncteurs, les images inverses de
$(h_U, \mathcal{F}_U \to \mathcal{Q}_U)$ par $s$ et $t$ coïncident. Puisque
$T = U/R$ (Espaces, lemme \ref{spaces-lemma-space-presentation}), nous
obtenons un morphisme $h : T \to B$ tel que $h_U = h \circ p$. D'après
Descente sur les espaces, proposition
\ref{spaces-descent-proposition-fpqc-descent-quasi-coherent}, le quotient
$\mathcal{Q}_U$ descend en un quotient $\mathcal{F}_T \to \mathcal{Q}$ sur
$X_T$. Comme $U \to T$ est surjectif et plat, il résulte de Morphismes
d'espaces, lemme \ref{spaces-morphisms-lemma-flat-permanence}, que
$\mathcal{Q}$ est plat sur $T$.

\medskip\noindent
Réciproquement, soit $(h, \mathcal{F}_T \to \mathcal{Q})$ un couple sur $T$.
Nous obtenons alors une transformation naturelle
$v : T \to \text{Q}_{\mathcal{F}/X/B}$ en envoyant un morphisme
$a : T' \to T$, où $T'$ est un schéma, sur
$(h \circ a, \mathcal{F}_{T'} \to a^*\mathcal{Q})$. Nous omettons de vérifier
que cette construction et celle du paragraphe précédent sont inverses l'une
de l'autre.

\medskip\noindent
Dans le cas de $\text{Q}^{fp}_{\mathcal{F}/X/B}$, ajoutons ceci : étant donné
un morphisme $h : T \to B$, un faisceau quasi-cohérent sur $X_T$ est de
présentation finie comme $\mathcal{O}_{X_T}$-module si et seulement si son
image inverse sur $X_U$ est de présentation finie comme
$\mathcal{O}_{X_U}$-module. Cela résulte du fait que $X_U \to X_T$ est
surjectif et étale, et de Descente sur les espaces, lemme
\ref{spaces-descent-lemma-finite-presentation-descends}.
\end{proof}

\begin{lemma}
\label{lemma-q-sheaf-in-X}
Dans la situation \ref{situation-q}, soit $\{X_i \to X\}_{i \in I}$ un
recouvrement fpqc et, pour chaque $i, j \in I$, soit
$\{X_{ijk} \to X_i \times_X X_j\}$ un recouvrement fpqc. Désignons par
$\mathcal{F}_i$, resp.\ $\mathcal{F}_{ijk}$, l'image inverse de
$\mathcal{F}$ sur $X_i$, resp.\ sur $X_{ijk}$. Pour tout schéma $T$ sur $B$,
le diagramme
$$
\xymatrix{
Q_{\mathcal{F}/X/B}(T) \ar[r] &
\prod\nolimits_i
Q_{\mathcal{F}_i/X_i/B}(T)
\ar@<1ex>[r]^-{\text{pr}_0^*} \ar@<-1ex>[r]_-{\text{pr}_1^*}
&
\prod\nolimits_{i, j, k}
Q_{\mathcal{F}_{ijk}/X_{ijk}/B}(T)
}
$$
présente la première flèche comme l'égalisateur des deux autres. La même
assertion vaut pour le foncteur $\text{Q}^{fp}_{\mathcal{F}/X/B}$.
\end{lemma}

\begin{proof}
Soit $\mathcal{F}_{i, T} \to \mathcal{Q}_i$ un élément de l'égalisateur de
$\text{pr}_0^*$ et $\text{pr}_1^*$. D'après la remarque
\ref{remark-q-sheaf}, nous obtenons une surjection
$\mathcal{F}_T \to \mathcal{Q}$ de $\mathcal{O}_{X_T}$-modules
quasi-cohérents, dont la restriction à $X_{i, T}$ redonne
$\mathcal{F}_i \to \mathcal{Q}_i$. D'après Morphismes d'espaces, lemme
\ref{spaces-morphisms-lemma-flat-permanence}, nous voyons que $\mathcal{Q}$ est
plat sur $T$, comme voulu. Dans le cas du foncteur
$\text{Q}^{fp}_{\mathcal{F}/X/B}$, c'est-à-dire si $\mathcal{Q}_i$ est de
présentation finie, $\mathcal{Q}$ est lui aussi de présentation finie d'après
Descente sur les espaces, lemme
\ref{spaces-descent-lemma-finite-presentation-descends}.
\end{proof}

\begin{lemma}
\label{lemma-q-limit-preserving}
Dans la situation \ref{situation-q}, supposons en outre que (a) $f$ soit
quasi-compact et quasi-séparé, et que (b) $\mathcal{F}$ soit de présentation
finie. Alors le foncteur $\text{Q}^{fp}_{\mathcal{F}/X/B}$ commute aux limites
au sens suivant : si $T = \lim T_i$ est une limite filtrante de schémas
affines sur $B$, alors
$\text{Q}^{fp}_{\mathcal{F}/X/B}(T) =
\colim \text{Q}^{fp}_{\mathcal{F}/X/B}(T_i)$.
\end{lemma}

\begin{proof}
Soit $T = \lim T_i$ comme dans l'énoncé du lemme. Choisissons $i_0 \in I$ et
remplaçons $I$ par $\{i \in I \mid i \geq i_0\}$. Nous pouvons poser
$B = S = T_{i_0}$ et remplacer $X$ par $X_{T_0}$ et $\mathcal{F}$ par son
image inverse sur $X_{T_0}$. Alors $X_T = \lim X_{T_i}$ ; voir Limites
d'espaces, lemme
\ref{spaces-limits-lemma-directed-inverse-system-has-limit}. Soit
$\mathcal{F}_T \to \mathcal{Q}$ un élément de
$\text{Q}^{fp}_{\mathcal{F}/X/B}(T)$. D'après Limites d'espaces, lemme
\ref{spaces-limits-lemma-descend-modules-finite-presentation}, il existe un
$i$ et une application $\mathcal{F}_{T_i} \to \mathcal{Q}_i$ de
$\mathcal{O}_{X_{T_i}}$-modules de présentation finie dont l'image inverse sur
$X_T$ est l'application quotient donnée.

\medskip\noindent
Il nous reste à vérifier que, quitte à augmenter $i$, l'application
$\mathcal{F}_{T_i} \to \mathcal{Q}_i$ est surjective et $\mathcal{Q}_i$ est
plat sur $T_i$. Pour cela, choisissons un schéma affine $U$ et un morphisme
étale surjectif $U \to X$ (voir Propriétés des espaces, lemme
\ref{spaces-properties-lemma-quasi-compact-affine-cover}). Nous pouvons
vérifier la surjectivité et la platitude sur $T_i$ après image inverse par le
recouvrement étale $U_{T_i} \to X_{T_i}$ (par définition). Nous sommes ainsi
ramenés au cas où $X = \Spec(B_0)$ est un schéma affine de présentation finie
sur $B = S = T_0 = \Spec(A_0)$. En écrivant $T_i = \Spec(A_i)$, puis
$T = \Spec(A)$ avec $A = \colim A_i$, nous sommes parvenus au problème
d'algèbre suivant. Soit $M_i \to N_i$ une application de
$B_0 \otimes_{A_0} A_i$-modules de présentation finie telle que
$M_i \otimes_{A_i} A \to N_i \otimes_{A_i} A$ soit surjective et que
$N_i \otimes_{A_i} A$ soit plat sur $A$. Montrer qu'il existe
$i' \geq i$ tel que
$M_i \otimes_{A_i} A_{i'} \to N_i \otimes_{A_i} A_{i'}$ soit surjective et
que $N_i \otimes_{A_i} A_{i'}$ soit plat sur $A$. La première assertion
résulte d'Algèbre, lemme
\ref{algebra-lemma-module-map-property-in-colimit}, et la seconde d'Algèbre,
lemme \ref{algebra-lemma-flat-finite-presentation-limit-flat}.
\end{proof}

\begin{lemma}
\label{lemma-q-RS-star}
Dans la situation \ref{situation-q}. Soit
$$
\xymatrix{
Z \ar[r] \ar[d] & Z' \ar[d] \\
Y \ar[r] & Y'
}
$$
une somme amalgamée dans la catégorie des schémas sur $B$, où
$Z \to Z'$ est un épaississement et $Z \to Y$ est affine ; voir Plus sur les
morphismes, lemme
\ref{more-morphisms-lemma-pushout-along-thickening}. Alors l'application
naturelle
$$
Q_{\mathcal{F}/X/B}(Y') \longrightarrow
Q_{\mathcal{F}/X/B}(Y) \times_{Q_{\mathcal{F}/X/B}(Z)} Q_{\mathcal{F}/X/B}(Z')
$$
est bijective. Si $X \to B$ est localement de présentation finie, la même
assertion vaut pour $Q^{fp}_{\mathcal{F}/X/B}$.
\end{lemma}

\begin{proof}
Construisons une application inverse. Supposons donc donnés
$\mathcal{F}_Y \to \mathcal{A}$,
$\mathcal{F}_{Z'} \to \mathcal{B}'$ et un isomorphisme
$\mathcal{A}|_{X_Z} \to \mathcal{B}'|_{X_Z}$ compatible aux surjections
données. Nous appliquons alors Sommes amalgamées d'espaces, lemme
\ref{spaces-pushouts-lemma-space-over-pushout-flat-modules}, pour obtenir sur
$X_{Y'}$ un module quasi-cohérent $\mathcal{A}'$ plat sur $Y'$. Comme ce
faisceau est construit comme un produit fibré (voir la démonstration du lemme
cité), il existe une application canonique
$\mathcal{F}_{Y'} \to \mathcal{A}'$. On peut voir que cette application est
surjective parce qu'elle se factorise en
$$
\begin{matrix}
\mathcal{F}_{Y'} \\
\downarrow \\
(X_Y \to X_{Y'})_*\mathcal{F}_Y
\times_{(X_Z \to X_{Y'})_*\mathcal{F}_Z}
(X_{Z'} \to X_{Y'})_*\mathcal{F}_{Z'} \\
\downarrow \\
\mathcal{A}' =
(X_Y \to X_{Y'})_*\mathcal{A}
\times_{(X_Z \to X_{Y'})_*\mathcal{A}|_{X_Z}}
(X_{Z'} \to X_{Y'})_*\mathcal{B}'
\end{matrix}
$$
la première flèche étant surjective d'après Plus d'algèbre, lemme
\ref{more-algebra-lemma-module-over-fibre-product-bis}, et la seconde d'après
Plus d'algèbre, lemme
\ref{more-algebra-lemma-surjection-module-over-fibre-product}.

\medskip\noindent
Dans le cas de $Q^{fp}_{\mathcal{F}/X/B}$, il nous suffit de montrer que la
construction ci-dessus produit un module de présentation finie. Cela est
expliqué dans Plus d'algèbre, remarque
\ref{more-algebra-remark-relative-modules-over-fibre-product}, dans le cadre de
l'algèbre commutative. Le présent cas des modules sur des espaces algébriques
s'en déduit par localisation étale.
\end{proof}

\begin{remark}[Obstructions aux quotients]
\label{remark-q-obs}
Dans la situation \ref{situation-q}, {\bf supposons} que $\mathcal{F}$ soit
plat sur $B$. Soit $T \subset T'$ un épaississement du premier ordre de
schémas sur $B$, de faisceau d'idéaux $\mathcal{J}$. Alors
$X_T \subset X_{T'}$ est un épaississement du premier ordre d'espaces
algébriques, dont le faisceau d'idéaux $\mathcal{I}$ est un quotient de
$f_T^*\mathcal{J}$. Nous considérerons les faisceaux sur $X_{T'}$, resp.\ sur
$T'$, comme des faisceaux sur $X_T$, resp.\ sur $T$, au moyen de l'équivalence
fondamentale décrite dans Plus sur les morphismes d'espaces, section
\ref{spaces-more-morphisms-section-thickenings}. Soit
$$
0 \to \mathcal{K} \to \mathcal{F}_T \to \mathcal{Q} \to 0
$$
une suite définissant un élément $x$ de $Q_{\mathcal{F}/X/B}(T)$. Comme
$\mathcal{F}_{T'}$ est plat sur $T'$, nous avons une suite exacte courte
$$
0 \to f_T^*\mathcal{J} \otimes_{\mathcal{O}_{X_T}} \mathcal{F}_T
\xrightarrow{i} \mathcal{F}_{T'} \xrightarrow{\pi} \mathcal{F}_T \to 0
$$
et nous avons
$f_T^*\mathcal{J} \otimes_{\mathcal{O}_{X_T}} \mathcal{F}_T =
\mathcal{I} \otimes_{\mathcal{O}_{X_T}} \mathcal{F}_T$ ; voir Théorie des
déformations, lemme \ref{defos-lemma-deform-module-ringed-topoi}. Pour un
$\mathcal{O}_{X_T}$-module $\mathcal{G}$, nous emploierons l'abréviation
$
f_T^*\mathcal{J} \otimes_{\mathcal{O}_{X_T}} \mathcal{G} =
\mathcal{G} \otimes_{\mathcal{O}_T} \mathcal{J}
$.
Comme $\mathcal{Q}$ est plat sur $T$, nous obtenons une suite exacte courte
$$
0 \to
\mathcal{K} \otimes_{\mathcal{O}_T} \mathcal{J} \to
\mathcal{F}_T \otimes_{\mathcal{O}_T} \mathcal{J} \to
\mathcal{Q} \otimes_{\mathcal{O}_T} \mathcal{J} \to
\to 0
$$
En combinant ce qui précède, nous obtenons une extension canonique
$$
0 \to \mathcal{Q} \otimes_{\mathcal{O}_T} \mathcal{J} \to
\pi^{-1}(\mathcal{K})/i(\mathcal{K} \otimes_{\mathcal{O}_T} \mathcal{J}) \to
\mathcal{K} \to 0
$$
de $\mathcal{O}_{X_T}$-modules. Elle définit une classe canonique
$$
o_x(T') \in
\Ext^1_{\mathcal{O}_{X_T}}(\mathcal{K},
\mathcal{Q} \otimes_{\mathcal{O}_T} \mathcal{J})
$$
Si $o_x(T')$ est nul, nous obtenons une scission de la suite exacte courte qui
le définit, autrement dit un $\mathcal{O}_{X_{T'}}$-sous-module
$\mathcal{K}' \subset \pi^{-1}(\mathcal{K})$ s'insérant dans une suite exacte
courte
$0 \to \mathcal{K} \otimes_{\mathcal{O}_T} \mathcal{J} \to
\mathcal{K}' \to \mathcal{K} \to 0$. Il résulte alors du lemme cité ci-dessus
que $\mathcal{Q}' = \mathcal{F}_{T'}/\mathcal{K}'$ est un relèvement de $x$
en un élément de $Q_{\mathcal{F}/X/B}(T')$. Réciproquement, le lecteur voit
que l'existence d'un relèvement implique que $o_x(T')$ est nul. De plus, si
$x \in Q_{\mathcal{F}/X/B}^{fp}(T)$, alors automatiquement
$x' \in Q_{\mathcal{F}/X/B}^{fp}(T')$, d'après Théorie des déformations,
lemme \ref{defos-lemma-deform-fp-module-ringed-topoi}. Si nous avons un jour
besoin de cette remarque, nous en ferons un lemme, formulerons précisément le
résultat et en donnerons une démonstration détaillée (en fait, tout ce qui
précède vaut dans le cadre de topos annelés arbitraires).
\end{remark}

\begin{remark}[Déformations des quotients]
\label{remark-q-defos}
Dans la situation \ref{situation-q}, {\bf supposons} que $\mathcal{F}$ soit
plat sur $B$. Nous poursuivons la discussion de la remarque
\ref{remark-q-obs}. Supposons $o_x(T') = 0$. Nous affirmons alors que
l'ensemble des relèvements $x' \in Q_{\mathcal{F}/X/B}(T')$ est un espace
principal homogène sous le groupe
$$
\Hom_{\mathcal{O}_{X_T}}(\mathcal{K},
\mathcal{Q} \otimes_{\mathcal{O}_T} \mathcal{J})
$$
En effet, étant donné un quotient
$\mathcal{F}_{T'} \to \mathcal{Q}'$ plat sur $T'$ qui relève le quotient
$\mathcal{Q}$, nous obtenons un diagramme commutatif à lignes et colonnes
exactes
$$
\xymatrix{
& 0 \ar[d] & 0 \ar[d] & 0 \ar[d] \\
0 \ar[r] &
\mathcal{K} \otimes \mathcal{J} \ar[r] \ar[d] &
\mathcal{F}_T \otimes \mathcal{J} \ar[r] \ar[d] &
\mathcal{Q} \otimes \mathcal{J} \ar[r] \ar[d] &
0 \\
0 \ar[r] &
\mathcal{K}' \ar[r] \ar[d] &
\mathcal{F}_{T'} \ar[r] \ar[d] &
\mathcal{Q}' \ar[r] \ar[d] &
0 \\
0 \ar[r] &
\mathcal{K} \ar[d] \ar[r] &
\mathcal{F}_T \ar[d] \ar[r] &
\mathcal{Q} \ar[d] \ar[r] &
0 \\
& 0 & 0 & 0
}
$$
(pour le voir, utiliser les observations faites dans la remarque précédente).
Étant donnée une application
$\varphi : \mathcal{K} \to \mathcal{Q} \otimes \mathcal{J}$, nous pouvons
considérer le sous-faisceau
$\mathcal{K}'_\varphi \subset \mathcal{F}_{T'}$ formé des sections locales
$s$ dont l'image dans $\mathcal{F}_T$ est une section locale $k$ de
$\mathcal{K}$ et dont l'image dans $\mathcal{Q}'$ est la section locale
$\varphi(k)$ de $\mathcal{Q} \otimes \mathcal{J}$. Posons alors
$\mathcal{Q}'_\varphi = \mathcal{F}_{T'}/\mathcal{K}'_\varphi$.
Réciproquement, tout second relèvement de $x$ correspond à l'un des quotients
ainsi construits. Si nous avons un jour besoin de cette remarque, nous en
ferons un lemme, formulerons précisément le résultat et en donnerons une
démonstration détaillée (en fait, tout ce qui précède vaut dans le cadre de
topos annelés arbitraires).
\end{remark}










\section{Le foncteur de Quot}
\label{section-quot}

\noindent
Dans cette section, nous démontrons que le foncteur de Quot est un espace
algébrique.

\begin{situation}
\label{situation-quot}
Soit $S$ un schéma. Soit $f : X \to B$ un morphisme d'espaces algébriques sur
$S$. Supposons que $f$ soit de présentation finie. Soit $\mathcal{F}$ un
$\mathcal{O}_X$-module quasi-cohérent. Pour tout schéma $T$ sur $B$, nous
désignerons par $X_T$ le changement de base de $X$ à $T$ et par
$\mathcal{F}_T$ l'image inverse de $\mathcal{F}$ par le morphisme de
projection $X_T = X \times_S T \to X$. Pour un tel $T$, posons
$$
\Quotfunctor_{\mathcal{F}/X/B}(T) =
\left\{
\begin{matrix}
\text{quotients }\mathcal{F}_T \to \mathcal{Q}\text{ où }
\mathcal{Q}\text{ est un}\\
\mathcal{O}_{X_T}\text{-module quasi-cohérent de présentation finie, plat sur }T\\
\text{à support propre sur }T
\end{matrix}
\right\}
$$
D'après Catégories dérivées des espaces, lemme
\ref{spaces-perfect-lemma-base-change-module-support-proper-over-base}, c'est
un sous-foncteur du foncteur $Q^{fp}_{\mathcal{F}/X/B}$ étudié dans la section
\ref{section-functor-quotients}. Nous obtenons donc un foncteur
\begin{equation}
\label{equation-quot}
\Quotfunctor_{\mathcal{F}/X/B} : (\Sch/B)^{opp} \longrightarrow \textit{Sets}
\end{equation}
C'est le {\it foncteur de Quot} associé à $\mathcal{F}/X/B$.
\end{situation}

\noindent
Dans la situation \ref{situation-quot}, nous regardons parfois
$\Quotfunctor_{\mathcal{F}/X/B}$ comme un foncteur
$(\Sch/S)^{opp} \to \textit{Sets}$ muni d'un morphisme
$\Quotfunctor_{\mathcal{F}/X/B} \to B$. En effet, si $T$ est un schéma sur
$S$, un élément de $\Quotfunctor_{\mathcal{F}/X/B}(T)$ est un couple
$(h, \mathcal{Q})$, où $h$ est un morphisme $h : T \to B$ et où $Q$ est un
quotient $\mathcal{F}_T \to \mathcal{Q}$ de présentation finie et plat sur
$T$, sur $X_T = X \times_{B, h} T$, à support propre sur $T$. En particulier,
lorsque nous disons que $\Quotfunctor_{\mathcal{F}/X/B}$ est un espace
algébrique, nous entendons que le foncteur correspondant
$(\Sch/S)^{opp} \to \textit{Sets}$ est un espace algébrique.

\begin{lemma}
\label{lemma-quot-sheaf}
Dans la situation \ref{situation-quot}. Le foncteur
$\Quotfunctor_{\mathcal{F}/X/B}$ satisfait à la propriété de faisceau pour la
topologie fpqc.
\end{lemma}

\begin{proof}
Dans le lemme \ref{lemma-q-sheaf}, nous avons vu que le foncteur
$\text{Q}^{fp}_{\mathcal{F}/X/S}$ est un faisceau. Rappelons que, pour un
schéma $T$ sur $S$, le sous-ensemble
$\Quotfunctor_{\mathcal{F}/X/S}(T) \subset \text{Q}_{\mathcal{F}/X/S}(T)$
sélectionne les quotients dont le support est propre sur $T$. Cela définit un
sous-faisceau d'après Descente sur les espaces, lemme
\ref{spaces-descent-lemma-descending-property-proper}, combiné à Morphismes
d'espaces, lemme \ref{spaces-morphisms-lemma-flat-pullback-support}, qui montre
que la formation du support schématique commute au changement de base plat.
\end{proof}

\noindent
Vérification de cohérence : $\Quotfunctor_{\mathcal{F}/X/B}$ joue le même rôle
parmi les espaces algébriques sur $S$.

\begin{lemma}
\label{lemma-extend-quot-to-spaces}
Dans la situation \ref{situation-quot}. Soit $T$ un espace algébrique sur $S$.
Nous avons
$$
\Mor_{\Sh((\Sch/S)_{fppf})}(T,  \Quotfunctor_{\mathcal{F}/X/B}) =
\left\{
\begin{matrix}
(h, \mathcal{F}_T \to \mathcal{Q}) \text{ où }
h : T \to B \text{ et où}\\
\mathcal{Q}\text{ est de présentation finie et}\\
\text{plate sur }T\text{ à support propre sur }T
\end{matrix}
\right\}
$$
où $\mathcal{F}_T$ désigne l'image inverse de $\mathcal{F}$ sur l'espace
algébrique $X \times_{B, h} T$.
\end{lemma}

\begin{proof}
Observons que les membres de gauche et de droite de l'égalité sont des
sous-ensembles des membres de gauche et de droite de la seconde égalité du
lemme \ref{lemma-extend-q-to-spaces}. Pour voir que ces sous-ensembles se
correspondent par l'identification donnée dans la démonstration de ce lemme, il
suffit de montrer ceci : étant donnés $h : T \to B$, un morphisme étale
surjectif $U \to T$ et un $\mathcal{O}_{X_T}$-module quasi-cohérent de type
fini $\mathcal{Q}$, les conditions suivantes sont équivalentes :
\begin{enumerate}
\item le support schématique de $\mathcal{Q}$ est propre sur $T$, et
\item le support schématique de $(X_U \to X_T)^*\mathcal{Q}$ est propre sur
$U$.
\end{enumerate}
Cela résulte de Descente sur les espaces, lemme
\ref{spaces-descent-lemma-descending-property-proper}, combiné à Morphismes
d'espaces, lemme \ref{spaces-morphisms-lemma-flat-pullback-support}, qui montre
que la formation du support schématique commute au changement de base plat.
\end{proof}

\begin{proposition}
\label{proposition-quot}
Soit $S$ un schéma. Soit $f : X \to B$ un morphisme d'espaces algébriques sur
$S$. Soit $\mathcal{F}$ un faisceau quasi-cohérent sur $X$. Si $f$ est de
présentation finie et séparé, alors
$\Quotfunctor_{\mathcal{F}/X/B}$ est un espace algébrique. Si $\mathcal{F}$
est de présentation finie, alors
$\Quotfunctor_{\mathcal{F}/X/B} \to B$ est localement de présentation finie.
\end{proposition}

\begin{proof}
D'après le lemme \ref{lemma-quot-sheaf},
$\Quotfunctor_{\mathcal{F}/X/B}$ est un faisceau pour la topologie fppf. Soit
$\textit{Quot}_{\mathcal{F}/X/B}$ le champ en groupoïdes correspondant à
$\Quotfunctor_{\mathcal{F}/X/S}$ ; voir Champs algébriques, section
\ref{algebraic-section-split}. D'après Champs algébriques, proposition
\ref{algebraic-proposition-algebraic-stack-no-automorphisms}, il suffit de
montrer que $\textit{Quot}_{\mathcal{F}/X/B}$ est un champ algébrique.
Considérons le $1$-morphisme de champs en groupoïdes
$$
\textit{Quot}_{\mathcal{F}/X/S}
\longrightarrow
\Cohstack_{X/B}
$$
sur $(\Sch/S)_{fppf}$, qui associe au quotient
$\mathcal{F}_T \to \mathcal{Q}$ le module $\mathcal{Q}$. D'après le théorème
\ref{theorem-coherent-algebraic-general}, nous savons que
$\Cohstack_{X/B}$ est un champ algébrique. D'après Champs algébriques, lemme
\ref{algebraic-lemma-representable-morphism-to-algebraic}, il suffit de montrer
que ce $1$-morphisme est représentable par des espaces algébriques.

\medskip\noindent
Soit $T$ un schéma sur $S$ et soit l'objet $(h, \mathcal{G})$ de
$\Cohstack_{X/B}$ sur $T$ qui correspond à un $1$-morphisme
$\xi : (\Sch/T)_{fppf} \to \Cohstack_{X/B}$. Le $2$-produit fibré
$$
\mathcal{Z} =
(\Sch/T)_{fppf}
\times_{\xi, \Cohstack_{X/B}}
\textit{Quot}_{\mathcal{F}/X/S}
$$
est un champ en sétoïdes ; voir Champs, lemme
\ref{stacks-lemma-2-fibre-product-gives-stack-in-setoids}. Le faisceau
d'ensembles correspondant (c'est-à-dire le foncteur ; voir Champs, lemmes
\ref{stacks-lemma-2-fibre-product-gives-stack-in-setoids} et
\ref{stacks-lemma-when-stack-in-sets}) associe à un schéma $T'/T$ l'ensemble
des surjections $u : \mathcal{F}_{T'} \to \mathcal{G}_{T'}$ de modules
quasi-cohérents sur $X_{T'}$. Nous voyons donc que $\mathcal{Z}$ est
représentable par un sous-espace ouvert (d'après Platitude sur les espaces,
lemme \ref{spaces-flat-lemma-F-surj-open}) de l'espace algébrique
$\mathit{Hom}(\mathcal{F}_T, \mathcal{G})$ de la proposition
\ref{proposition-hom}.
\end{proof}

\begin{remark}[Quot au moyen des axiomes d'Artin]
\label{remark-quot-via-artins-axioms}
Soit $S$ un schéma noethérien dont tous les anneaux locaux sont des anneaux G.
Soit $X$ un espace algébrique sur $S$ dont le morphisme structural
$f : X \to S$ est de présentation finie et séparé. Soit $\mathcal{F}$ un
faisceau quasi-cohérent de présentation finie sur $X$, plat sur $S$. Nous
esquissons dans cette remarque comment employer les axiomes d'Artin pour
démontrer que $\Quotfunctor_{\mathcal{F}/X/S}$ est un espace algébrique
localement de présentation finie sur $S$, sans utiliser l'algébricité du champ
des faisceaux cohérents comme dans la démonstration de la proposition
\ref{proposition-quot}.

\medskip\noindent
Vérifions les conditions énumérées dans Axiomes d'Artin, proposition
\ref{artin-proposition-spaces-diagonal-representable}. On peut voir comme suit
la représentabilité de la diagonale de
$\Quotfunctor_{\mathcal{F}/X/S}$ : supposons donnés deux quotients
$\mathcal{F}_T \to \mathcal{Q}_i$, $i = 1, 2$. Désignons par
$\mathcal{K}_1$ le noyau du premier. Il faut alors montrer que le lieu de $T$
au-dessus duquel $u : \mathcal{K}_1 \to \mathcal{Q}_2$ devient nul est
représentable. Cela résulte par exemple de Platitude sur les espaces, lemme
\ref{spaces-flat-lemma-F-zero-closed-proper}, ou d'une discussion antérieure
du faisceau $\mathit{Hom}$ dans ce chapitre. Les axiomes [0] (faisceau), [1]
(limites), [2] (Rim--Schlessinger) résultent des lemmes
\ref{lemma-quot-sheaf}, \ref{lemma-q-limit-preserving} et
\ref{lemma-q-RS-star} (moyennant un travail supplémentaire pour traiter la
condition de propreté). L'axiome [3] (dimension finie des espaces tangents)
résulte de la description des déformations infinitésimales dans la remarque
\ref{remark-q-defos} et de la finitude de la cohomologie des faisceaux
cohérents sur les espaces algébriques propres sur des corps (Cohomologie des
espaces, lemme \ref{spaces-cohomology-lemma-proper-pushforward-coherent}).
L'axiome [4] (effectivité des objets formels) résulte du théorème d'existence
de Grothendieck (Plus sur les morphismes d'espaces, théorème
\ref{spaces-more-morphisms-theorem-grothendieck-existence}). Comme toujours,
l'axiome [5] (ouverture de la versalité) est le plus délicat à vérifier. On
peut par exemple utiliser la théorie de l'obstruction décrite dans la remarque
\ref{remark-q-obs} et la description des déformations de la remarque
\ref{remark-q-defos}, en appliquant le critère d'Axiomes d'Artin, lemme
\ref{artin-lemma-get-openness-obstruction-theory}. Comparer avec la seconde
démonstration du lemme \ref{lemma-coherent-defo-thy}.
\end{remark}









\section{Le foncteur de Hilbert}
\label{section-hilb}

\noindent
Dans cette section, nous démontrons que le foncteur Hilb est un espace
algébrique.

\begin{situation}
\label{situation-hilb}
Soit $S$ un schéma. Soit $f : X \to B$ un morphisme d'espaces algébriques sur
$S$. Supposons que $f$ soit de présentation finie. Pour tout schéma $T$ sur
$B$, nous désignerons par $X_T$ le changement de base de $X$ à $T$. Pour un
tel $T$, posons
$$
\Hilbfunctor_{X/B}(T) =
\left\{
\begin{matrix}
\text{sous-espaces fermés }Z \subset X_T\text{ tels que }Z \to T\\
\text{soit de présentation finie, plat et propre}
\end{matrix}
\right\}
$$
Comme le changement de base conserve les propriétés requises (Espaces, lemme
\ref{spaces-lemma-base-change-immersions}, et Morphismes d'espaces, lemmes
\ref{spaces-morphisms-lemma-base-change-finite-presentation},
\ref{spaces-morphisms-lemma-base-change-flat} et
\ref{spaces-morphisms-lemma-base-change-proper}), nous obtenons un foncteur
\begin{equation}
\label{equation-hilb}
\Hilbfunctor_{X/B} : (\Sch/B)^{opp} \longrightarrow \textit{Sets}
\end{equation}
C'est le {\it foncteur de Hilbert} associé à $X/B$.
\end{situation}

\noindent
Dans la situation \ref{situation-hilb}, nous regardons parfois
$\Hilbfunctor_{X/B}$ comme un foncteur
$(\Sch/S)^{opp} \to \textit{Sets}$ muni d'un morphisme
$\Hilbfunctor_{X/S} \to B$. En effet, si $T$ est un schéma sur $S$, un
élément de $\Hilbfunctor_{X/B}(T)$ est un couple $(h, Z)$ où $h$ est un
morphisme $h : T \to B$ et où $Z \subset X_T = X \times_{B, h} T$ est un
sous-schéma fermé, plat, propre et de présentation finie sur $T$. En
particulier, lorsque nous disons que $\Hilbfunctor_{X/B}$ est un espace
algébrique, nous entendons que le foncteur correspondant
$(\Sch/S)^{opp} \to \textit{Sets}$ est un espace algébrique.

\medskip\noindent
Bien entendu, le foncteur de Hilbert n'est qu'un cas particulier du foncteur
de Quot.

\begin{lemma}
\label{lemma-hilb-is-quot}
Dans la situation \ref{situation-hilb}, nous avons
$\Hilbfunctor_{X/B} = \Quotfunctor_{\mathcal{O}_X/X/B}$.
\end{lemma}

\begin{proof}
Soit $T$ un schéma sur $B$. Étant donné un élément
$Z \in \Hilbfunctor_{X/B}(T)$, nous pouvons considérer le quotient
$\mathcal{O}_{X_T} \to i_*\mathcal{O}_Z$, où $i : Z \to X_T$ est le
morphisme d'inclusion. Remarquons que $i_*\mathcal{O}_Z$ est quasi-cohérent.
Comme $Z \to T$ et $X_T \to T$ sont de présentation finie, nous voyons que
$i$ est de présentation finie (Morphismes d'espaces, lemme
\ref{spaces-morphisms-lemma-finite-presentation-permanence}) ; ainsi,
$i_*\mathcal{O}_Z$ est un $\mathcal{O}_{X_T}$-module de présentation finie
(Descente sur les espaces, lemme
\ref{spaces-descent-lemma-finite-finitely-presented-module}). Comme
$Z \to T$ est propre, nous voyons que $i_*\mathcal{O}_Z$ est à support propre
sur $T$ (au sens défini dans Catégories dérivées des espaces, section
\ref{spaces-perfect-section-proper-over-base}). Comme $\mathcal{O}_Z$ est plat
sur $T$ et que $i$ est affine, nous voyons que $i_*\mathcal{O}_Z$ est plat sur
$T$ (nous omettons un petit argument). Par conséquent,
$\mathcal{O}_{X_T} \to i_*\mathcal{O}_Z$ est un élément de
$\Quotfunctor_{\mathcal{O}_X/X/B}(T)$.

\medskip\noindent
Réciproquement, étant donné un élément
$\mathcal{O}_{X_T} \to \mathcal{Q}$ de
$\Quotfunctor_{\mathcal{O}_X/X/B}(T)$, nous pouvons considérer l'immersion
fermée $i : Z \to X_T$ correspondant au faisceau d'idéaux quasi-cohérent
$\mathcal{I} = \Ker(\mathcal{O}_{X_T} \to \mathcal{Q})$ (Morphismes
d'espaces, lemme \ref{spaces-morphisms-lemma-closed-immersion-ideals}). Par
construction de $Z$, nous avons $\mathcal{Q} = i_*\mathcal{O}_Z$. Nous pouvons
alors lire les arguments précédents à rebours pour voir que $Z$ définit un
élément de $\Hilbfunctor_{X/B}(T)$. Par exemple, $\mathcal{I}$ est
quasi-cohérent de type fini (Modules sur les sites, lemme
\ref{sites-modules-lemma-kernel-surjection-finite-onto-finite-presentation}) ;
donc $i : Z \to X_T$ est de présentation finie (Morphismes d'espaces, lemme
\ref{spaces-morphisms-lemma-closed-immersion-finite-presentation}) ; donc
$Z \to T$ est de présentation finie (Morphismes d'espaces, lemme
\ref{spaces-morphisms-lemma-composition-finite-presentation}). La propreté de
$Z \to T$ résulte de la discussion dans Catégories dérivées des espaces,
section \ref{spaces-perfect-section-proper-over-base}. La platitude de
$Z \to T$ résulte de la platitude de $\mathcal{Q}$ sur $T$.

\medskip\noindent
Nous omettons la vérification (immédiate) que les deux constructions ci-dessus
sont inverses l'une de l'autre.
\end{proof}

\noindent
Vérification de cohérence : le faisceau $\Hilbfunctor_{X/B}$ joue le même rôle
parmi les espaces algébriques sur $S$.

\begin{lemma}
\label{lemma-extend-hilb-to-spaces}
Dans la situation \ref{situation-hilb}. Soit $T$ un espace algébrique sur $S$.
Nous avons
$$
\Mor_{\Sh((\Sch/S)_{fppf})}(T, \Hilbfunctor_{X/B}) =
\left\{
\begin{matrix}
(h, Z)\text{ où }h : T \to B,\ Z \subset X_T \\
\text{de présentation finie, plat, propre sur }T
\end{matrix}
\right\}
$$
où $X_T = X \times_{B, h} T$.
\end{lemma}

\begin{proof}
D'après le lemme \ref{lemma-hilb-is-quot}, nous avons
$\Hilbfunctor_{X/B} = \Quotfunctor_{\mathcal{O}_X/X/B}$. Nous pouvons donc
appliquer le lemme \ref{lemma-extend-quot-to-spaces} pour voir que le membre de
gauche est en bijection avec l'ensemble des surjections
$\mathcal{O}_{X_T} \to \mathcal{Q}$ qui sont de présentation finie, plates
sur $T$ et à support propre sur $T$. En raisonnant exactement comme dans la
démonstration du lemme \ref{lemma-hilb-is-quot}, nous voyons que ces quotients
correspondent exactement aux immersions fermées $Z \to X_T$ telles que
$Z \to T$ soit propre, plat et de présentation finie.
\end{proof}

\begin{proposition}
\label{proposition-hilb}
Soit $S$ un schéma. Soit $f : X \to B$ un morphisme d'espaces algébriques sur
$S$. Si $f$ est de présentation finie et séparé, alors
$\Hilbfunctor_{X/B}$ est un espace algébrique localement de présentation finie
sur $B$.
\end{proposition}

\begin{proof}
Conséquence immédiate du lemme \ref{lemma-hilb-is-quot} et de la proposition
\ref{proposition-quot}.
\end{proof}






\section{Le champ de Picard}
\label{section-picard-stack}

\noindent
Le champ de Picard d'un morphisme d'espaces algébriques a été introduit dans
Exemples de champs, section \ref{examples-stacks-section-picard-stack}. Nous
déduirons du lemme suivant que c'est un sous-champ ouvert du champ des faisceaux
cohérents (dans les cas favorables).

\begin{lemma}
\label{lemma-picard-stack-open-in-coh}
Soit $S$ un schéma. Soit $f : X \to B$ un morphisme d'espaces algébriques sur
$S$ qui est plat, de présentation finie et propre. L'application naturelle
$$
\Picardstack_{X/B} \longrightarrow \Cohstack_{X/B}
$$
est représentable par des immersions ouvertes.
\end{lemma}

\begin{proof}
Observons que l'application envoie simplement un triplet
$(T, g, \mathcal{L})$ comme dans Exemples de champs, section
\ref{examples-stacks-section-picard-stack}, sur le même triplet
$(T, g, \mathcal{L})$, mais considéré maintenant comme un triplet du genre
décrit dans la situation \ref{situation-coherent}. Cela fonctionne parce que
le $\mathcal{O}_{X_T}$-module inversible $\mathcal{L}$ est certainement un
$\mathcal{O}_{X_T}$-module de présentation finie, qu'il est plat sur $T$ parce
que $X_T \to T$ est plat, et que son support est propre sur $T$ puisque
$X_T \to T$ est propre (Morphismes d'espaces, lemmes
\ref{spaces-morphisms-lemma-base-change-flat} et
\ref{spaces-morphisms-lemma-base-change-proper}). L'énoncé a donc un sens.

\medskip\noindent
Cela dit, il est clair que le contenu du lemme est le suivant : étant donné un
objet $(T, g, \mathcal{F})$ de $\Cohstack_{X/B}$, il existe un sous-schéma
ouvert $U \subset T$ tel que, pour un morphisme de schémas $T' \to T$, les
conditions suivantes soient équivalentes :
\begin{enumerate}
\item[(a)] $T' \to T$ se factorise par $U$,
\item[(b)] l'image inverse $\mathcal{F}_{T'}$ de $\mathcal{F}$ par
$X_{T'} \to X_T$ est inversible.
\end{enumerate}
Soit $W \subset |X_T|$ l'ensemble des points $x \in |X_T|$ tels que
$\mathcal{F}$ soit localement libre au voisinage de $x$. D'après Plus sur les
morphismes d'espaces, lemme
\ref{spaces-more-morphisms-lemma-finite-free-open}.
$W$ est ouvert et la formation de $W$ commute à tout changement de base. Il
est clair que, si $T' \to T$ satisfait à (b), alors
$|X_{T'}| \to |X_T|$ est à valeurs dans $W$. Nous pouvons donc, pour
construire $U$, remplacer $T$ par l'ouvert
$T \setminus f_T(|X_T| \setminus W)$. Après cette opération, nous sommes dans
la situation où $\mathcal{F}$ est localement libre de type fini. Dans ce cas,
nous obtenons une décomposition en réunion disjointe
$X_T = X_0 \amalg X_1 \amalg X_2 \amalg \ldots$ de sous-espaces ouverts et
fermés telle que la restriction de $\mathcal{F}$ à $X_i$ soit localement libre
de rang $i$. Alors il est clair que
$$
U = T \setminus f_T(|X_0| \cup |X_2| \cup |X_3| \cup \ldots )
$$
convient. (Remarquons que si nous supposons $T$ quasi-compact, alors $X_T$ est
quasi-compact ; seul un nombre fini de $X_i$ sont donc non vides et $U$ est
bien ouvert.)
\end{proof}

\begin{proposition}
\label{proposition-pic}
Soit $S$ un schéma. Soit $f : X \to B$ un morphisme d'espaces algébriques sur
$S$. Si $f$ est plat, de présentation finie et propre, alors
$\Picardstack_{X/B}$ est un champ algébrique.
\end{proposition}

\begin{proof}
Conséquence immédiate du lemme \ref{lemma-picard-stack-open-in-coh}, du lemme
de Champs algébriques
\ref{algebraic-lemma-representable-morphism-to-algebraic}, et soit du théorème
\ref{theorem-coherent-algebraic}, soit du théorème
\ref{theorem-coherent-algebraic-general}
\end{proof}

\section{Le foncteur de Picard}
\label{section-picard-functor}

\noindent
Dans cette section, nous revenons sur le foncteur de Picard étudié dans
Schémas de Picard des courbes, section \ref{pic-section-picard-functor}.
La discussion sera plus générale, car nous voulons étudier le foncteur de
Picard d'un morphisme d'espaces algébriques comme dans la section consacrée au
champ de Picard, voir la section \ref{section-picard-stack}.

\medskip\noindent
Soit $S$ un schéma et soit $X$ un espace algébrique sur $S$.
Un faisceau inversible sur $X$ est un $\mathcal{O}_X$-module inversible
sur $X_\etale$, voir
Modules sur les sites, définition \ref{sites-modules-definition-invertible-sheaf}.
Le groupe des classes d'isomorphisme de modules inversibles est noté
$\Pic(X)$, voir
Modules sur les sites, définition \ref{sites-modules-definition-pic}.
Étant donné un morphisme $f : X \to Y$ d'espaces algébriques sur $S$,
l'image inverse définit un homomorphisme de groupes $\Pic(Y) \to \Pic(X)$.
La correspondance
$X \leadsto \Pic(X)$ est un foncteur contravariant de la catégorie
des schémas vers la catégorie des groupes abéliens. Ce foncteur n'est pas
représentable, mais il se trouve qu'une variante relative de cette construction
l'est parfois.

\begin{situation}
\label{situation-pic}
Soit $S$ un schéma.
Soit $f : X \to B$ un morphisme d'espaces algébriques sur $S$.
Nous définissons
$$
\Picardfunctor_{X/B} : (\Sch/B)^{opp} \longrightarrow \textit{Sets}
$$
comme le faisceau fppf associé au foncteur qui, à un schéma $T$
sur $B$, associe le groupe $\Pic(X_T)$.
\end{situation}

\noindent
Dans la situation \ref{situation-pic}, nous considérons parfois
$\Picardfunctor_{X/B}$ comme un foncteur $(\Sch/S)^{opp} \to \textit{Sets}$
muni d'un morphisme $\Picardfunctor_{X/B} \to B$. De ce point de vue,
nous définissons $\Picardfunctor_{X/B}$ comme le faisceau fppf associé au
foncteur
$$
T/S \longmapsto \{(h, \mathcal{L}) \mid
h : T \to B,\ \mathcal{L} \in \Pic(X \times_{B, h} T)\}
$$
En particulier, lorsque nous disons que $\Picardfunctor_{X/B}$ est un espace
algébrique, nous entendons que le foncteur correspondant
$(\Sch/S)^{opp} \to \textit{Sets}$ est un espace algébrique.

\medskip\noindent
Une remarque souvent utilisée est que, si $T$ est un schéma sur $B$, alors
$\Picardfunctor_{X_T/T}$ est la restriction de $\Picardfunctor_{X/B}$ à
$(\Sch/T)_{fppf}$.

\begin{lemma}
\label{lemma-pic-over-pic}
Dans la situation \ref{situation-pic},
le foncteur $\Picardfunctor_{X/B}$ est le faisceau associé au
foncteur $T \mapsto \Ob(\Picardstack_{X/B, T})/\cong$.
\end{lemma}

\begin{proof}
Puisque la catégorie fibre $\Picardstack_{X/B, T}$ du champ de Picard
$\Picardstack_{X/B}$ au-dessus de $T$ est la catégorie des faisceaux inversibles
sur $X_T$ (voir la section \ref{section-picard-stack} et
Exemples de champs, section \ref{examples-stacks-section-picard-stack}),
cela découle immédiatement des définitions.
\end{proof}

\noindent
Il n'est en fait pas trivial de déterminer la valeur de $\Picardfunctor_{X/B}$
sur les schémas $T$ au-dessus de $B$. Voici un lemme qui aide à accomplir cette
tâche.

\begin{lemma}
\label{lemma-flat-geometrically-connected-fibres}
Dans la situation \ref{situation-pic}.
Si $\mathcal{O}_T \to f_{T, *}\mathcal{O}_{X_T}$ est un isomorphisme
pour tout schéma $T$ sur $B$, alors
$$
0 \to \Pic(T) \to \Pic(X_T) \to \Picardfunctor_{X/B}(T)
$$
est une suite exacte pour tout $T$.
\end{lemma}

\begin{proof}
Nous pouvons remplacer $B$ par $T$ et $X$ par $X_T$, et supposer que $B = T$
afin d'alléger les notations. Soit $\mathcal{N}$ un
$\mathcal{O}_B$-module inversible. Si $f^*\mathcal{N} \cong \mathcal{O}_X$, alors
l'hypothèse donne $f_*f^*\mathcal{N} \cong f_*\mathcal{O}_X \cong \mathcal{O}_B$.
Comme $\mathcal{N}$ est localement trivial, l'application canonique
$\mathcal{N} \to f_*f^*\mathcal{N}$ est localement un isomorphisme (car
$\mathcal{O}_B \to f_*f^*\mathcal{O}_B$ est un isomorphisme par hypothèse).
Nous en concluons que
$\mathcal{N} \to f_*f^*\mathcal{N} \to \mathcal{O}_B$ est un isomorphisme,
et voyons que $\mathcal{N}$ est trivial. Cela prouve que la première flèche
est injective.

\medskip\noindent
Soit $\mathcal{L}$ un $\mathcal{O}_X$-module inversible appartenant au
noyau de $\Pic(X) \to \Picardfunctor_{X/B}(B)$. Il existe alors un
recouvrement fppf $\{B_i \to B\}$ tel que l'image inverse de $\mathcal{L}$
soit le faisceau inversible trivial sur $X_{B_i}$. Choisissons une section
trivialisante $s_i$. Alors $\text{pr}_0^*s_i$ et $\text{pr}_1^*s_j$ sont toutes
deux des sections trivialisantes de $\mathcal{L}$ sur
$X_{B_i \times_B B_j}$ et diffèrent donc par une unité multiplicative
$$
f_{ij} \in
\Gamma(X_{S_i \times_B B_j}, \mathcal{O}_{X_{B_i \times_B B_j}}^*) =
\Gamma(B_i \times_B B_j, \mathcal{O}_{B_i \times_N B_j}^*)
$$
(l'égalité résultant de notre hypothèse sur l'image directe des faisceaux
structuraux). Bien entendu, ces éléments satisfont à la condition de cocycle
sur $B_i \times_B B_j \times_B B_k$ ; ils définissent donc une donnée de
descente sur des faisceaux inversibles relativement au recouvrement fppf
$\{B_i \to B\}$. D'après Descente, proposition
\ref{descent-proposition-fpqc-descent-quasi-coherent}, il existe un
$\mathcal{O}_B$-module inversible $\mathcal{N}$ muni de trivialisations sur les
$B_i$ dont la donnée de descente associée est $\{f_{ij}\}$. (La proposition
s'applique parce que $B$ est un schéma, par le remplacement effectué au début
de la démonstration.) Alors $f^*\mathcal{N} \cong \mathcal{L}$, car le
foncteur des données de descente vers les modules est pleinement fidèle.
\end{proof}

\begin{lemma}
\label{lemma-flat-geometrically-connected-fibres-with-section}
Dans la situation \ref{situation-pic}, soit $\sigma : B \to X$ une section.
Supposons que $\mathcal{O}_T \to f_{T, *}\mathcal{O}_{X_T}$ soit un isomorphisme
pour tout $T$ sur $B$. Alors
$$
0 \to \Pic(T) \to \Pic(X_T) \to \Picardfunctor_{X/B}(T) \to 0
$$
est une suite exacte scindée, une rétraction étant donnée par
$\sigma_T^* : \Pic(X_T) \to \Pic(T)$.
\end{lemma}

\begin{proof}
Notons $K(T) = \Ker(\sigma_T^* : \Pic(X_T) \to \Pic(T))$.
Comme $\sigma$ est une section de $f$, nous voyons que $\Pic(X_T)$ est la somme
directe de $\Pic(T)$ et de $K(T)$.
Ainsi, d'après le lemme \ref{lemma-flat-geometrically-connected-fibres},
nous avons $K(T) \subset \Picardfunctor_{X/B}(T)$ pour tout $T$. De plus,
il ressort clairement de la construction que $\Picardfunctor_{X/B}$ est le
faisceau associé au préfaisceau $K$. Pour achever la démonstration, il suffit
de montrer que $K$ satisfait à la condition de faisceau pour les recouvrements
fppf, ce que nous faisons dans le paragraphe suivant.

\medskip\noindent
Soit $\{T_i \to T\}$ un recouvrement fppf. Soient $\mathcal{L}_i$ des éléments
de $K(T_i)$ qui ont les mêmes images dans $K(T_i \times_T T_j)$ pour tous
$i$ et $j$. Choisissons, pour tout $i$, un isomorphisme
$\alpha_i : \mathcal{O}_{T_i} \to \sigma_{T_i}^*\mathcal{L}_i$.
Choisissons un isomorphisme
$$
\varphi_{ij} :
\mathcal{L}_i|_{X_{T_i \times_T T_j}}
\longrightarrow
\mathcal{L}_j|_{X_{T_i \times_T T_j}}
$$
Si l'application
$$
\alpha_j|_{T_i \times_T T_j} \circ
\sigma_{T_i \times_T T_j}^*\varphi_{ij} \circ
\alpha_i|_{T_i \times_T T_j} :
\mathcal{O}_{T_i \times_T T_j} \to \mathcal{O}_{T_i \times_T T_j}
$$
n'est pas la multiplication par $1$, mais par un certain $u_{ij}$, nous pouvons
multiplier $\varphi_{ij}$ par $u_{ij}^{-1}$ pour y remédier. Cela fait,
considérons l'endomorphisme
$$
\varphi_{ki}|_{X_{T_i \times_T T_j \times_T T_k}} \circ
\varphi_{jk}|_{X_{T_i \times_T T_j \times_T T_k}} \circ
\varphi_{ij}|_{X_{T_i \times_T T_j \times_T T_k}}
\quad\text{sur}\quad
\mathcal{L}_i|_{X_{T_i \times_T T_j \times_T T_k}}
$$
qui est donné par la multiplication par une certaine section $f_{ijk}$ du
faisceau structural de $X_{T_i \times_T T_j \times_T T_k}$. Par notre choix
de $\varphi_{ij}$, nous voyons que l'image inverse de cette application par
$\sigma$ est la multiplication par $1$. Notre hypothèse sur les fonctions sur
$X$ donne alors $f_{ijk} = 1$. Nous obtenons ainsi une donnée de descente pour
le recouvrement fppf
$\{X_{T_i} \to X\}$. D'après Descente sur les espaces, proposition
\ref{spaces-descent-proposition-fpqc-descent-quasi-coherent}, il existe un
$\mathcal{O}_{X_T}$-module inversible $\mathcal{L}$ et un isomorphisme
$\alpha : \mathcal{O}_T \to \sigma_T^*\mathcal{L}$ dont les images inverses
sur $X_{T_i}$ redonnent $(\mathcal{L}_i, \alpha_i)$ (nous omettons un petit
détail). Ainsi, $\mathcal{L}$ définit un objet de $K(T)$, comme voulu.
\end{proof}

\noindent
Dans la situation \ref{situation-pic}, soit $\sigma : B \to X$ une section.
Nous notons $\Picardstack_{X/B, \sigma}$ la catégorie définie comme suit :
\begin{enumerate}
\item Un objet est un quadruplet $(T, h, \mathcal{L}, \alpha)$, où
$(T, h, \mathcal{L})$ est un objet de $\Picardstack_{X/B}$ au-dessus de $T$ et
$\alpha : \mathcal{O}_T \to \sigma_T^*\mathcal{L}$ est un isomorphisme.
\item Un morphisme $(g, \varphi) : (T, h, \mathcal{L}, \alpha)
\to (T', h', \mathcal{L}', \alpha')$ est donné par un morphisme de schémas
$g : T \to T'$ tel que $h = h' \circ g$ et un isomorphisme
$\varphi : (g')^*\mathcal{L}' \to \mathcal{L}$ tel que
$\sigma_T^*\varphi \circ g^*\alpha' = \alpha$.
Ici, $g' : X_{T'} \to X_T$ est le changement de base de $g$.
\end{enumerate}
Il existe un foncteur d'oubli naturel fidèle
$$
\Picardstack_{X/B, \sigma} \longrightarrow
\Picardstack_{X/B}
$$
Nous considérons ainsi $\Picardstack_{X/B, \sigma}$ comme une catégorie
sur $(\Sch/S)_{fppf}$.

\begin{lemma}
\label{lemma-pic-with-section-stack}
Dans la situation \ref{situation-pic}, soit $\sigma : B \to X$ une section.
Alors $\Picardstack_{X/B, \sigma}$, défini ci-dessus, est un champ en
groupoïdes sur $(\Sch/S)_{fppf}$.
\end{lemma}

\begin{proof}
Nous savons déjà que $\Picardstack_{X/B}$ est un champ en groupoïdes
sur $(\Sch/S)_{fppf}$, d'après
Exemples de champs, lemme \ref{examples-stacks-lemma-picard-stack}.
Montrons la descente des objets pour $\Picardstack_{X/B, \sigma}$.
Soit $\{T_i \to T\}$ un recouvrement fppf, soit
$\xi_i = (T_i, h_i, \mathcal{L}_i, \alpha_i)$ un objet de
$\Picardstack_{X/B, \sigma}$ au-dessus de $T_i$, et soit
$\varphi_{ij} : \text{pr}_0^*\xi_i \to \text{pr}_1^*\xi_j$
une donnée de descente. En appliquant le résultat pour $\Picardstack_{X/B}$,
nous voyons que nous pouvons supposer donné un objet $(T, h, \mathcal{L})$
de $\Picardstack_{X/B}$ au-dessus de $T$ dont l'image inverse est $\xi_i$
pour tout $i$. Nous obtenons alors
$$
\alpha_i : \mathcal{O}_{T_i} \to \sigma_{T_i}^*\mathcal{L}_i =
(T_i \to T)^*\sigma_T^*\mathcal{L}
$$
Puisque les applications $\varphi_{ij}$ sont compatibles aux $\alpha_i$,
nous voyons que les images inverses de $\alpha_i$ et $\alpha_j$ sont la même
application sur $T_i \times_T T_j$. Par descente des faisceaux quasi-cohérents
(Descente, proposition \ref{descent-proposition-fpqc-descent-quasi-coherent},
nous voyons que les $\alpha_i$ sont les restrictions d'une unique application
$\alpha : \mathcal{O}_T \to \sigma_T^*\mathcal{L}$, comme voulu.
Nous omettons la démonstration de la descente des morphismes.
\end{proof}

\begin{lemma}
\label{lemma-compare-pic-with-section}
Dans la situation \ref{situation-pic}, soit $\sigma : B \to X$ une section.
Le morphisme $\Picardstack_{X/B, \sigma} \to \Picardstack_{X/B}$
est représentable, surjectif et lisse.
\end{lemma}

\begin{proof}
Soit $T$ un schéma et soit $(\Sch/T)_{fppf} \to \Picardstack_{X/B}$
donné par l'objet $\xi = (T, h, \mathcal{L})$ de $\Picardstack_{X/B}$
au-dessus de $T$. Nous devons montrer que
$$
(\Sch/T)_{fppf} \times_{\xi, \Picardstack_{X/B}} \Picardstack_{X/B, \sigma}
$$
est représentable par un schéma $V$ et que le morphisme correspondant
$V \to T$ est surjectif et lisse. Voir
Champs algébriques, sections \ref{algebraic-section-representable-morphism},
\ref{algebraic-section-morphisms-representable-by-algebraic-spaces} et
\ref{algebraic-section-representable-properties}.
Le foncteur d'oubli $\Picardstack_{X/B, \sigma} \to \Picardstack_{X/B}$
est fidèle sur les catégories fibres et, pour $T'/T$, l'ensemble des classes
d'isomorphisme est l'ensemble des isomorphismes
$$
\alpha' : \mathcal{O}_{T'} \longrightarrow (T' \to T)^*\sigma_T^*\mathcal{L}
$$
Voir Champs algébriques, lemme
\ref{algebraic-lemma-criterion-map-representable-spaces-fibred-in-groupoids}.
Nous savons que ce foncteur est représentable par un schéma affine $U$ de
présentation finie sur $T$, d'après la proposition \ref{proposition-isom}
(appliquée à $\text{id} : T \to T$, à $\mathcal{O}_T$ et à
$\sigma^*\mathcal{L}$). En travaillant localement pour la topologie de Zariski
sur $T$, nous pouvons supposer que $\sigma_T^*\mathcal{L}$ est isomorphe à
$\mathcal{O}_T$ ; nous voyons alors que notre foncteur est représentable par
$\mathbf{G}_m \times T$ sur $T$. Ainsi,
$U \to T$ est, localement pour la topologie de Zariski sur $T$, de la forme
de la projection $\mathbf{G}_m \times T \to T$, laquelle est bien lisse et
surjective.
\end{proof}

\begin{lemma}
\label{lemma-flat-geometrically-connected-fibres-with-section-functor-stack}
Dans la situation \ref{situation-pic}, soit $\sigma : B \to X$ une section.
Si $\mathcal{O}_T \to f_{T, *}\mathcal{O}_{X_T}$ est un isomorphisme
pour tout $T$ sur $B$, alors
$\Picardstack_{X/B, \sigma} \to (\Sch/S)_{fppf}$
est fibré en sétoïdes et l'ensemble de ses classes d'isomorphisme au-dessus de
$T$ est donné par
$$
\coprod\nolimits_{h : T \to B}
\Ker(\sigma_T^* : \Pic(X \times_{B, h} T) \to \Pic(T))
$$
\end{lemma}

\begin{proof}
Si $\xi = (T, h, \mathcal{L}, \alpha)$ est un objet de
$\Picardstack_{X/B, \sigma}$ au-dessus de $T$, alors un automorphisme
$\varphi$ de $\xi$ est donné par la multiplication par une section globale
inversible $u$ du faisceau structural de $X_T$ telle qu'en outre
$\sigma_T^*u = 1$. Alors $u = 1$ par notre hypothèse selon laquelle
$\mathcal{O}_T \to f_{T, *}\mathcal{O}_{X_T}$ est un isomorphisme.
Ainsi, $\Picardstack_{X/B, \sigma}$ est fibré en sétoïdes sur
$(\Sch/S)_{fppf}$. Étant donnés $T$ et $h : T \to B$, l'ensemble des classes
d'isomorphisme de couples $(\mathcal{L}, \alpha)$ est le même que l'ensemble
des classes d'isomorphisme des $\mathcal{L}$ tels que
$\sigma_T^*\mathcal{L} \cong \mathcal{O}_T$ (sans spécification de
l'isomorphisme). Cela est clair, car deux choix de $\alpha$ diffèrent par une
unité globale sur $T$, ce qui revient au même qu'une unité globale sur $X_T$.
\end{proof}

\begin{proposition}
\label{proposition-pic-functor}
Soit $S$ un schéma. Soit $f : X \to B$ un morphisme d'espaces algébriques
sur $S$. Supposons que
\begin{enumerate}
\item $f$ soit plat, de présentation finie et propre, et
\item $\mathcal{O}_T \to f_{T, *}\mathcal{O}_{X_T}$ soit un isomorphisme
pour tout schéma $T$ sur $B$.
\end{enumerate}
Alors $\Picardfunctor_{X/B}$ est un espace algébrique.
\end{proposition}

\noindent
Dans la situation de la proposition, le champ algébrique
$\Picardstack_{X/B}$ est une gerbe au-dessus de l'espace algébrique
$\Picardfunctor_{X/B}$. Après développement de la théorie générale des gerbes,
cela fournit une démonstration plus courte de la proposition (mais faisant
appel à une théorie plus générale).

\begin{proof}
Il existe un morphisme surjectif, plat et de présentation finie
$B' \to B$ d'espaces algébriques tel que le changement de base
$X' = X \times_B B'$ sur $B'$ possède une section : nous pouvons en effet
prendre $B' = X$.
Observons que $\Picardfunctor_{X'/B'} = B' \times_B \Picardfunctor_{X/B}$.
Ainsi, $\Picardfunctor_{X'/B'} \to \Picardfunctor_{X/B}$ est représentable
par des espaces algébriques, surjectif, plat et de présentation finie.
Par conséquent, si nous pouvons montrer que $\Picardfunctor_{X'/B'}$ est un
espace algébrique, il s'ensuit que $\Picardfunctor_{X/B}$ est un espace
algébrique d'après Bootstrap, théorème
\ref{bootstrap-theorem-final-bootstrap}.
Nous nous ramenons ainsi au cas décrit dans le paragraphe suivant.

\medskip\noindent
En plus des hypothèses de la proposition, supposons donnée une section
$\sigma : B \to X$. D'après la proposition \ref{proposition-pic},
$\Picardstack_{X/B}$ est un champ algébrique.
D'après le lemme \ref{lemma-compare-pic-with-section} et
Champs algébriques, lemme
\ref{algebraic-lemma-representable-morphism-to-algebraic},
$\Picardstack_{X/B, \sigma}$ est un champ algébrique.
D'après le lemme
\ref{lemma-flat-geometrically-connected-fibres-with-section-functor-stack}
et Champs algébriques, lemme
\ref{algebraic-lemma-characterize-representable-by-space},
nous voyons que $T \mapsto \Ker(\sigma_T^* : \Pic(X_T) \to \Pic(T))$
est un espace algébrique.
D'après le lemme \ref{lemma-flat-geometrically-connected-fibres-with-section},
ce foncteur est le même que $\Picardfunctor_{X/B}$.
\end{proof}

\begin{lemma}
\label{lemma-diagonal-pic}
Sous les hypothèses et avec les notations de la proposition
\ref{proposition-pic-functor}.
Alors la diagonale
$\Picardfunctor_{X/B} \to \Picardfunctor_{X/B} \times_B \Picardfunctor_{X/B}$
est représentable par des immersions. Autrement dit,
$\Picardfunctor_{X/B} \to B$ est localement séparé.
\end{lemma}

\begin{proof}
Soit $T$ un schéma sur $B$ et soient $s, t \in \Picardfunctor_{X/B}(T)$.
Nous voulons montrer qu'il existe un sous-schéma localement fermé
$Z \subset T$ tel que $s|_Z = t|_Z$ et tel qu'un morphisme $T' \to T$ se
factorise par $Z$ si et seulement si $s|_{T'} = t|_{T'}$.

\medskip\noindent
Nous commençons par ramener le problème général au cas où $s$ et $t$
proviennent de modules inversibles sur $X_T$. Nous suggérons au lecteur de
sauter cette étape. Choisissons un recouvrement fppf
$\{T_i \to T\}_{i \in I}$ tel que $s|_{T_i}$ et $t|_{T_i}$ proviennent de
$\Pic(X_{T_i})$ pour tout $i$.
Supposons que nous puissions démontrer le résultat pour tous les couples
$s|_{T_i}, t|_{T_i}$. Nous obtenons alors des sous-schémas localement fermés
$Z_i \subset T_i$ possédant la propriété universelle voulue.
Il s'ensuit que $Z_i$ et $Z_j$ ont la même image inverse schématique dans
$T_i \times_T T_j$.
Cela détermine une donnée de descente sur $Z_i/T_i$.
Comme $Z_i \to T_i$ est localement quasi-fini, il résulte de
Compléments sur les morphismes, lemme
\ref{more-morphisms-lemma-separated-locally-quasi-finite-morphisms-fppf-descend}
que nous obtenons un morphisme localement quasi-fini $Z \to T$ dont le
changement de base redonne $Z_i \to T_i$. Alors $Z \to T$ est une immersion
d'après Descente, lemme
\ref{descent-lemma-descending-fppf-property-immersion}.
Enfin, puisque $\Picardfunctor_{X/B}$ est un faisceau fppf, nous concluons que
$s|_Z = t|_Z$ et que $Z$ satisfait à la propriété universelle mentionnée
ci-dessus.

\medskip\noindent
Supposons que $s$ et $t$ proviennent de modules inversibles
$\mathcal{V}$, $\mathcal{W}$ sur $X_T$.
Posons $\mathcal{L} = \mathcal{V} \otimes \mathcal{W}^{\otimes -1}$
Nous cherchons un sous-schéma localement fermé $Z$ de $T$ tel que
$T' \to T$ se factorise par $Z$ si et seulement si
$\mathcal{L}_{X_{T'}}$ est l'image inverse d'un faisceau inversible sur $T'$,
voir le lemme \ref{lemma-flat-geometrically-connected-fibres}.
L'existence de $Z$ résulte donc de
Compléments sur les morphismes d'espaces, lemme
\ref{spaces-more-morphisms-lemma-diagonal-picard-flat-proper}.
\end{proof}

\section{Morphismes relatifs}
\label{section-relative-morphisms}

\noindent
Nous poursuivons la discussion de Critères de représentabilité, section
\ref{criteria-section-relative-morphisms}.
Dans cette section, partant d'un schéma $S$ et de morphismes d'espaces
algébriques $Z \to B$ et $X \to B$ sur $S$, nous avons construit un foncteur
$$
\mathit{Mor}_B(Z, X) : (\Sch/B)^{opp} \longrightarrow \textit{Sets}, \quad
T \longmapsto \{f : Z_T \to X_T\}
$$
Nous considérons parfois $\mathit{Mor}_B(Z, X)$ comme un foncteur
$(\Sch/S)^{opp} \to \textit{Sets}$ muni d'un morphisme
$\mathit{Mor}_B(Z, X) \to B$.
En effet, si $T$ est un schéma sur $S$, alors un élément de
$\mathit{Mor}_B(Z, X)(T)$ est un couple $(f, h)$, où $h$ est un morphisme
$h : T \to B$ et où
$f : Z \times_{B, h} T \to X \times_{B, h} T$
est un morphisme d'espaces algébriques sur $T$. En particulier, lorsque nous
disons que $\mathit{Mor}_B(Z, X)$ est un espace algébrique, nous entendons que
le foncteur correspondant $(\Sch/S)^{opp} \to \textit{Sets}$ est un espace
algébrique.

\begin{lemma}
\label{lemma-Mor-into-Hilb}
Soit $S$ un schéma. Considérons des morphismes d'espaces algébriques
$Z \to B$ et $X \to B$ sur $S$.
Si $X \to B$ est séparé et si $Z \to B$ est de présentation finie,
plat et propre, alors il existe une transformation injective naturelle de
foncteurs
$$
\mathit{Mor}_B(Z, X) \longrightarrow \Hilbfunctor_{Z \times_B X/B}
$$
qui envoie un morphisme $f : Z_T \to X_T$ sur son graphe.
\end{lemma}

\begin{proof}
Étant donné un schéma $T$ sur $B$ et un morphisme $f_T : Z_T \to X_T$
sur $T$, le graphe de $f$ est le morphisme
$\Gamma_f = (\text{id}, f) : Z_T \to Z_T \times_T X_T = (Z \times_B X)_T$.
Rappelons que les propriétés d'être séparé, plat, propre ou de présentation
finie sont des propriétés des morphismes d'espaces algébriques qui sont stables
par changement de base (Morphismes d'espaces, lemmes
\ref{spaces-morphisms-lemma-base-change-separated},
\ref{spaces-morphisms-lemma-base-change-flat},
\ref{spaces-morphisms-lemma-base-change-proper} et
\ref{spaces-morphisms-lemma-base-change-finite-presentation}).
Ainsi, $\Gamma_f$ est une immersion fermée d'après
Morphismes d'espaces, lemme \ref{spaces-morphisms-lemma-semi-diagonal}.
De plus, $\Gamma_f(Z_T)$ est plat, propre et de présentation finie sur $T$.
Donc $\Gamma_f(Z_T)$ définit un élément de
$\Hilbfunctor_{Z \times_B X/B}(T)$.
Pour montrer que la transformation est injective, il suffit de montrer que deux
morphismes ayant le même graphe sont égaux. Cela est vrai parce que, si
$Y \subset (Z \times_B X)_T$ est le graphe d'un morphisme $f$, nous pouvons
retrouver $f$ en composant l'inverse de
$\text{pr}_1|_Y : Y \to Z_T$ avec $\text{pr}_2|_Y$.
\end{proof}

\begin{lemma}
\label{lemma-Mor-into-Hilb-open}
Hypothèses et notations comme dans le lemme \ref{lemma-Mor-into-Hilb}.
La transformation
$\mathit{Mor}_B(Z, X) \longrightarrow \Hilbfunctor_{Z \times_B X/B}$
est représentable par des immersions ouvertes.
\end{lemma}

\begin{proof}
Soit $T$ un schéma sur $B$ et soit $Y \subset (Z \times_B X)_T$
un élément de $\Hilbfunctor_{Z \times_B X/B}(T)$. Nous voyons alors que
$Y$ est le graphe d'un morphisme $Z_T \to X_T$ sur $T$ si et seulement si
$k = \text{pr}_1|_Y : Y \to Z_T$ est un isomorphisme. D'après
Compléments sur les morphismes d'espaces, lemme
\ref{spaces-more-morphisms-lemma-where-isomorphism}, il existe un sous-schéma
ouvert $V \subset T$ tel que, pour tout morphisme de schémas $T' \to T$,
$k_{T'} : Y_{T'} \to Z_{T'}$ soit un isomorphisme si et seulement si
$T' \to T$ se factorise par $V$.
Cela prouve le lemme.
\end{proof}

\begin{proposition}
\label{proposition-Mor}
Soit $S$ un schéma. Soient $Z \to B$ et $X \to B$ des morphismes d'espaces
algébriques sur $S$. Supposons que $X \to B$ soit de présentation finie et
séparé, et que $Z \to B$ soit de présentation finie, plat et propre. Alors
$\mathit{Mor}_B(Z, X)$ est un espace algébrique localement de présentation
finie sur $B$.
\end{proposition}

\begin{proof}
Conséquence immédiate du lemme \ref{lemma-Mor-into-Hilb-open}
et de la proposition \ref{proposition-hilb}.
\end{proof}







\section{Le champ des espaces algébriques}
\label{section-stack-of-spaces}

\noindent
Cette section poursuit la discussion commencée dans
Exemples de champs, sections
\ref{examples-stacks-section-stack-of-spaces},
\ref{examples-stacks-section-stack-of-finite-type-spaces} et
\ref{examples-stacks-section-stack-in-groupoids-of-finite-type-spaces}.
En travaillant sur $\mathbf{Z}$, la discussion qui s'y trouve montre que nous
avons un champ en groupoïdes
$$
p'_{ft} : \Spacesstack'_{ft} \longrightarrow \Sch_{fppf}
$$
qui paramètre les familles (non plates) d'espaces algébriques de type fini.
Plus précisément, un objet\footnote{Nous effectuons toujours un remplacement
comme dans Exemples de champs, lemme
\ref{examples-stacks-lemma-stack-ft-spaces}.}
de $\Spacesstack'_{ft}$ est un morphisme de type fini $X \to S$
d'un espace algébrique $X$ vers un schéma $S$, et un morphisme
$(X' \to S') \to (X \to S)$ est donné par un couple $(f, g)$, où
$f : X' \to X$ est un morphisme d'espaces algébriques et
$g : S' \to S$ est un morphisme de schémas, qui s'inscrivent dans un diagramme
commutatif
$$
\xymatrix{
X' \ar[d] \ar[r]_f & X \ar[d] \\
S' \ar[r]^g & S
}
$$
induisant un isomorphisme $X' \to S' \times_S X$ ; autrement dit, le
diagramme est cartésien dans la catégorie des espaces algébriques.
Le foncteur $p'_{ft}$ envoie $(X \to S)$ sur $S$ et $(f, g)$ sur $g$.
Nous définissons une sous-catégorie pleine
$$
\Spacesstack'_{fp, flat, proper} \subset
\Spacesstack'_{ft}
$$
formée des objets $X \to S$ de $\Spacesstack'_{ft}$ tels que
$X \to S$ soit de présentation finie, plat et propre.
Nous notons
$$
p'_{fp, flat, proper} :
\Spacesstack'_{fp, flat, proper}
\longrightarrow
\Sch_{fppf}
$$
la restriction du foncteur $p'_{ft}$ à la sous-catégorie indiquée.
Nous rappelons d'abord les résultats déjà obtenus dans les références
ci-dessus, puis nous commençons à y ajouter d'autres résultats.

\begin{lemma}
\label{lemma-spaces-fibred-in-groupoids}
La catégorie $\Spacesstack'_{ft}$ est fibrée en groupoïdes
sur $\Sch_{fppf}$. Il en est de même de
$\Spacesstack'_{fp, flat, proper}$.
\end{lemma}

\begin{proof}
Nous l'avons vu dans
Exemples de champs, section
\ref{examples-stacks-section-stack-in-groupoids-of-finite-type-spaces}
pour $\Spacesstack'_{ft}$, et cela entraîne aisément le résultat dans l'autre
cas. Cependant, démontrons-le aussi directement en vérifiant les conditions
(1) et (2) de Catégories, définition
\ref{categories-definition-fibred-groupoids}.

\medskip\noindent
Condition (1). Soit $X \to S$ un objet de $\Spacesstack'_{ft}$ et soit
$S' \to S$ un morphisme de schémas. Posons alors $X' = S' \times_S X$.
Notons que $X' \to S'$ est de type fini d'après
Morphismes d'espaces, lemme
\ref{spaces-morphisms-lemma-base-change-finite-type},
ce qui donne un morphisme $(X' \to S') \to (X \to S)$
au-dessus de $S' \to S$.
On raisonne de même dans l'autre cas en utilisant
Morphismes d'espaces, lemmes
\ref{spaces-morphisms-lemma-base-change-finite-presentation},
\ref{spaces-morphisms-lemma-base-change-flat} et
\ref{spaces-morphisms-lemma-base-change-proper}.

\medskip\noindent
Condition (2). Considérons des morphismes
$(f, g) : (X' \to S') \to (X \to S)$ et
$(a, b) : (Y \to T) \to (X \to S)$ de $\Spacesstack'_{ft}$.
Étant donné un morphisme $h : T \to S'$ tel que $g \circ h = b$, nous devons
montrer qu'il existe un unique morphisme
$(k, h) : (Y \to T) \to (X' \to S')$ de $\Spacesstack'_{ft}$ tel que
$(f, g) \circ (k, h) = (a, b)$.
Cela résulte clairement de l'égalité $X' = S' \times_S X$.
Le même raisonnement vaut donc pour toute sous-catégorie pleine de
$\Spacesstack'_{ft}$ qui satisfait à (1).
\end{proof}

\begin{lemma}
\label{lemma-spaces-diagonal}
La diagonale
$$
\Delta : \Spacesstack'_{fp, flat, proper} \longrightarrow
\Spacesstack'_{fp, flat, proper} \times \Spacesstack'_{fp, flat, proper}
$$
est représentable par des espaces algébriques.
\end{lemma}

\begin{proof}
Nous utiliserons le critère (2) de
Champs algébriques, lemme \ref{algebraic-lemma-representable-diagonal}.
Soit $S$ un schéma et soient $X$ et $Y$ des espaces algébriques de présentation
finie sur $S$, plats sur $S$ et propres sur $S$. Nous devons montrer que le
foncteur
$$
\mathit{Isom}_S(X, Y) : (\Sch/S)_{fppf} \longrightarrow \textit{Sets}, \quad
T \longmapsto \{f : X_T \to Y_T \text{ est un isomorphisme}\}
$$
est un espace algébrique. Un argument élémentaire montre que
$\mathit{Isom}_S(X, Y)$ s'insère dans un produit fibré
$$
\xymatrix{
\mathit{Isom}_S(X, Y) \ar[r] \ar[d] & S \ar[d]_{(\text{id}, \text{id})} \\
\mathit{Mor}_S(X, Y) \times \mathit{Mor}_S(Y, X) \ar[r] &
\mathit{Mor}_S(X, X) \times \mathit{Mor}_S(Y, Y)
}
$$
La flèche du bas envoie $(\varphi, \psi)$ sur
$(\psi \circ \varphi, \varphi \circ \psi)$.
D'après la proposition \ref{proposition-Mor}, les foncteurs de la ligne
inférieure sont des espaces algébriques sur $S$.
Le résultat découle donc du fait que la catégorie des espaces algébriques sur
$S$ admet les produits fibrés.
\end{proof}

\begin{lemma}
\label{lemma-spaces-stack}
La catégorie $\Spacesstack'_{ft}$ est un champ en groupoïdes
sur $\Sch_{fppf}$. Il en est de même de
$\Spacesstack'_{fp, flat, proper}$.
\end{lemma}

\begin{proof}
La raison pour laquelle ce lemme est vrai tient dans le slogan suivant : toute
donnée de descente fppf pour les espaces algébriques est effective, voir
Bootstrap, section \ref{bootstrap-section-applications}.
Plus précisément, le lemme pour $\Spacesstack'_{ft}$ découle de
Exemples de champs, lemme
\ref{examples-stacks-lemma-stack-of-finite-type-spaces},
comme nous l'avons vu dans Exemples de champs, section
\ref{examples-stacks-section-stack-in-groupoids-of-finite-type-spaces}.
Cependant, rappelons la démonstration. Nous devons vérifier les conditions
(1), (2) et (3) de Champs, définition
\ref{stacks-definition-stack-in-groupoids}.

\medskip\noindent
Nous avons vu la propriété (1) au lemme
\ref{lemma-spaces-fibred-in-groupoids}.

\medskip\noindent
La propriété (2) résulte du lemme \ref{lemma-spaces-diagonal} dans le cas de
$\Spacesstack'_{fp, flat, proper}$.
Dans le cas de $\Spacesstack'_{ft}$, elle découle de
Exemples de champs, lemme
\ref{examples-stacks-lemma-pre-stack-of-spaces}
(et c'est réellement la référence « correcte »).

\medskip\noindent
La condition (3) pour $\Spacesstack'_{ft}$ se vérifie comme suit. Supposons
donnés
\begin{enumerate}
\item un recouvrement fppf $\{U_i \to U\}_{i \in I}$ dans $\Sch_{fppf}$,
\item pour tout $i \in I$, un espace algébrique $X_i$ de type fini sur
$U_i$, et
\item pour tous $i, j \in I$, un isomorphisme
$\varphi_{ij} : X_i \times_U U_j \to U_i \times_U X_j$ d'espaces algébriques
sur $U_i \times_U U_j$, satisfaisant à la condition de cocycle sur
$U_i \times_U U_j \times_U U_k$.
\end{enumerate}
Nous devons montrer qu'il existe un espace algébrique $X$ de type fini sur $U$
et des isomorphismes $X_{U_i} \cong X_i$ sur $U_i$ redonnant les
isomorphismes $\varphi_{ij}$. Cela résulte de
Bootstrap, lemme \ref{bootstrap-lemma-descend-algebraic-space}, partie (2).
D'après Descente sur les espaces, lemme
\ref{spaces-descent-lemma-descending-property-finite-type},
$X \to U$ est de type fini.
Dans le cas de $\Spacesstack'_{fp, flat, proper}$, on utilise en outre
Descente sur les espaces, lemmes
\ref{spaces-descent-lemma-descending-property-finite-presentation},
\ref{spaces-descent-lemma-descending-property-flat} et
\ref{spaces-descent-lemma-descending-property-proper}
à la dernière étape.
\end{proof}

\noindent
Vérification de cohérence : les champs $\Spacesstack'_{ft}$ et
$\Spacesstack'_{fp, flat, proper}$ jouent le même rôle parmi les espaces
algébriques.

\begin{lemma}
\label{lemma-extend-spaces-to-spaces}
Soit $T$ un espace algébrique sur $\mathbf{Z}$. Soit $\mathcal{S}_T$
le champ algébrique correspondant (Champs algébriques, sections
\ref{algebraic-section-split},
\ref{algebraic-section-representable-by-algebraic-spaces} et
\ref{algebraic-section-stacks-spaces}).
Nous avons une équivalence de catégories
$$
\left\{
\begin{matrix}
\text{morphismes d'espaces algébriques }\\
X \to T\text{ de type fini}
\end{matrix}
\right\}
\longrightarrow
\Mor_{\textit{Cat}/\Sch_{fppf}}(\mathcal{S}_T, \Spacesstack'_{ft})
$$
et une équivalence de catégories
$$
\left\{
\begin{matrix}
\text{morphismes d'espaces algébriques }X \to T\\
\text{de présentation finie, plats et propres}
\end{matrix}
\right\}
\longrightarrow
\Mor_{\textit{Cat}/\Sch_{fppf}}(\mathcal{S}_T,
\Spacesstack'_{fp, flat, proper})
$$
\end{lemma}

\begin{proof}
Nous allons déduire ce lemme du fait qu'il vaut pour les schémas
(essentiellement par construction des champs) et du fait que les données de
descente fppf pour les espaces algébriques sur des espaces algébriques sont
effectives. Nous encourageons vivement le lecteur à omettre la démonstration.

\medskip\noindent
La construction de gauche à droite dans chacune des flèches est immédiate :
étant donné $X \to T$ de type fini, le foncteur
$\mathcal{S}_T \to  \Spacesstack'_{ft}$ associe à $U/T$ le changement de base
$X_U \to U$. Nous allons expliquer comment construire un quasi-inverse.

\medskip\noindent
Si $T$ est un schéma, il existe un quasi-inverse d'après le lemme de
$2$-Yoneda, voir Catégories, lemme
\ref{categories-lemma-yoneda-2category}.
Soit $p : U \to T$ un morphisme surjectif \'etale, où $U$ est un schéma.
Soit $R = U \times_T U$, de projections $s, t : R \to U$.
Observons que nous obtenons des morphismes
$$
\xymatrix{
\mathcal{S}_{U \times_T U \times_T U} \ar@<2ex>[r] \ar[r] \ar@<-2ex>[r] &
\mathcal{S}_R \ar@<1ex>[r] \ar@<-1ex>[r] &
\mathcal{S}_U \ar[r] &
\mathcal{S}_T
}
$$
qui satisfont strictement à diverses compatibilités.

\medskip\noindent
Soit $G : \mathcal{S}_T \to \Spacesstack'_{ft}$ un foncteur sur
$\Sch_{fppf}$. La restriction de $G$ à $\mathcal{S}_U$ par l'application
ci-dessus correspond, par le lemme de $2$-Yoneda, à un morphisme de type fini
$X_U \to U$ d'espaces algébriques. Puisque $p \circ s = p \circ t$, nous
voyons que $R \times_{s, U} X_U$ et $R \times_{t, U} X_U$ correspondent tous
deux à la restriction de $G$ à $\mathcal{S}_R$. Nous obtenons donc un
isomorphisme canonique
$\varphi : X_U \times_{U, t} R \to R \times_{s, U} X_U$ sur $R$.
Cet isomorphisme satisfait à la condition de cocycle en vertu des diverses
compatibilités du diagramme ci-dessus.
Nous avons ainsi une donnée de descente qui est effective d'après
Bootstrap, lemme \ref{bootstrap-lemma-descend-algebraic-space}, partie (2).
Autrement dit, nous obtenons un objet $X \to T$ de la catégorie du membre de
droite. Nous omettons de vérifier que la construction $G \leadsto X$ est
fonctorielle et qu'elle est quasi-inverse de l'autre construction.
Dans le cas de $\Spacesstack'_{fp, flat, proper}$, on utilise en outre
Descente sur les espaces, lemmes
\ref{spaces-descent-lemma-descending-property-finite-presentation},
\ref{spaces-descent-lemma-descending-property-flat} et
\ref{spaces-descent-lemma-descending-property-proper}
à la dernière étape pour voir que $X \to T$ est de présentation finie,
plat et propre.
\end{proof}

\begin{remark}
\label{remark-spaces-base-change}
Soit $B$ un espace algébrique sur $\Spec(\mathbf{Z})$.
Soit $B\textit{-Spaces}'_{ft}$ la catégorie formée des couples
$(X \to S, h : S \to B)$, où $X \to S$ est un objet de
$\Spacesstack'_{ft}$ et $h : S \to B$ est un morphisme.
Un morphisme $(X' \to S', h') \to (X \to S, h)$ dans
$B\textit{-Spaces}'_{ft}$ est un morphisme $(f, g)$ dans
$\Spacesstack'_{ft}$ tel que $h \circ g = h'$.
Dans cette situation, le diagramme
$$
\xymatrix{
B\textit{-Spaces}'_{ft} \ar[r] \ar[d] & \Spacesstack'_{ft} \ar[d] \\
(\Sch/B)_{fppf} \ar[r] & \Sch_{fppf}
}
$$
est un carré $2$-produit fibré. Cette remarque triviale sera parfois utile
pour déduire des résultats du cas absolu $\Spacesstack'_{ft}$ au cas des
familles sur un espace algébrique de base donné.
Bien entendu, une construction analogue vaut pour
$B\textit{-Spaces}'_{fp, flat, proper}$
\end{remark}

\begin{lemma}
\label{lemma-spaces-limits}
Le champ
$p'_{fp, flat, proper} :
\Spacesstack'_{fp, flat, proper} \to \Sch_{fppf}$ commute aux limites
(Axiomes d'Artin, définition \ref{artin-definition-limit-preserving}).
\end{lemma}

\begin{proof}
Soit $T = \lim T_i$ la limite d'un système projectif filtrant de schémas
affines. D'après Limites d'espaces, lemme
\ref{spaces-limits-lemma-descend-finite-presentation},
la catégorie des espaces algébriques de présentation finie sur $T$ est la
limite inductive des catégories d'espaces algébriques de présentation finie
sur les $T_i$.
Pour achever la démonstration, utilisons le fait que la platitude et la
propreté descendent à travers la limite, voir
Limites d'espaces, lemmes
\ref{spaces-limits-lemma-descend-flat} et
\ref{spaces-limits-lemma-eventually-proper}.
\end{proof}

\begin{lemma}
\label{lemma-spaces-RS-star}
Soit
$$
\xymatrix{
T \ar[r] \ar[d] & T' \ar[d] \\
S \ar[r] & S'
}
$$
une somme amalgamée dans la catégorie des schémas, où
$T \to T'$ est un épaississement et $T \to S$ est affine, voir
Compléments sur les morphismes, lemme
\ref{more-morphisms-lemma-pushout-along-thickening}.
Alors le foncteur sur les catégories fibres
$$
\begin{matrix}
\Spacesstack'_{fp, flat, proper, S'} \\
\downarrow \\
\Spacesstack'_{fp, flat, proper, S}
\times_{\Spacesstack'_{fp, flat, proper, T}}
\Spacesstack'_{fp, flat, proper, T'}
\end{matrix}
$$
est une équivalence.
\end{lemma}

\begin{proof}
Le foncteur est une équivalence si l'on retranche « propre » de la liste des
conditions et que l'on remplace « de présentation finie » par
« localement de présentation finie », voir Sommes amalgamées d'espaces, lemme
\ref{spaces-pushouts-lemma-equivalence-categories-spaces-pushout-flat}.
Il suffit donc de montrer qu'étant donné un morphisme $X' \to S'$ d'un espace
algébrique vers $S'$ qui est plat et localement de présentation finie,
$X' \to S'$ est propre si et seulement si
$S \times_{S'} X' \to S$ et $T' \times_{S'} X' \to T'$ sont propres.
Une implication résulte du fait que la propreté est préservée par changement de
base (Morphismes d'espaces, lemme
\ref{spaces-morphisms-lemma-base-change-proper}) et l'autre, du fait que la
propreté de $S \times_{S'} X' \to S$ entraîne celle de $X' \to S'$, d'après
Compléments sur les morphismes d'espaces, lemme
\ref{spaces-more-morphisms-lemma-thicken-property-morphisms-cartesian}.
\end{proof}

\begin{lemma}
\label{lemma-spaces-tangent-space}
Soit $k$ un corps et soit $x = (X \to \Spec(k))$ un objet de
$\mathcal{X} = \Spacesstack'_{fp, flat, proper}$ au-dessus de $\Spec(k)$.
\begin{enumerate}
\item Si $k$ est de type fini sur $\mathbf{Z}$, alors les espaces vectoriels
$T\mathcal{F}_{\mathcal{X}, k, x}$ et
$\text{Inf}(\mathcal{F}_{\mathcal{X}, k, x})$
(voir Axiomes d'Artin, section \ref{artin-section-tangent-spaces})
sont de dimension finie, et
\item en général, les espaces vectoriels $T_x(k)$ et $\text{Inf}_x(k)$
(voir Axiomes d'Artin, section \ref{artin-section-inf})
sont de dimension finie.
\end{enumerate}
\end{lemma}

\begin{proof}
La discussion d'Axiomes d'Artin, section \ref{artin-section-tangent-spaces},
ne s'applique qu'aux corps de type fini sur le schéma de base
$\Spec(\mathbf{Z})$.
Notre champ satisfait à (RS*) d'après le lemme \ref{lemma-spaces-RS-star},
et nous pouvons appliquer
Axiomes d'Artin, lemme \ref{artin-lemma-properties-lift-RS-star},
pour obtenir les espaces vectoriels $T_x(k)$ et $\text{Inf}_x(k)$ mentionnés en
(2). De plus, dans le cas de type fini, ces espaces coïncident avec ceux
mentionnés en (1), d'après Axiomes d'Artin, remarque
\ref{artin-remark-compare-deformation-spaces}.
Cela étant réglé, nous pouvons commencer la démonstration.
Observons que l'épaississement du premier ordre
$\Spec(k) \to \Spec(k[\epsilon]) = \Spec(k[k])$
a pour module conormal $k$. Ainsi, les formules d'après
Théorie des déformations, lemme \ref{defos-lemma-deform-spaces},
qui décrivent les déformations infinitésimales de $X$ et les automorphismes
infinitésimaux de $X$, deviennent
$$
T_x(k) = \Ext^1_{\mathcal{O}_X}(\NL_{X/k}, \mathcal{O}_X)
\quad\text{et}\quad
\text{Inf}_x(k) = \Ext^0_{\mathcal{O}_X}(\NL_{X/k}, \mathcal{O}_X)
$$
D'après Compléments sur les morphismes d'espaces, lemme
\ref{spaces-more-morphisms-lemma-netherlander-fp}, et le fait que $X$ est
noethérien, nous voyons que les faisceaux de cohomologie de $\NL_{X/k}$ sont
cohérents et nuls sauf aux degrés $0$ et $-1$.
D'après Catégories dérivées d'espaces, lemme
\ref{spaces-perfect-lemma-ext-finite}, les groupes $\Ext$ affichés sont des
espaces vectoriels de dimension finie sur $k$, ce qui achève la démonstration.
\end{proof}

\noindent
Prenons garde au fait que l'ouverture de la versalité (démontrée dans le lemme
suivant) est un peu étrange, car notre champ ne satisfait pas à l'effectivité
formelle, voir Exemples, section
\ref{examples-section-proper-spaces-not-algebraic}.
Plus tard, nous appliquerons l'ouverture de la versalité à des sous-champs
convenables de $\Spacesstack'_{fp, flat, proper}$ qui satisfont à l'effectivité
formelle, afin de conclure que ces champs sont algébriques.

\begin{lemma}
\label{lemma-spaces-defo-thy}
Le champ en groupoïdes
$\mathcal{X} = \Spacesstack'_{fp, flat, proper}$ satisfait à l'ouverture de la
versalité sur $\Spec(\mathbf{Z})$.
De même, après changement de base (remarque \ref{remark-spaces-base-change}),
l'ouverture de la versalité vaut sur tout schéma de base noethérien $S$.
\end{lemma}

\begin{proof}
Pour la démonstration « usuelle » de ce fait, nous renvoyons à la discussion de
la remarque qui suit cette démonstration. Nous allons le démontrer à l'aide de
Axiomes d'Artin, lemme \ref{artin-lemma-SGE-implies-openness-versality}.
Nous avons déjà vu que la diagonale de $\mathcal{X}$ est représentable par des
espaces algébriques, qu'il satisfait à (RS*) et qu'il commute aux
limites ; voir les lemmes \ref{lemma-spaces-diagonal},
\ref{lemma-spaces-RS-star} et
\ref{lemma-spaces-limits}.
Il nous suffit donc de voir que $\mathcal{X}$ satisfait à l'effectivité
formelle forte formulée dans
Axiomes d'Artin, lemme \ref{artin-lemma-SGE-implies-openness-versality}.

\medskip\noindent
Soit $(R_n)$ un système projectif d'anneaux tel que $R_n \to R_m$ soit
surjectif à noyau de carré nul pour tous $n \geq m$. Soit
$X_n \to \Spec(R_n)$ un morphisme de présentation finie, plat et propre, où
$X_n$ est un espace algébrique, et soit $X_{n + 1} \to X_n$ un morphisme
sur $\Spec(R_{n + 1})$ induisant un isomorphisme
$X_n = X_{n + 1} \times_{\Spec(R_{n + 1})} \Spec(R_n)$.
Nous devons trouver un morphisme plat, propre et de présentation finie
$X \to \Spec(\lim R_n)$ dont la source est un espace algébrique, tel que
$X_n$ soit le changement de base de $X$ pour tout $n$.

\medskip\noindent
Soit $I_n = \Ker(R_n \to R_1)$. Nous pouvons considérer
$(X_1 \subset X_n) \to (\Spec(R_1) \subset \Spec(R_n))$
comme un morphisme d'épaississements du premier ordre. (Nous prions le lecteur
de parcourir une partie du contenu consacré aux épaississements d'espaces
algébriques dans Compléments sur les morphismes d'espaces, section
\ref{spaces-more-morphisms-section-thickenings}, avant de poursuivre.)
Le faisceau structural de $X_n$ est une extension
$$
0 \to \mathcal{O}_{X_1} \otimes_{R_1} I_n \to
\mathcal{O}_{X_n} \to \mathcal{O}_{X_1} \to 0
$$
au-dessus de $0 \to I_n \to R_n \to R_1$, voir
Compléments sur les morphismes d'espaces, lemme
\ref{spaces-more-morphisms-lemma-deform}.
Considérons l'extension
$$
0 \to \lim \mathcal{O}_{X_1} \otimes_{R_1} I_n \to
\lim \mathcal{O}_{X_n} \to \mathcal{O}_{X_1} \to 0
$$
au-dessus de $0 \to \lim I_n \to \lim R_n \to R_1 \to 0$.
La suite affichée est exacte, car le $R^1\lim$ du système de noyaux est nul
d'après Catégories dérivées d'espaces, lemme
\ref{spaces-perfect-lemma-Rlim-quasi-coherent}.
Observons que l'application
$$
\mathcal{O}_{X_1} \otimes_{R_1} \lim I_n \longrightarrow
\lim \mathcal{O}_{X_1} \otimes_{R_1} I_n
$$
induit un isomorphisme après application du foncteur $DQ_X$, voir
Catégories dérivées d'espaces, lemme
\ref{spaces-perfect-lemma-pullback-and-limits}.
Nous obtenons donc une unique extension
$$
0 \to \mathcal{O}_{X_1} \otimes_{R_1} \lim I_n \to
\mathcal{O}' \to \mathcal{O}_{X_1} \to 0
$$
au-dessus de $0 \to \lim I_n \to \lim R_n \to R_1 \to 0$,
par l'équivalence de catégories de
Théorie des déformations, lemme
\ref{defos-lemma-thickening-over-thickening-space-quasi-coherent}.
Le faisceau $\mathcal{O}'$ détermine un épaississement du premier ordre
d'espaces algébriques $X_1 \subset X$ au-dessus de
$\Spec(R_1) \subset \Spec(\lim R_n)$, d'après
Compléments sur les morphismes d'espaces, lemme
\ref{spaces-more-morphisms-lemma-first-order-thickening}.
Observons que $X \to \Spec(\lim R_n)$ est plat d'après le lemme déjà utilisé
de Compléments sur les morphismes d'espaces,
\ref{spaces-more-morphisms-lemma-deform}.
D'après Compléments sur les morphismes d'espaces, lemme
\ref{spaces-more-morphisms-lemma-deform-property},
$X \to \Spec(\lim R_n)$ est propre et de présentation finie.
Cela achève la démonstration.
\end{proof}

\begin{remark}
\label{remark-spaces-defo-thy}
Le lemme \ref{lemma-spaces-defo-thy} peut aussi se démontrer soit à l'aide de
Axiomes d'Artin, lemme \ref{artin-lemma-dual-openness}
(comme dans la première démonstration du
lemme \ref{lemma-coherent-defo-thy}), soit à l'aide d'une théorie des
obstructions comme dans Axiomes d'Artin, lemme
\ref{artin-lemma-get-openness-obstruction-theory}
(comme dans la seconde démonstration du
lemme \ref{lemma-coherent-defo-thy}).
Dans les deux cas, on utilise la théorie des déformations et des obstructions
développée dans Cotangent, section
\ref{cotangent-section-deformations-ringed-topoi}, pour traduire les propriétés
nécessaires des déformations et des obstructions en groupes $\Ext$ auxquels
Catégories dérivées d'espaces, lemme
\ref{spaces-perfect-lemma-compute-ext}, peut être appliqué.
La seconde méthode (qui utilise une théorie des obstructions et donc le
complexe cotangent tout entier) est peut-être la méthode « standard » employée
dans la plupart des références.
\end{remark}

\section{Le champ des schémas propres polarisés}
\label{section-polarized}

\noindent
Pour étudier le champ des schémas propres polarisés, il suffit de travailler
sur $\mathbf{Z}$, puisque nous pourrons ensuite effectuer l'image inverse sur
tout schéma ou espace algébrique voulu (voir la remarque
\ref{remark-polarized-base-change}).

\begin{situation}
\label{situation-polarized}
Nous définissons une catégorie $\Polarizedstack$ comme suit. Ses objets sont
les couples $(X \to S, \mathcal{L})$ où
\begin{enumerate}
\item $X \to S$ est un morphisme de schémas propre, plat et de présentation
finie, et
\item $\mathcal{L}$ est un $\mathcal{O}_X$-module inversible relativement
ample sur $X/S$
(Morphismes, définition \ref{morphisms-definition-relatively-ample}).
\end{enumerate}
Un morphisme $(X' \to S', \mathcal{L}') \to (X \to S, \mathcal{L})$
entre objets est donné par un triplet $(f, g, \varphi)$, où
$f : X' \to X$ et $g : S' \to S$ sont des morphismes de schémas qui
s'inscrivent dans un diagramme commutatif
$$
\xymatrix{
X' \ar[d] \ar[r]_f & X \ar[d] \\
S' \ar[r]^g & S
}
$$
induisant un isomorphisme $X' \to S' \times_S X$ ; autrement dit, le
diagramme est cartésien, et
$\varphi : f^*\mathcal{L} \to \mathcal{L}'$ est un isomorphisme.
La composition est définie de manière évidente (voir
Exemples de champs, sections
\ref{examples-stacks-section-stack-of-spaces} et
\ref{examples-stacks-section-stack-of-quasi-coherent-sheaves}).
Le foncteur d'oubli
$$
p : \Polarizedstack \longrightarrow \Sch_{fppf},\quad
(X \to S, \mathcal{L}) \longmapsto S
$$
est celui par lequel nous considérons $\Polarizedstack$ comme une catégorie
sur $\Sch_{fppf}$
(voir la section \ref{section-conventions} pour les notations).
\end{situation}

\noindent
Dans la section précédente, nous avons déjà accompli une partie substantielle
du travail sur le champ $\Spacesstack'_{fp, flat, proper}$ des espaces
algébriques de présentation finie, plats et propres. Pour utiliser ces
résultats, nous considérons le foncteur d'oubli
\begin{equation}
\label{equation-over-proper-spaces}
\Polarizedstack \longrightarrow
\Spacesstack'_{fp, flat, proper},\quad
(X \to S, \mathcal{L}) \longmapsto (X \to S)
\end{equation}
Ce foncteur sera un outil utile dans la suite.
Observons que, si $(X \to S)$ appartient à l'image essentielle de
(\ref{equation-over-proper-spaces}), alors $X$ et $S$ sont des schémas.

\begin{lemma}
\label{lemma-polarized-fibred-in-groupoids}
La catégorie $\Polarizedstack$ est fibrée en groupoïdes sur
$\Spacesstack'_{fp, flat, proper}$.
La catégorie $\Polarizedstack$ est fibrée en groupoïdes sur $\Sch_{fppf}$.
\end{lemma}

\begin{proof}
Nous vérifions les conditions (1) et (2) de
Catégories, définition \ref{categories-definition-fibred-groupoids}.

\medskip\noindent
Condition (1). Soit $(X \to S, \mathcal{L})$ un objet de
$\Polarizedstack$ et soit $(X' \to S') \to (X \to S)$ un morphisme de
$\Spacesstack'_{fp, flat, proper}$. Soit alors $\mathcal{L}'$ l'image inverse
de $\mathcal{L}$ sur $X'$.
Observons que $X, S, S'$ sont des schémas ; par conséquent, $X'$ est aussi un
schéma (en tant que produit fibré de schémas). Alors $\mathcal{L}'$ est ample
sur $X'/S'$ d'après
Morphismes, lemme \ref{morphisms-lemma-ample-base-change}.
Nous obtenons ainsi un morphisme
$(X' \to S', \mathcal{L}') \to (X \to S, \mathcal{L})$
au-dessus de $(X' \to S') \to (X \to S)$.

\medskip\noindent
Condition (2). Considérons des morphismes
$(f, g, \varphi) : (X' \to S', \mathcal{L}') \to (X \to S, \mathcal{L})$ et
$(a, b, \psi) : (Y \to T, \mathcal{N}) \to (X \to S, \mathcal{L})$
de $\Polarizedstack$. Étant donné un morphisme
$(k, h) : (Y \to T) \to (X' \to S')$ de
$\Spacesstack'_{fp, flat, proper}$ tel que
$(f, g) \circ (k, h) = (a, b)$, nous devons montrer qu'il existe un unique
morphisme
$(k, h, \chi) : (Y \to T, \mathcal{N}) \to (X' \to S', \mathcal{L}')$
de $\Polarizedstack$ tel que
$(f, g, \varphi) \circ (k, h, \chi) = (a, b, \psi)$.
Nous pouvons simplement prendre
$$
\chi = \psi \circ (k^*\varphi)^{-1}
$$
Cela prouve la condition (2). Une composée de foncteurs définissant des
catégories fibrées définit une catégorie fibrée, voir
Catégories, lemme \ref{categories-lemma-fibred-over-fibred}.
Nous voyons ainsi que $\Polarizedstack$ est fibrée en groupoïdes sur
$\Sch_{fppf}$ (à strictement parler, nous devrions vérifier que les catégories
fibres sont des groupoïdes et appliquer
Catégories, lemme \ref{categories-lemma-fibred-groupoids}).
\end{proof}

\begin{lemma}
\label{lemma-polarized-stack}
La catégorie $\Polarizedstack$ est un champ en groupoïdes sur
$\Spacesstack'_{fp, flat, proper}$ (muni de la topologie induite, voir
Champs, définition \ref{stacks-definition-topology-inherited}).
La catégorie $\Polarizedstack$ est un champ en groupoïdes sur $\Sch_{fppf}$.
\end{lemma}

\begin{proof}
Nous démontrons que $\Polarizedstack$ est un champ en groupoïdes sur
$\Spacesstack'_{fp, flat, proper}$ en vérifiant les conditions (1), (2) et (3)
de Champs, définition \ref{stacks-definition-stack-in-groupoids}.
Nous avons déjà vu (1) au lemme
\ref{lemma-polarized-fibred-in-groupoids}.

\medskip\noindent
Un recouvrement de $\Spacesstack'_{fp, flat, proper}$ s'obtient comme suit :
soit $X \to S$ un objet de $\Spacesstack'_{fp, flat, proper}$.
Supposons que $\{S_i \to S\}_{i \in I}$ soit un recouvrement de
$\Sch_{fppf}$. Posons $X_i = S_i \times_S X$. Alors
$\{(X_i \to S_i) \to (X \to S)\}_{i \in I}$ est un recouvrement de
$\Spacesstack'_{fp, flat, proper}$, et tout recouvrement de
$\Spacesstack'_{fp, flat, proper}$ est isomorphe à l'un de ceux-ci.
Posons $S_{ij} = S_i \times_S S_j$ et $X_{ij} = S_{ij} \times_S X$, de sorte
que $(X_{ij} \to S_{ij}) =
(X_i \to S_i) \times_{(X \to S)} (X_j \to S_j)$.
Supposons ensuite que $\mathcal{L}, \mathcal{N}$ soient des faisceaux
inversibles amples sur $X/S$, de sorte que
$(X \to S, \mathcal{L})$ et $(X \to S, \mathcal{N})$ soient deux objets de
$\Polarizedstack$ au-dessus de l'objet $(X \to S)$.
Pour vérifier la descente des morphismes, supposons donnés des morphismes
$(\text{id}, \text{id}, \varphi_i)$ de
$(X_i \to S_i, \mathcal{L}|_{X_i})$ vers
$(X_i \to S_i, \mathcal{N}|_{X_i})$
dont les changements de base en morphismes de
$(X_{ij} \to S_{ij}, \mathcal{L}|_{X_{ij}})$ vers
$(X_{ij} \to S_{ij}, \mathcal{N}|_{X_{ij}})$
coïncident. Alors
$\varphi_i : \mathcal{L}|_{X_i} \to \mathcal{N}|_{X_i}$
sont des isomorphismes de modules inversibles sur $X_i$ tels que les
restrictions de $\varphi_i$ et $\varphi_j$ à $X_{ij}$ soient les mêmes
isomorphismes.
Par descente des faisceaux quasi-cohérents
(Descente sur les espaces, proposition
\ref{spaces-descent-proposition-fpqc-descent-quasi-coherent}),
nous obtenons un unique isomorphisme $\varphi : \mathcal{L} \to \mathcal{N}$
dont la restriction à $X_i$ redonne $\varphi_i$.

\medskip\noindent
La descente des objets se démontre exactement de la même manière.
Supposons en effet que
$\{(X_i \to S_i) \to (X \to S)\}_{i \in I}$ soit un recouvrement de
$\Spacesstack'_{fp, flat, proper}$ comme ci-dessus.
Supposons donnés des objets $(X_i \to S_i, \mathcal{L}_i)$ de
$\Polarizedstack$ au-dessus de $(X_i \to S_i)$ et une donnée de descente
$$
(\text{id}, \text{id}, \varphi_{ij}) :
(X_{ij} \to S_{ij}, \mathcal{L}_i|_{X_{ij}})
\to
(X_{ij} \to S_{ij}, \mathcal{L}_j|_{X_{ij}})
$$
satisfaisant à la condition de cocycle évidente sur
$(X_{ijk} \to S_{ijk})$ pour tout triplet d'indices.
Alors, par descente des faisceaux quasi-cohérents
(Descente sur les espaces, proposition
\ref{spaces-descent-proposition-fpqc-descent-quasi-coherent}),
nous obtenons un unique $\mathcal{O}_X$-module inversible $\mathcal{L}$ et des
isomorphismes $\mathcal{L}|_{X_i} \to \mathcal{L}_i$ redonnant la donnée de
descente $\varphi_{ij}$.
Pour montrer que $(X \to S, \mathcal{L})$ est un objet de
$\Polarizedstack$, il nous reste à prouver que $\mathcal{L}$ est ample.
Cela résulte de Descente sur les espaces, lemme
\ref{spaces-descent-lemma-descending-property-ample}.

\medskip\noindent
Puisque nous avons déjà vu que $\Spacesstack'_{fp, flat, proper}$ est un champ
en groupoïdes sur $\Sch_{fppf}$ (lemme \ref{lemma-spaces-stack}), il s'ensuit
maintenant formellement que $\Polarizedstack$ est un champ en groupoïdes sur
$\Sch_{fppf}$. Voir Champs, lemme \ref{stacks-lemma-stack-over-stack}.
\end{proof}

\noindent
Vérification de cohérence : le champ $\Polarizedstack$ joue le même rôle parmi
les espaces algébriques.

\begin{lemma}
\label{lemma-extend-polarized-to-spaces}
Soit $T$ un espace algébrique sur $\mathbf{Z}$. Soit $\mathcal{S}_T$
le champ algébrique correspondant (Champs algébriques, sections
\ref{algebraic-section-split},
\ref{algebraic-section-representable-by-algebraic-spaces} et
\ref{algebraic-section-stacks-spaces}).
Nous avons une équivalence de catégories
$$
\left\{
\begin{matrix}
(X \to T, \mathcal{L})\text{ où }X \to T\text{ est un morphisme}\\
\text{d'espaces algébriques, propre, plat et de}\\
\text{présentation finie, et }\mathcal{L}\text{ ample sur }X/T
\end{matrix}
\right\}
\longrightarrow
\Mor_{\textit{Cat}/\Sch_{fppf}}(\mathcal{S}_T, \Polarizedstack)
$$
\end{lemma}

\begin{proof}
Omis. Indications : raisonner exactement comme dans la démonstration du
lemme \ref{lemma-extend-spaces-to-spaces} et utiliser
Descente sur les espaces, proposition
\ref{spaces-descent-proposition-fpqc-descent-quasi-coherent}, pour faire
descendre le faisceau inversible dans la construction du foncteur
quasi-inverse. La propriété d'amplitude relative descend d'après
Descente sur les espaces, lemme
\ref{spaces-descent-lemma-descending-property-ample}.
\end{proof}

\begin{remark}
\label{remark-polarized-base-change}
Soit $B$ un espace algébrique sur $\Spec(\mathbf{Z})$.
Soit $B\textit{-Polarized}$ la catégorie formée des triplets
$(X \to S, \mathcal{L}, h : S \to B)$, où
$(X \to S, \mathcal{L})$ est un objet de $\Polarizedstack$ et
$h : S \to B$ est un morphisme.
Un morphisme
$(X' \to S', \mathcal{L}', h') \to (X \to S, \mathcal{L}, h)$ dans
$B\textit{-Polarized}$ est un morphisme $(f, g, \varphi)$ dans
$\Polarizedstack$ tel que $h \circ g = h'$.
Dans cette situation, le diagramme
$$
\xymatrix{
B\textit{-Polarized} \ar[r] \ar[d] & \Polarizedstack \ar[d] \\
(\Sch/B)_{fppf} \ar[r] & \Sch_{fppf}
}
$$
est un carré $2$-produit fibré. Cette remarque triviale sera parfois utile
pour déduire des résultats du cas absolu $\Polarizedstack$ au cas des familles
sur un espace algébrique de base donné.
\end{remark}

\begin{lemma}
\label{lemma-polarized-to-spaces-algebraic}
Le foncteur (\ref{equation-over-proper-spaces}) définit un $1$-morphisme
$$
\Polarizedstack \to \Spacesstack'_{fp, flat, proper}
$$
de champs en groupoïdes sur $\Sch_{fppf}$ qui est algébrique au sens de
Critères de représentabilité, définition
\ref{criteria-definition-algebraic}.
\end{lemma}

\begin{proof}
D'après les lemmes \ref{lemma-spaces-stack} et
\ref{lemma-polarized-stack}, l'énoncé a un sens. Pour le démontrer,
choisissons un schéma $S$ et un objet
$\xi = (X \to S)$ de $\Spacesstack'_{fp, flat, proper}$ au-dessus de $S$.
Nous devons montrer que
$$
\mathcal{X} = (\Sch/S)_{fppf} \times_{\xi, \Spacesstack'_{fp, flat, proper}}
\Polarizedstack
$$
est un champ algébrique sur $S$. Observons qu'un objet de $\mathcal{X}$ est
donné par un couple $(T/S, \mathcal{L})$, où $T$ est un schéma sur $S$ et où
$\mathcal{L}$ est un $\mathcal{O}_{X_T}$-module inversible qui est ample sur
$X_T/T$. Les morphismes sont définis de manière évidente.
En particulier, nous voyons immédiatement que nous avons une inclusion
$$
\mathcal{X} \subset \Picardstack_{X/S}
$$
de catégories sur $(\Sch/S)_{fppf}$, induisant l'égalité sur les ensembles de
morphismes. Puisque $\Picardstack_{X/S}$ est un champ algébrique d'après la
proposition \ref{proposition-pic}, il suffit de montrer que l'inclusion
ci-dessus est représentable par des immersions ouvertes. C'est exactement le
contenu de Descente sur les espaces, lemme
\ref{spaces-descent-lemma-ample-in-neighbourhood}.
\end{proof}

\begin{lemma}
\label{lemma-polarized-diagonal}
La diagonale
$$
\Delta : \Polarizedstack \longrightarrow
\Polarizedstack \times \Polarizedstack
$$
est représentable par des espaces algébriques.
\end{lemma}

\begin{proof}
C'est une conséquence formelle des lemmes
\ref{lemma-polarized-to-spaces-algebraic} et
\ref{lemma-spaces-diagonal}.
Voir Critères de représentabilité, lemme
\ref{criteria-lemma-diagonals-and-algebraic-morphisms}.
\end{proof}

\begin{lemma}
\label{lemma-polarized-limits}
Le champ en groupoïdes $\Polarizedstack$ commute aux limites
(Axiomes d'Artin, définition \ref{artin-definition-limit-preserving}).
\end{lemma}

\begin{proof}
Soit $I$ un ensemble filtrant et soit $(A_i, \varphi_{ii'})$ un système
d'anneaux indexé par $I$. Posons $S = \Spec(A)$ et $S_i = \Spec(A_i)$.
Nous devons montrer que, sur les catégories fibres, nous avons
$$
\Polarizedstack_S = \colim \Polarizedstack_{S_i}
$$
Nous savons que la catégorie des schémas de présentation finie sur $S$ est la
limite inductive des catégories des schémas de présentation finie sur les
$S_i$, voir
Limites, lemme \ref{limits-lemma-descend-finite-presentation}.
De plus, étant donné $X_i \to S_i$ de présentation finie, de limite
$X \to S$, la catégorie des $\mathcal{O}_X$-modules inversibles
$\mathcal{L}$ est la limite inductive des catégories des
$\mathcal{O}_{X_i}$-modules inversibles $\mathcal{L}_i$, voir
Limites, lemmes \ref{limits-lemma-descend-modules-finite-presentation} et
\ref{limits-lemma-descend-invertible-modules}.
Si $X \to S$ est propre et plat, alors, pour $i$ assez grand, le morphisme
$X_i \to S_i$ est également propre et plat, voir
Limites, lemmes \ref{limits-lemma-eventually-proper} et
\ref{limits-lemma-descend-flat-finite-presentation}.
Enfin, si $\mathcal{L}$ est ample sur $X$, alors $\mathcal{L}_i$ est ample sur
$X_i$ pour $i$ assez grand, voir
Limites, lemme \ref{limits-lemma-limit-ample}.
La réunion de ces faits achève la démonstration.
\end{proof}

\begin{lemma}
\label{lemma-polarized-RS-star}
Dans la situation \ref{situation-coherent}. Soit
$$
\xymatrix{
T \ar[r] \ar[d] & T' \ar[d] \\
S \ar[r] & S'
}
$$
une somme amalgamée dans la catégorie des schémas, où
$T \to T'$ est un épaississement et $T \to S$ est affine, voir
Compléments sur les morphismes, lemme
\ref{more-morphisms-lemma-pushout-along-thickening}.
Alors le foncteur sur les catégories fibres
$$
\Polarizedstack_{S'}
\longrightarrow
\Polarizedstack_S \times_{\Polarizedstack_T} \Polarizedstack_{T'}
$$
est une équivalence.
\end{lemma}

\begin{proof}
D'après Compléments sur les morphismes, lemme
\ref{more-morphisms-lemma-equivalence-categories-schemes-over-pushout-flat},
il existe une équivalence
$$
\textit{flat-lfp}_{S'}
\longrightarrow
\textit{flat-lfp}_S \times_{\textit{flat-lfp}_T} \textit{flat-lfp}_{T'}
$$
où $\textit{flat-lfp}_S$ désigne la catégorie des schémas plats et localement
de présentation finie sur $S$.
Soit $X'/S'$ dans le membre de gauche correspondant au triplet
$(X/S, Y'/T', \varphi)$ du membre de droite.
Posons $Y = T \times_{T'} Y'$, qui est isomorphe à $T \times_S X$ par
$\varphi$. Alors Compléments sur les morphismes, lemme
\ref{more-morphisms-lemma-scheme-over-pushout-flat-modules}, montre que nous
avons une équivalence
$$
\textit{QCoh-flat}_{X'/S'}
\longrightarrow
\textit{QCoh-flat}_{X/S}
\times_{\textit{QCoh-flat}_{Y/T}} \textit{QCoh-flat}_{Y'/T'}
$$
où $\textit{QCoh-flat}_{X/S}$ désigne la catégorie des
$\mathcal{O}_X$-modules quasi-cohérents plats sur $S$.
Puisque $X \to S$, $Y \to T$, $X' \to S'$ et $Y' \to T'$ sont plats, cela
s'applique en particulier aux modules inversibles et donne une équivalence de
catégories
$$
\textit{Pic}(X')
\longrightarrow
\textit{Pic}(X) \times_{\textit{Pic}(Y)} \textit{Pic}(Y')
$$
où $\textit{Pic}(X)$ désigne la catégorie des $\mathcal{O}_X$-modules
inversibles. Il y a ici un petit point à vérifier : il faut montrer que, si un
objet $\mathcal{F}'$ de $\textit{QCoh-flat}_{X'/S'}$ a pour images inverses des
modules inversibles sur $X$ et $Y'$, alors $\mathcal{F}'$ est un
$\mathcal{O}_{X'}$-module inversible.
Il résulte du lemme cité que $\mathcal{F}'$ est un
$\mathcal{O}_{X'}$-module de présentation finie.
D'après Compléments sur les morphismes, lemme
\ref{more-morphisms-lemma-flat-and-free-at-point-fibre}, il suffit de vérifier
que la restriction de $\mathcal{F}'$ aux fibres de $X' \to S'$ est inversible.
Mais les fibres de $X' \to S'$ sont les mêmes que celles de $X \to S$ ; ces
restrictions sont donc inversibles.

\medskip\noindent
Compte tenu de ce qui précède, nous obtenons une équivalence de catégories si
nous supprimons l'hypothèse (pour la catégorie des objets au-dessus de $S$) que
$X \to S$ soit propre et l'hypothèse que $\mathcal{L}$ soit ample.
Il est maintenant clair que, si $X' \to S'$ est propre, alors
$X \to S$ et $Y' \to T'$ sont propres (Morphismes, lemme
\ref{morphisms-lemma-base-change-proper}).
Réciproquement, si $X \to S$ et $Y' \to T'$ sont propres, alors
$X' \to S'$ est propre d'après
Compléments sur les morphismes, lemme
\ref{more-morphisms-lemma-thicken-property-morphisms-cartesian}.
De même, si $\mathcal{L}'$ est ample sur $X'/S'$, alors
$\mathcal{L}'|_X$ est ample sur $X/S$ et
$\mathcal{L}'|_{Y'}$ est ample sur $Y'/T'$
(Morphismes, lemme \ref{morphisms-lemma-ample-base-change}).
Enfin, si $\mathcal{L}'|_X$ est ample sur $X/S$ et si
$\mathcal{L}'|_{Y'}$ est ample sur $Y'/T'$, alors
$\mathcal{L}'$ est ample sur $X'/S'$ d'après
Compléments sur les morphismes, lemme
\ref{more-morphisms-lemma-thicken-property-relatively-ample}.
\end{proof}

\begin{lemma}
\label{lemma-polarized-tangent-space}
Soit $k$ un corps et soit $x = (X \to \Spec(k), \mathcal{L})$
un objet de $\mathcal{X} = \Polarizedstack$ au-dessus de $\Spec(k)$.
\begin{enumerate}
\item Si $k$ est de type fini sur $\mathbf{Z}$, alors les espaces vectoriels
$T\mathcal{F}_{\mathcal{X}, k, x}$ et
$\text{Inf}(\mathcal{F}_{\mathcal{X}, k, x})$
(voir Axiomes d'Artin, section \ref{artin-section-tangent-spaces})
sont de dimension finie, et
\item en général, les espaces vectoriels $T_x(k)$ et $\text{Inf}_x(k)$
(voir Axiomes d'Artin, section \ref{artin-section-inf})
sont de dimension finie.
\end{enumerate}
\end{lemma}

\begin{proof}
La discussion d'Axiomes d'Artin, section \ref{artin-section-tangent-spaces},
ne s'applique qu'aux corps de type fini sur le schéma de base
$\Spec(\mathbf{Z})$. Notre champ satisfait à (RS*) d'après le
lemme \ref{lemma-polarized-RS-star}, et nous pouvons appliquer
Axiomes d'Artin, lemme \ref{artin-lemma-properties-lift-RS-star},
pour obtenir les espaces vectoriels $T_x(k)$ et $\text{Inf}_x(k)$ mentionnés en
(2). De plus, dans le cas de type fini, ces espaces coïncident avec ceux
mentionnés dans la partie (1), d'après Axiomes d'Artin, remarque
\ref{artin-remark-compare-deformation-spaces}.
Cela étant réglé, nous pouvons commencer la démonstration.

\medskip\noindent
Une démonstration consiste à employer un argument analogue à celui de la
démonstration du lemme \ref{lemma-spaces-tangent-space} ; cela nous obligerait
à développer une théorie des déformations pour les couples formés d'un schéma
et d'un module quasi-cohérent.
Une autre démonstration utiliserait le résultat du
lemme \ref{lemma-spaces-tangent-space}, l'algébricité de
$\Polarizedstack \to \Spacesstack'_{fp, flat, proper}$ et un calcul de l'espace
de déformations d'un module inversible. Nous allons toutefois plutôt traduire
la question en une question de déformation d'algèbres graduées sur $k$, puis en
déduire le résultat.

\medskip\noindent
Soit $\mathcal{C}_k$ la catégorie des $k$-algèbres locales artiniennes $A$ de
corps résiduel $k$. Notre objet $x$ de $\mathcal{X}$ au-dessus de $k$ donne une
catégorie de prédéformations
$p : \mathcal{F} \to \mathcal{C}_k$, voir
Axiomes d'Artin, section \ref{artin-section-predeformation-categories}.
Ainsi, $\mathcal{F}(A)$ est la catégorie des triplets
$(X_A, \mathcal{L}_A, \alpha)$, où $(X_A, \mathcal{L}_A)$ est un objet de
$\Polarizedstack$ au-dessus de $A$ et où $\alpha$ est un isomorphisme
$(X_A, \mathcal{L}_A) \times_{\Spec(A)} \Spec(k) \cong (X, \mathcal{L})$.
D'autre part, soit $q : \mathcal{G} \to \mathcal{C}_k$ la catégorie cofibrée
en groupoïdes définie dans
Problèmes de déformations, exemple
\ref{examples-defos-example-graded-algebras}.
Choisissons $d_0 \gg 0$ (nous verrons plus bas quelle grandeur suffit).
Soit $P$ la $k$-algèbre graduée
$$
P = k \oplus \bigoplus\nolimits_{d \geq d_0} H^0(X, \mathcal{L}^{\otimes d})
$$
Alors $y = (k, P)$ est un objet de $\mathcal{G}(k)$.
Soit $\mathcal{G}_y$ la catégorie de prédéformations de
Théorie formelle des déformations, remarque
\ref{formal-defos-remark-localize-cofibered-groupoid}.
Étant donné $(X_A, \mathcal{F}_A, \alpha)$ comme ci-dessus, posons
$$
Q = A \oplus \bigoplus\nolimits_{d \geq d_0} H^0(X_A, \mathcal{L}_A^{\otimes d})
$$
L'isomorphisme $\alpha$ induit une application $\beta : Q \to P$.
Par la théorie des déformations des schémas projectifs
(Compléments sur les morphismes, lemme
\ref{more-morphisms-lemma-deform-projective}), nous obtenons un
$1$-morphisme
$$
\mathcal{F} \longrightarrow \mathcal{G}_y,\quad
(X_A, \mathcal{F}_A, \alpha) \longmapsto (Q, \beta : Q \to P)
$$
de catégories cofibrées en groupoïdes sur $\mathcal{C}_k$.
En fait, ce foncteur est une équivalence, avec pour quasi-inverse celui donné
par $Q \mapsto \underline{\text{Proj}}_A(Q)$.
En effet, le schéma $X_A = \underline{\text{Proj}}_A(Q)$ est plat sur $A$
d'après Diviseurs, lemme \ref{divisors-lemma-relative-proj-flat}.
Posons $\mathcal{L}_A = \mathcal{O}_{X_A}(1)$ ; celui-ci est plat sur $A$
d'après le même lemme. Nous tirons de $\beta$ un isomorphisme
$(X_A, \mathcal{L}_A) \times_{\Spec(A)} \Spec(k) = (X, \mathcal{L})$.
Nous pouvons alors déduire toutes les propriétés voulues du couple
$(X_A, \mathcal{L}_A)$ des propriétés correspondantes de
$(X, \mathcal{L})$ à l'aide des techniques de
Compléments sur les morphismes, sections
\ref{more-morphisms-section-morphisms-thickenings} et
\ref{more-morphisms-section-deform}.
Nous omettons certains détails.

\medskip\noindent
En conclusion, nous voyons que
$T\mathcal{F} = T\mathcal{G}_y = T_y\mathcal{G}$ et
$\text{Inf}(\mathcal{F}) = \text{Inf}_y(\mathcal{G})$.
Ces espaces vectoriels sont de dimension finie d'après
Problèmes de déformations, lemme
\ref{examples-defos-lemma-graded-algebras-TI}, ce qui achève la démonstration.
\end{proof}

\begin{lemma}[Effectivité formelle forte pour les schémas polarisés]
\label{lemma-polarized-strong-effectiveness}
\begin{slogan}
Le théorème d'algébrisation de Grothendieck reste valable dans le cadre non
noethérien si l'on suppose la platitude et la présentation finie.
\end{slogan}
Soit $(R_n)$ un système projectif d'anneaux à applications de transition
surjectives dont les noyaux sont localement nilpotents. Posons
$R = \lim R_n$.
Posons $S_n = \Spec(R_n)$ et $S = \Spec(R)$. Considérons un diagramme
commutatif
$$
\xymatrix{
X_1 \ar[r]_{i_1} \ar[d] & X_2 \ar[r]_{i_2} \ar[d] & X_3 \ar[r] \ar[d] &
\ldots \\
S_1 \ar[r] & S_2 \ar[r] & S_3 \ar[r] & \ldots
}
$$
de schémas dont les carrés sont cartésiens. Supposons donnés
$(\mathcal{L}_n, \varphi_n)$, où chaque $\mathcal{L}_n$ est un faisceau
inversible sur $X_n$ et où
$\varphi_n : i_n^*\mathcal{L}_{n + 1} \to \mathcal{L}_n$ est un isomorphisme.
Si
\begin{enumerate}
\item $X_n \to S_n$ est propre, plat et de présentation finie, et
\item $\mathcal{L}_1$ est ample sur $X_1$,
\end{enumerate}
alors il existe un morphisme de schémas $X \to S$ propre, plat et de
présentation finie, un $\mathcal{O}_X$-module inversible ample
$\mathcal{L}$, ainsi que des isomorphismes
$X_n \cong X \times_S S_n$ et
$\mathcal{L}_n \cong \mathcal{L}|_{X_n}$ compatibles aux morphismes $i_n$ et
$\varphi_n$.
\end{lemma}

\begin{proof}
Choisissons $d_0$ pour $X_1 \to S_1$ et $\mathcal{L}_1$ comme dans
Compléments sur les morphismes, lemme
\ref{more-morphisms-lemma-deform-projective}.
Pour tout $n \geq 1$, posons
$$
A_n = R_n \oplus
\bigoplus\nolimits_{d \geq d_0} H^0(X_n, \mathcal{L}_n^{\otimes d})
$$
D'après le lemme, chaque $A_n$ est une $R_n$-algèbre graduée de présentation
finie dont les composantes homogènes $(A_n)_d$ sont des $R_n$-modules
projectifs finis, telle que $X_n = \text{Proj}(A_n)$ et
$\mathcal{L}_n = \mathcal{O}_{\text{Proj}(A_n)}(1)$.
Le lemme garantit aussi que les applications
$$
A_1 \leftarrow A_2 \leftarrow A_3 \leftarrow \ldots
$$
induisent des isomorphismes $A_n = A_m \otimes_{R_m} R_n$ pour $n \leq m$.
Nous posons
$$
B = \bigoplus\nolimits_{d \geq 0} B_d
\quad\text{avec}\quad
B_d = \lim_n (A_n)_d
$$
D'après Compléments d'algèbre, lemme
\ref{more-algebra-lemma-lim-finite-projective-gives-finite-projective},
$B_d$ est un $R$-module projectif fini et $B \otimes_R R_n = A_n$.
Ainsi, le schéma
$$
X = \text{Proj}(B)
\quad\text{et}\quad
\mathcal{L} = \mathcal{O}_X(1)
$$
est plat sur $S$, et $\mathcal{L}$ est un $\mathcal{O}_X$-module
quasi-cohérent plat sur $S$, voir
Diviseurs, lemme \ref{divisors-lemma-relative-proj-flat}.
Comme la formation de Proj commute au changement de base
(Constructions, lemme \ref{constructions-lemma-base-change-map-proj}),
nous obtenons des isomorphismes canoniques
$$
X \times_S S_n = X_n
\quad\text{et}\quad
\mathcal{L}|_{X_n} \cong \mathcal{L}_n
$$
compatibles aux applications de transition du système.
Nous pouvons donc considérer $X_1 \subset X$ comme un sous-schéma fermé.
Nous montrerons plus bas que $B$ est de présentation finie sur $R$.
D'après Diviseurs, lemmes \ref{divisors-lemma-relative-proj-proper} et
\ref{divisors-lemma-relative-proj-finite-presentation},
cela implique que $X \to S$ est de présentation finie et propre, et que
$\mathcal{L} = \mathcal{O}_X(1)$ est de présentation finie en tant que
$\mathcal{O}_X$-module.
Puisque la restriction de $\mathcal{L}$ au changement de base
$X_1 \to S_1$ est inversible, il résulte de
Compléments sur les morphismes, lemme
\ref{more-morphisms-lemma-finite-free-open}, que $\mathcal{L}$ est inversible
sur un voisinage ouvert de $X_1$ dans $X$.
Comme $X \to S$ est fermé et comme $\Ker(R \to R_1)$ est contenu dans le
radical de Jacobson
(Compléments d'algèbre, lemme \ref{more-algebra-lemma-limit-henselian}),
nous voyons que tout voisinage ouvert de $X_1$ dans $X$ est égal à $X$.
Ainsi, $\mathcal{L}$ est inversible. Enfin, l'ensemble des points de $S$ où
$\mathcal{L}$ est ample sur la fibre est ouvert dans $S$
(Compléments sur les morphismes, lemme
\ref{more-morphisms-lemma-ample-in-neighbourhood})
et contient $S_1$ ; il est donc égal à $S$. Ainsi, $X \to S$ et
$\mathcal{L}$ possèdent toutes les propriétés requises dans l'énoncé du lemme.

\medskip\noindent
Démontrons l'assertion ci-dessus.
Choisissons une présentation
$A_1 = R_1[X_1, \ldots, X_s]/(F_1, \ldots, F_t)$ où les $X_i$ sont des
variables de degrés $d_i$ et les $F_j$ des polynômes homogènes en les $X_i$ de
degré $e_j$.
Nous pouvons alors choisir une application
$$
\Psi : R[X_1, \ldots, X_s] \longrightarrow B
$$
relevant l'application $R_1[X_1, \ldots, X_s] \to A_1$.
Puisque chaque $B_d$ est projectif fini sur $R$, nous concluons du lemme de
Nakayama (Algèbre, lemme \ref{algebra-lemma-NAK}, en utilisant de nouveau le
fait que $\Ker(R \to R_1)$ est contenu dans le radical de Jacobson de $R$) que
$\Psi$ est surjective. Comme $- \otimes_R R_1$ est exact à droite, nous pouvons
trouver $G_1, \ldots, G_t \in \Ker(\Psi)$ dont les images sont
$F_1, \ldots, F_t$ dans $R_1[X_1, \ldots, X_s]$.
Observons que $\Ker(\Psi)_d$ est un $R$-module projectif fini pour tout
$d \geq 0$, en tant que noyau de la surjection
$R[X_1, \ldots, X_s]_d \to B_d$ de $R$-modules projectifs finis.
Nous concluons encore une fois du lemme de Nakayama que $\Ker(\Psi)$ est
engendré par $G_1, \ldots, G_t$.
\end{proof}

\begin{lemma}
\label{lemma-polarized-existence}
Considérons le champ $\Polarizedstack$ sur le schéma de base
$\Spec(\mathbf{Z})$. Alors tout objet formel est effectif.
\end{lemma}

\begin{proof}
Pour les définitions des notions du lemme, voir
Axiomes d'Artin, section \ref{artin-section-formal-objects}.
Il ressort des définitions que le lemme résulte immédiatement du lemme plus
général \ref{lemma-polarized-strong-effectiveness}.
\end{proof}

\begin{lemma}
\label{lemma-polarized-defo-thy}
Le champ en groupoïdes $\Polarizedstack$ satisfait à l'ouverture de la
versalité sur $\Spec(\mathbf{Z})$.
De même, après changement de base (remarque
\ref{remark-polarized-base-change}), l'ouverture de la versalité vaut sur tout
schéma de base noethérien $S$.
\end{lemma}

\begin{proof}
Cela résulte de
Axiomes d'Artin, lemme \ref{artin-lemma-SGE-implies-openness-versality},
et des lemmes \ref{lemma-polarized-diagonal},
\ref{lemma-polarized-RS-star},
\ref{lemma-polarized-limits} et
\ref{lemma-polarized-strong-effectiveness}.
Pour la démonstration « usuelle » de ce fait, voir la discussion de la remarque
qui suit cette démonstration.
\end{proof}

\begin{remark}
\label{remark-polarized-defo-thy}
Le lemme \ref{lemma-polarized-defo-thy} peut également se démontrer à l'aide
d'une théorie des obstructions comme dans
Axiomes d'Artin, lemme \ref{artin-lemma-get-openness-obstruction-theory}
(comme dans la seconde démonstration du
lemme \ref{lemma-coherent-defo-thy}).
Pour cela, il faut généraliser la théorie des déformations et des obstructions
développée dans
Cotangent, section \ref{cotangent-section-deformations-ringed-topoi}, au cas des
couples formés d'espaces algébriques et de modules quasi-cohérents.
Une autre possibilité consiste à utiliser le fait que le $1$-morphisme
$\Polarizedstack \to \Spacesstack'_{fp, flat, proper}$ est algébrique
(lemme \ref{lemma-polarized-to-spaces-algebraic}) et le fait que nous
connaissons l'ouverture de la versalité pour le but
(lemme \ref{lemma-spaces-defo-thy} et remarque
\ref{remark-spaces-defo-thy}).
\end{remark}

\begin{theorem}[Algébricité du champ des schémas polarisés]
\label{theorem-polarized-algebraic}
Le champ $\Polarizedstack$ (situation \ref{situation-polarized})
est algébrique. En fait, pour tout espace algébrique $B$, le champ
$B\textit{-Polarized}$ (remarque \ref{remark-polarized-base-change})
est algébrique.
\end{theorem}

\begin{proof}
Le cas absolu résulte de
Axiomes d'Artin, lemme \ref{artin-lemma-diagonal-representable},
et des lemmes \ref{lemma-polarized-diagonal},
\ref{lemma-polarized-RS-star},
\ref{lemma-polarized-limits},
\ref{lemma-polarized-existence} et
\ref{lemma-polarized-defo-thy}.
Le cas sur $B$ en résulte, compte tenu de la description de
$B\textit{-Polarized}$ comme $2$-produit fibré dans la remarque
\ref{remark-polarized-base-change} et du fait que les champs algébriques
admettent les $2$-produits fibrés, voir
Champs algébriques, lemme \ref{algebraic-lemma-2-fibre-product}.
\end{proof}

\section{Le champ des courbes}
\label{section-curves}

\noindent
Dans cette section, nous prouvons que le champ des courbes est algébrique.
Pour une discussion plus approfondie des espaces de modules de courbes, nous
renvoyons le lecteur à Modules de courbes, section
\ref{moduli-curves-section-introduction}.

\medskip\noindent
Dans le Projet Stacks, une courbe est une variété de dimension $1$.
Cependant, lorsque nous parlons de familles de courbes, nous autorisons souvent
les fibres à être réductibles et/ou non réduites. Dans cette section, le champ
des courbes « paramétrera les schémas propres de dimension $\leq 1$ ».
Il se trouve toutefois que, pour obtenir la bonne notion de famille, il nous
faut autoriser l'espace total de notre famille à être un espace algébrique.
Cela conduit à la définition suivante.

\begin{situation}
\label{situation-curves}
Nous définissons une catégorie $\Curvesstack$ comme suit :
\begin{enumerate}
\item Les objets sont les {\it familles de courbes}. Plus précisément, un
objet est un morphisme $f : X \to S$ où la base $S$ est un schéma,
l'{\it espace total} $X$ est un espace algébrique, et où
$f$ est plat, propre, de présentation finie et de dimension relative
$\leq 1$ (Morphismes d'espaces, définition
\ref{spaces-morphisms-definition-relative-dimension}).
\item Un morphisme $(X' \to S') \to (X \to S)$ entre objets est donné par un
couple $(f, g)$, où $f : X' \to X$ est un morphisme d'espaces algébriques et
$g : S' \to S$ est un morphisme de schémas, qui s'inscrivent dans un diagramme
commutatif
$$
\xymatrix{
X' \ar[d] \ar[r]_f & X \ar[d] \\
S' \ar[r]^g & S
}
$$
induisant un isomorphisme $X' \to S' \times_S X$ ; autrement dit, le
diagramme est cartésien.
\end{enumerate}
Le foncteur d'oubli
$$
p : \Curvesstack \longrightarrow \Sch_{fppf},\quad
(X \to S) \longmapsto S
$$
est celui par lequel nous considérons $\Curvesstack$ comme une catégorie sur
$\Sch_{fppf}$
(voir la section \ref{section-conventions} pour les notations).
\end{situation}

\noindent
Il résulte de Espaces sur les corps, lemme
\ref{spaces-over-fields-lemma-codim-1-point-in-schematic-locus},
et, plus généralement, de
Compléments sur les morphismes d'espaces, lemme
\ref{spaces-more-morphisms-lemma-projective-over-complete},
que, si $S$ est le spectre d'un corps, d'un anneau local artinien ou d'un
anneau local noethérien complet, alors, pour toute famille de courbes
$X \to S$, l'espace total $X$ est un schéma.
En revanche, il existe des familles de courbes sur $\mathbf{A}^1_k$ dont
l'espace total n'est pas un schéma, voir
Exemples, section \ref{examples-section-family-of-curves}.

\medskip\noindent
Il est clair que
\begin{equation}
\label{equation-curves-over-proper-spaces}
\Curvesstack \subset \Spacesstack'_{fp, flat, proper}
\end{equation}
et qu'un objet $X \to S$ de $\Spacesstack'_{fp, flat, proper}$ appartient à
$\Curvesstack$ si et seulement si $X \to S$ est de dimension relative
$\leq 1$. Nous utiliserons ce fait pour vérifier les axiomes d'Artin pour
$\Curvesstack$.

\begin{lemma}
\label{lemma-curves-fibred-in-groupoids}
La catégorie $\Curvesstack$ est fibrée en groupoïdes sur $\Sch_{fppf}$.
\end{lemma}

\begin{proof}
En utilisant le plongement (\ref{equation-curves-over-proper-spaces}),
la description de son image et le fait correspondant pour
$\Spacesstack'_{fp, flat, proper}$
(lemme \ref{lemma-spaces-fibred-in-groupoids}), nous nous ramenons à l'énoncé
suivant : étant donné un morphisme
$$
\xymatrix{
X' \ar[r] \ar[d] & X \ar[d] \\
S' \ar[r] & S
}
$$
dans $\Spacesstack'_{fp, flat, proper}$ (rappelons que cela implique en
particulier que le diagramme est cartésien), si $X \to S$ est de dimension
relative $\leq 1$, alors $X' \to S'$ est de dimension relative $\leq 1$.
Cela résulte de Morphismes d'espaces, lemme
\ref{spaces-morphisms-lemma-dimension-fibre-after-base-change}.
\end{proof}

\begin{lemma}
\label{lemma-curves-stack}
La catégorie $\Curvesstack$ est un champ en groupoïdes sur $\Sch_{fppf}$.
\end{lemma}

\begin{proof}
En utilisant le plongement (\ref{equation-curves-over-proper-spaces}),
la description de son image et le fait correspondant pour
$\Spacesstack'_{fp, flat, proper}$ (lemme \ref{lemma-spaces-stack}),
nous nous ramenons à l'énoncé suivant : étant donné un objet $X \to S$ de
$\Spacesstack'_{fp, flat, proper}$ et un recouvrement fppf
$\{S_i \to S\}_{i \in I}$, les assertions suivantes sont équivalentes :
\begin{enumerate}
\item $X \to S$ est de dimension relative $\leq 1$, et
\item pour tout $i$, le changement de base $X_i \to S_i$ est de dimension
relative $\leq 1$.
\end{enumerate}
Cela résulte de Morphismes d'espaces, lemme
\ref{spaces-morphisms-lemma-dimension-fibre-after-base-change}.
\end{proof}

\begin{lemma}
\label{lemma-curves-diagonal}
La diagonale
$$
\Delta : \Curvesstack \longrightarrow \Curvesstack \times \Curvesstack
$$
est représentable par des espaces algébriques.
\end{lemma}

\begin{proof}
Cela résulte immédiatement du plongement pleinement fidèle
(\ref{equation-curves-over-proper-spaces}) et du fait correspondant pour
$\Spacesstack'_{fp, flat, proper}$
(lemme \ref{lemma-spaces-diagonal}).
\end{proof}

\begin{remark}
\label{remark-curves-base-change}
Soit $B$ un espace algébrique sur $\Spec(\mathbf{Z})$.
Soit $B\text{-}\Curvesstack$ la catégorie formée des couples
$(X \to S, h : S \to B)$, où $X \to S$ est un objet de $\Curvesstack$ et
$h : S \to B$ est un morphisme.
Un morphisme $(X' \to S', h') \to (X \to S, h)$ dans
$B\text{-}\Curvesstack$ est un morphisme $(f, g)$ dans $\Curvesstack$ tel que
$h \circ g = h'$.
Dans cette situation, le diagramme
$$
\xymatrix{
B\text{-}\Curvesstack \ar[r] \ar[d] & \Curvesstack \ar[d] \\
(\Sch/B)_{fppf} \ar[r] & \Sch_{fppf}
}
$$
est un carré $2$-produit fibré. Cette remarque triviale sera parfois utile
pour déduire des résultats du cas absolu $\Curvesstack$ au cas des familles de
courbes sur un espace algébrique de base donné.
\end{remark}

\begin{lemma}
\label{lemma-curves-limits}
Le champ $\Curvesstack \to \Sch_{fppf}$ commute aux limites
(Axiomes d'Artin, définition \ref{artin-definition-limit-preserving}).
\end{lemma}

\begin{proof}
En utilisant le plongement (\ref{equation-curves-over-proper-spaces}),
la description de son image et le fait correspondant pour
$\Spacesstack'_{fp, flat, proper}$ (lemme \ref{lemma-spaces-limits}),
nous nous ramenons à l'énoncé suivant :
soit $T = \lim T_i$ la limite d'un système projectif filtrant de schémas
affines. Soit $i \in I$ et soit $X_i \to T_i$ un objet de
$\Spacesstack'_{fp, flat, proper}$ au-dessus de $T_i$.
Supposons que $T \times_{T_i} X_i \to T$ soit de dimension relative
$\leq 1$. Alors, pour un certain $i' \geq i$, le morphisme
$T_{i'} \times_{T_i} X_i \to T_i$ est de dimension relative $\leq 1$.
Cela résulte de Limites d'espaces, lemme
\ref{spaces-limits-lemma-eventually-relative-dimension}.
\end{proof}

\begin{lemma}
\label{lemma-curves-RS-star}
Soit
$$
\xymatrix{
T \ar[r] \ar[d] & T' \ar[d] \\
S \ar[r] & S'
}
$$
une somme amalgamée dans la catégorie des schémas, où
$T \to T'$ est un épaississement et $T \to S$ est affine, voir
Compléments sur les morphismes, lemme
\ref{more-morphisms-lemma-pushout-along-thickening}.
Alors le foncteur sur les catégories fibres
$$
\Curvesstack_{S'}
\longrightarrow
\Curvesstack_S
\times_{\Curvesstack_T}
\Curvesstack_{T'}
$$
est une équivalence.
\end{lemma}

\begin{proof}
En utilisant le plongement (\ref{equation-curves-over-proper-spaces}),
la description de son image et le fait correspondant pour
$\Spacesstack'_{fp, flat, proper}$ (lemme \ref{lemma-spaces-RS-star}),
nous nous ramenons à l'énoncé suivant : étant donné un morphisme
$X' \to S'$ d'un espace algébrique vers $S'$ qui est de présentation finie,
plat et propre, alors $X' \to S'$ est de dimension relative $\leq 1$ si et
seulement si $S \times_{S'} X' \to S$ et
$T' \times_{S'} X' \to T'$ sont de dimension relative $\leq 1$.
Une implication résulte du fait que la propriété d'être de dimension relative
$\leq 1$ est préservée par changement de base
(Morphismes d'espaces, lemme
\ref{spaces-morphisms-lemma-dimension-fibre-after-base-change}).
L'autre résulte du fait que la propriété d'être de dimension relative
$\leq 1$ se vérifie sur les fibres et du fait que les fibres de $X' \to S'$
(aux points du schéma $S'$) sont les mêmes que celles de
$S \times_{S'} X' \to S$, puisque $S \to S'$ est un épaississement d'après
Compléments sur les morphismes, lemme
\ref{more-morphisms-lemma-pushout-along-thickening}.
\end{proof}

\begin{lemma}
\label{lemma-curves-tangent-space}
Soit $k$ un corps et soit $x = (X \to \Spec(k))$ un objet de
$\mathcal{X} = \Curvesstack$ au-dessus de $\Spec(k)$.
\begin{enumerate}
\item Si $k$ est de type fini sur $\mathbf{Z}$, alors les espaces vectoriels
$T\mathcal{F}_{\mathcal{X}, k, x}$ et
$\text{Inf}(\mathcal{F}_{\mathcal{X}, k, x})$
(voir Axiomes d'Artin, section \ref{artin-section-tangent-spaces})
sont de dimension finie, et
\item en général, les espaces vectoriels $T_x(k)$ et $\text{Inf}_x(k)$
(voir Axiomes d'Artin, section \ref{artin-section-inf})
sont de dimension finie.
\end{enumerate}
\end{lemma}

\begin{proof}
Cela résulte immédiatement du plongement pleinement fidèle
(\ref{equation-curves-over-proper-spaces}) et du fait correspondant pour
$\Spacesstack'_{fp, flat, proper}$
(lemme \ref{lemma-spaces-tangent-space}).
\end{proof}

\begin{lemma}
\label{lemma-curves-existence}
Considérons le champ $\Curvesstack$ sur le schéma de base
$\Spec(\mathbf{Z})$. Alors tout objet formel est effectif.
\end{lemma}

\begin{proof}
Pour les définitions des notions du lemme, voir
Axiomes d'Artin, section \ref{artin-section-formal-objects}.
Soit $(A, \mathfrak m, \kappa)$ un anneau local noethérien complet.
Soit $(X_n \to \Spec(A/\mathfrak m^n))$ un objet formel de
$\Curvesstack$ au-dessus de $A$.
D'après Compléments sur les morphismes d'espaces, lemme
\ref{spaces-more-morphisms-lemma-formal-algebraic-space-proper-reldim-1},
il existe un morphisme projectif $X \to \Spec(A)$ et un système compatible
d'isomorphismes
$X \times_{\Spec(A)} \Spec(A/\mathfrak m^n) \cong X_n$.
D'après Compléments sur les morphismes, lemme
\ref{more-morphisms-lemma-check-flatness-on-infinitesimal-nbhds},
$X \to \Spec(A)$ est plat. D'après Compléments sur les morphismes, lemme
\ref{more-morphisms-lemma-dimension-fibres-proper-flat},
$X \to \Spec(A)$ est de dimension relative $\leq 1$.
Cela prouve le lemme.
\end{proof}

\begin{lemma}
\label{lemma-curves-defo-thy}
Le champ en groupoïdes $\mathcal{X} = \Curvesstack$ satisfait à l'ouverture
de la versalité sur $\Spec(\mathbf{Z})$.
De même, après changement de base (remarque \ref{remark-curves-base-change}),
l'ouverture de la versalité vaut sur tout schéma de base noethérien $S$.
\end{lemma}

\begin{proof}
Cela résulte immédiatement du plongement pleinement fidèle
(\ref{equation-curves-over-proper-spaces}) et du fait correspondant pour
$\Spacesstack'_{fp, flat, proper}$
(lemme \ref{lemma-spaces-defo-thy}).
\end{proof}

\begin{theorem}[Algébricité du champ des courbes]
\label{theorem-curves-algebraic}
\begin{reference}
Voir \cite[Proposition 3.3, page 8]{dJHS} et
\cite[Appendix B by Jack Hall, Theorem B.1]{Smyth}.
\end{reference}
Le champ $\Curvesstack$ (situation \ref{situation-curves})
est algébrique. En fait, pour tout espace algébrique $B$, le champ
$B\text{-}\Curvesstack$ (remarque \ref{remark-curves-base-change})
est algébrique.
\end{theorem}

\begin{proof}
Le cas absolu résulte de
Axiomes d'Artin, lemme \ref{artin-lemma-diagonal-representable},
et des lemmes \ref{lemma-curves-diagonal},
\ref{lemma-curves-RS-star},
\ref{lemma-curves-limits},
\ref{lemma-curves-existence} et
\ref{lemma-curves-defo-thy}.
Le cas sur $B$ en résulte, compte tenu de la description de
$B\text{-}\Curvesstack$ comme $2$-produit fibré dans la remarque
\ref{remark-curves-base-change} et du fait que les champs algébriques admettent
les $2$-produits fibrés, voir
Champs algébriques, lemme \ref{algebraic-lemma-2-fibre-product}.
\end{proof}

\begin{lemma}
\label{lemma-curves-open-and-closed-in-spaces}
Le $1$-morphisme (\ref{equation-curves-over-proper-spaces})
$$
\Curvesstack \longrightarrow \Spacesstack'_{fp, flat, proper}
$$
est représentable par des immersions ouvertes et fermées.
\end{lemma}

\begin{proof}
Puisque (\ref{equation-curves-over-proper-spaces}) est un plongement pleinement
fidèle de catégories, il suffit de montrer ce qui suit : étant donné un objet
$X \to S$ de $\Spacesstack'_{fp, flat, proper}$, il existe un sous-schéma
ouvert et fermé $U \subset S$ tel qu'un morphisme $S' \to S$ se factorise par
$U$ si et seulement si le changement de base $X' \to S'$ de $X \to S$ est de
dimension relative $\leq 1$.
Cela résulte immédiatement de
Compléments sur les morphismes d'espaces, lemme
\ref{spaces-more-morphisms-lemma-dimension-fibres-proper-flat}.
\end{proof}

\begin{remark}
\label{remark-alternative-approach-curves}
Considérons le $2$-produit fibré
$$
\xymatrix{
\Curvesstack \times_{\Spacesstack'_{fp, flat, proper}}
\Polarizedstack \ar[r] \ar[d] &
\Polarizedstack \ar[d] \\
\Curvesstack \ar[r] &
\Spacesstack'_{fp, flat, proper}
}
$$
Ce produit fibré paramètre les courbes polarisées, c'est-à-dire les familles de
courbes munies d'un faisceau inversible relativement ample.
Il se trouve que la flèche verticale de gauche
$$
\textit{PolarizedCurves} \longrightarrow \Curvesstack
$$
est algébrique, lisse et surjective. En effet, ce $1$-morphisme est algébrique
(en tant que changement de base de la flèche du
lemme \ref{lemma-polarized-to-spaces-algebraic}), tout point appartient à son
image, et il n'y a pas d'obstruction à déformer les faisceaux inversibles sur
les courbes (voir la démonstration du lemme \ref{lemma-curves-existence}).
Cela donne une autre approche de l'algébricité de $\Curvesstack$.
En effet, d'après le lemme \ref{lemma-curves-open-and-closed-in-spaces},
$\textit{PolarizedCurves}$ est un sous-champ ouvert et fermé du champ
algébrique $\Polarizedstack$, et tout champ en groupoïdes qui est le but d'un
morphisme algébrique lisse provenant d'un champ algébrique est un champ
algébrique.
\end{remark}

\section{Espaces de modules de complexes sur un morphisme propre}
\label{section-moduli-complexes}

\noindent
Le titre et le contenu de cette section sont tirés de
\cite{lieblich-complexes}. Soit $S$ un schéma et soit
$f : X \to B$ un morphisme propre, plat et de présentation finie d'espaces
algébriques. Nous allons démontrer qu'il existe un champ algébrique
$$
\Complexesstack_{X/B}
$$
qui paramètre les « familles » d'objets de $D^b_{\textit{Coh}}$ des fibres
dont les auto-Ext négatifs s'annulent. Plus précisément, une famille est donnée
par un objet relativement parfait de la catégorie dérivée de l'espace total ;
cette notion quelque peu technique est étudiée dans
Compléments sur les morphismes d'espaces, section
\ref{spaces-more-morphisms-section-relatively-perfect}.

\medskip\noindent
Déjà lorsque $X$ est un espace algébrique propre sur un corps $k$, nous
obtenons un champ algébrique très intéressant. En effet, il existe un
plongement
$$
\Cohstack_{X/k} \longrightarrow \Complexesstack_{X/k}
$$
car, pour tout $\mathcal{O}$-module $\mathcal{F}$ (sur tout topos annelé),
nous avons $\Ext^i_\mathcal{O}(\mathcal{F}, \mathcal{F}) = 0$ pour $i < 0$.
Bien que cela montre assurément que notre champ est non vide, la véritable
motivation pour étudier $\Complexesstack_{X/k}$ est qu'il existe souvent des
objets de la catégorie dérivée $D^b_{\textit{Coh}}(\mathcal{O}_X)$ dont les
auto-Ext négatifs s'annulent et dont les faisceaux de cohomologie sont non nuls
en plus d'un degré.
Par exemple, $X$ pourrait être équivalent au sens dérivé à un autre espace
algébrique propre $Y$ sur $k$, c'est-à-dire que nous aurions une équivalence
$k$-linéaire
$$
F : D^b_{\textit{Coh}}(\mathcal{O}_Y)
\longrightarrow
D^b_{\textit{Coh}}(\mathcal{O}_X)
$$
Il existe des cas où cela se produit sans que $F$ soit donné par un
isomorphisme entre $X$ et $Y$ ; par exemple, dans le cas d'une variété
abélienne et de sa duale. Dans cette situation, $F$ induit un isomorphisme de
champs algébriques
$$
\Complexesstack_{Y/k}
\longrightarrow
\Complexesstack_{X/k}
$$
(insérer ici une référence future) et, en particulier, le champ des faisceaux
cohérents sur $Y$ s'envoie dans le champ des complexes sur $X$. En renversant
ce point de vue, si nous comprenons assez bien la géométrie de
$\Complexesstack_{X/k}$, nous pouvons tenter de l'utiliser pour étudier tous
les $Y$ qui lui sont équivalents au sens dérivé.

\begin{lemma}
\label{lemma-complexes-open-neg-exts-vanishing}
Soit $S$ un schéma.
Soit $f : X \to Y$ un morphisme d'espaces algébriques sur $S$.
Supposons que $f$ soit propre, plat et de présentation finie.
Soient $K, E \in D(\mathcal{O}_X)$. Supposons que $K$ soit pseudo-cohérent et
que $E$ soit $Y$-parfait (Compléments sur les morphismes d'espaces, définition
\ref{spaces-more-morphisms-definition-relatively-perfect}).
Pour un corps $k$ et un morphisme $y : \Spec(k) \to Y$, notons $K_y$, $E_y$
les images inverses sur la fibre $X_y$.
\begin{enumerate}
\item Il existe un ouvert $W \subset Y$ caractérisé par la propriété
$$
y \in |W|
\Leftrightarrow
\Ext^i_{\mathcal{O}_{X_y}}(K_y, E_y) = 0
\text{ pour }i < 0.
$$
\item Pour tout morphisme $V \to Y$ se factorisant par $W$, nous avons
$$
\Ext^i_{\mathcal{O}_{X_V}}(K_V, E_V) = 0
\quad\text{pour}\quad i < 0
$$
où $X_V$ est le changement de base de $X$, et $K_V$ et $E_V$ sont les images
inverses dérivées de $K$ et $E$ sur $X_V$.
\item Le foncteur $V \mapsto \Hom_{\mathcal{O}_{X_V}}(K_V, E_V)$
est un faisceau sur $(\textit{Spaces}/W)_{fppf}$ représentable par un espace
algébrique affine et de présentation finie sur $W$.
\end{enumerate}
\end{lemma}

\begin{proof}
Pour tout morphisme $V \to Y$, le complexe $K_V$ est pseudo-cohérent
(Cohomologie sur les sites, lemme
\ref{sites-cohomology-lemma-pseudo-coherent-pullback})
et $E_V$ est $V$-parfait (Compléments sur les morphismes d'espaces, lemme
\ref{spaces-more-morphisms-lemma-base-change-relatively-perfect}).
Observons aussi qu'étant donnés $y : \Spec(k) \to Y$, une extension de corps
$k'/k$ et le morphisme induit $y' : \Spec(k') \to Y$, nous avons
$$
\Ext^i_{\mathcal{O}_{X_{y'}}}(K_{y'}, E_{y'}) =
\Ext^i_{\mathcal{O}_{X_y}}(K_y, E_y) \otimes_k k'
$$
d'après Catégories dérivées de schémas, lemme
\ref{perfect-lemma-affine-morphism-and-hom-out-of-perfect}.
Ainsi, l'annulation en (1) est bien une propriété du point induit
$y \in |Y|$. Nous utiliserons ces deux observations sans plus les mentionner
dans la démonstration.

\medskip\noindent
Supposons d'abord que $Y$ soit un schéma affine. Nous pouvons alors appliquer
Compléments sur les morphismes d'espaces, lemme
\ref{spaces-more-morphisms-lemma-compute-ext-rel-perfect},
et trouver un $L \in D(\mathcal{O}_Y)$ pseudo-cohérent qui « calcule
universellement » $Rf_*R\SheafHom(K, E)$ au sens décrit dans ce lemme.
En développant les définitions, nous obtenons, pour un point $y \in Y$,
l'égalité
$$
\Ext^i_{\kappa(y)}(L \otimes_{\mathcal{O}_Y}^\mathbf{L} \kappa(y),
\kappa(y)) = \Ext^i_{\mathcal{O}_{X_y}}(K_y, E_y)
$$
Nous en concluons que
$$
H^i(L \otimes_{\mathcal{O}_Y}^\mathbf{L} \kappa(y)) = 0
\text{ pour } i > 0 \Leftrightarrow
\Ext^i_{\mathcal{O}_{X_y}}(K_y, E_y) = 0 \text{ pour }i < 0.
$$
D'après Catégories dérivées de schémas, lemme
\ref{perfect-lemma-jump-loci}, l'ensemble $W$ des $y \in Y$ où cela se produit
définit un ouvert de $Y$.
Cet ouvert $W$ satisfait alors à la condition de (1) pour tous les morphismes
provenant de spectres de corps, par l'« universalité » de $L$.

\medskip\noindent
Revenons au cas où $Y$ est un espace algébrique général.
Choisissons un recouvrement \'etale $\{V_i \to Y\}$ par des schémas affines
$V_i$. Nous voyons alors que l'image inverse du sous-ensemble $W \subset |Y|$
est le sous-ensemble correspondant $W_i \subset |V_i|$ pour
$X_{V_i}$, $K_{V_i}$, $E_{V_i}$.
D'après le paragraphe précédent, $W_i$ est ouvert, donc $W$ est ouvert.
Cela prouve (1) en général. En outre, les parties (2) et (3) sont entièrement
formulées en termes de la catégorie $\textit{Spaces}/W$ et des restrictions
$X_W$, $K_W$, $E_W$. Nous sommes donc ramenés au cas $W = Y$.

\medskip\noindent
Supposons $W = Y$. Nous affirmons que, pour tout espace algébrique $V$ sur
$Y$, le complexe $Rf_{V, *}R\SheafHom(K_V, E_V)$ a ses faisceaux de
cohomologie nuls en degrés $< 0$. Cela prouvera (2), car
$$
\Ext^i_{\mathcal{O}_{X_V}}(K_V, E_V) =
H^i(X_V, R\SheafHom(K_V, E_V)) = 
H^i(V, Rf_{V, *}R\SheafHom(K_V, E_V))
$$
d'après Cohomologie sur les sites, lemmes
\ref{sites-cohomology-lemma-section-RHom-over-U} et
\ref{sites-cohomology-lemma-Leray-unbounded}, et l'annulation des faisceaux de
cohomologie entraîne que le groupe de cohomologie $H^i$ est nul pour $i < 0$
d'après Catégories dérivées, lemme
\ref{derived-lemma-negative-vanishing}.

\medskip\noindent
Pour prouver l'affirmation, nous pouvons travailler localement pour la
topologie \'etale sur $V$. En particulier, nous pouvons supposer que $Y$ est
affine et que $W = Y$.
Soit $L \in D(\mathcal{O}_Y)$ comme dans le deuxième paragraphe de la
démonstration. Pour un espace algébrique $V$ sur $Y$, notons $L_V$ l'image
inverse dérivée de $L$ sur $V$. (Une propriété importante que nous utiliserons
est que $L$ « fonctionne » pour tous les espaces algébriques $V$ sur $Y$, et
pas seulement pour les $V$ affines.)
Comme $W = Y$, nous avons $H^i(L) = 0$ pour $i > 0$
(utiliser Compléments d'algèbre, lemme
\ref{more-algebra-lemma-lift-pseudo-coherent-from-residue-field},
pour passer des fibres aux germes). Donc $H^i(L_V) = 0$ pour $i > 0$.
La propriété qui définit $L$ est que
$$
Rf_{V, *}R\SheafHom(K_V, E_V) = R\SheafHom(L_V, \mathcal{O}_V)
$$
Puisque $L_V$ est concentré en degrés $\leq 0$, nous concluons que
$R\SheafHom(L_V, \mathcal{O}_V)$ est concentré en degrés $\geq 0$, ce qui
prouve l'affirmation. Cela achève la démonstration de (2).

\medskip\noindent
Supposons $W = Y$, sans faire d'hypothèse sur l'espace algébrique $Y$.
Comme nous avons (2), il résulte de
Espaces simpliciaux, lemme
\ref{spaces-simplicial-lemma-fppf-neg-ext-zero-hom}, que le foncteur $F$ donné
par $F(V) = \Hom_{\mathcal{O}_{X_V}}(K_V, E_V)$ est un
faisceau\footnote{Pour vérifier la propriété de faisceau relativement à un
recouvrement $\{V_i \to V\}_{i \in I}$, considérons d'abord
l'hyperrecouvrement fppf de {\v C}ech $a : V_\bullet \to V$, où
$V_n = \coprod_{i_0 \ldots i_n} V_{i_0} \times_V \ldots \times_V V_{i_n}$,
puis posons $U_\bullet = V_\bullet \times_{a, V} X_V$. Alors
$U_\bullet \to X_V$ est un hyperrecouvrement fppf auquel nous pouvons
appliquer Espaces simpliciaux, lemme
\ref{spaces-simplicial-lemma-fppf-neg-ext-zero-hom}.}
sur $(\textit{Spaces}/Y)_{fppf}$. Ainsi, pour prouver que $F$ est un espace
algébrique et que $F \to Y$ est affine et de présentation finie, nous pouvons
travailler localement pour la topologie \'etale sur $Y$ ; voir
Bootstrap, lemme \ref{bootstrap-lemma-locally-algebraic-space-finite-type}, et
Morphismes d'espaces, lemmes \ref{spaces-morphisms-lemma-affine-local} et
\ref{spaces-morphisms-lemma-finite-presentation-local}.
Nous concluons qu'il suffit de prouver que $F$ est un espace algébrique affine
de présentation finie sur $Y$ lorsque $Y$ est un schéma affine. Dans ce cas,
nous revenons à notre complexe pseudo-cohérent $L \in D(\mathcal{O}_Y)$.
Comme $H^i(L) = 0$ pour $i > 0$, nous pouvons représenter $L$ par un complexe
de la forme
$$
\ldots \to \mathcal{O}_Y^{\oplus m_1} \to \mathcal{O}_Y^{\oplus m_0}
\to 0 \to \ldots
$$
dont le dernier terme est en degré $0$, voir Compléments d'algèbre, lemme
\ref{more-algebra-lemma-pseudo-coherent}. En combinant les deux formules
affichées plus haut dans la démonstration, nous trouvons que
$$
F(V) =
\Ker(
\Hom_V(\mathcal{O}_V^{\oplus m_0}, \mathcal{O}_V)
\to 
\Hom_V(\mathcal{O}_V^{\oplus m_1}, \mathcal{O}_V)
)
$$
Autrement dit, il existe un diagramme produit fibré
$$
\xymatrix{
F \ar[d] \ar[r] & Y \ar[d]^0 \\
\mathbf{A}_Y^{m_0} \ar[r] & \mathbf{A}_Y^{m_1}
}
$$
qui prouve ce que nous voulions.
\end{proof}

\begin{lemma}
\label{lemma-complexes}
Soit $S$ un schéma. Soit $f : X \to Y$ un morphisme d'espaces algébriques
sur $S$. Supposons que $f$ soit propre, plat et de présentation finie.
Soit $E \in D(\mathcal{O}_X)$.
Supposons que
\begin{enumerate}
\item $E$ soit $S$-parfait (Compléments sur les morphismes d'espaces,
définition \ref{spaces-more-morphisms-definition-relatively-perfect}), et
\item pour tout point $s \in S$, nous ayons
$$
\Ext^i_{\mathcal{O}_{X_s}}(E_s, E_s) = 0
\quad\text{pour}\quad i < 0
$$
où $E_s$ est l'image inverse sur la fibre $X_s$.
\end{enumerate}
Alors
\begin{enumerate}
\item[(a)] (1) et (2) sont préservées par tout changement de base $V \to Y$,
\item[(b)] $\Ext^i_{\mathcal{O}_{X_V}}(E_V, E_V) = 0$ pour $i < 0$ et tout
$V$ sur $Y$,
\item[(c)] $V \mapsto \Hom_{\mathcal{O}_{X_V}}(E_V, E_V)$ est représentable
par un espace algébrique affine et de présentation finie sur $Y$.
\end{enumerate}
Ici, $X_V$ est le changement de base de $X$ et $E_V$ l'image inverse dérivée
de $E$ sur $X_V$.
\end{lemma}

\begin{proof}
Conséquence immédiate du lemme
\ref{lemma-complexes-open-neg-exts-vanishing}.
\end{proof}

\begin{situation}
\label{situation-complexes}
Soit $S$ un schéma. Soit $f : X \to B$ un morphisme d'espaces algébriques
sur $S$. Supposons que $f$ soit propre, plat et de présentation finie.
Nous notons $\Complexesstack_{X/B}$ la catégorie dont les objets sont les
triplets $(T, g, E)$ où
\begin{enumerate}
\item $T$ est un schéma sur $S$,
\item $g : T \to B$ est un morphisme sur $S$, et, en posant
$X_T = T \times_{g, B} X$,
\item $E$ est un objet de $D(\mathcal{O}_{X_T})$ satisfaisant aux conditions
(1) et (2) du lemme \ref{lemma-complexes}.
\end{enumerate}
Un morphisme $(T, g, E) \to (T', g', E')$ est donné par un couple
$(h, \varphi)$ où
\begin{enumerate}
\item $h : T \to T'$ est un morphisme de schémas sur $B$
(c'est-à-dire $g' \circ h = g$), et
\item $\varphi : L(h')^*E' \to E$ est un isomorphisme de
$D(\mathcal{O}_{X_T})$, où $h' : X_T \to X_{T'}$ est le changement de base de
$h$.
\end{enumerate}
\end{situation}

\noindent
Ainsi, $\Complexesstack_{X/B}$ est une catégorie et la règle
$$
p : \Complexesstack_{X/B} \longrightarrow (\Sch/S)_{fppf},
\quad
(T, g, E) \longmapsto T
$$
est un foncteur. Pour un schéma $T$ sur $S$, nous notons
$\Complexesstack_{X/B, T}$ la catégorie fibre de $p$ au-dessus de $T$.
Ces catégories fibres sont des groupoïdes.

\begin{lemma}
\label{lemma-complexes-fibred-in-groupoids}
Dans la situation \ref{situation-complexes}, le foncteur
$p : \Complexesstack_{X/B} \longrightarrow (\Sch/S)_{fppf}$
est fibré en groupoïdes.
\end{lemma}

\begin{proof}
Nous montrons que $p$ est fibré en groupoïdes en vérifiant les conditions (1)
et (2) de Catégories, définition
\ref{categories-definition-fibred-groupoids}.
Étant donnés un objet $(T', g', E')$ de $\Complexesstack_{X/B}$ et un morphisme
$h : T \to T'$ de schémas sur $S$, nous pouvons poser $g = h \circ g'$ et
$E = L(h')^*E'$, où $h' : X_T \to X_{T'}$ est le changement de base de $h$.
Il est alors clair que nous obtenons un morphisme
$(T, g, E) \to (T', g', E')$ de $\Complexesstack_{X/B}$ au-dessus de $h$.
Cela prouve (1).
Pour (2), supposons donnés des morphismes
$$
(h_1, \varphi_1) : (T_1, g_1, E_1) \to (T, g, E)
\quad\text{et}\quad
(h_2, \varphi_2) : (T_2, g_2, E_2) \to (T, g, E)
$$
de $\Complexesstack_{X/B}$ et un morphisme $h : T_1 \to T_2$ tel que
$h_2 \circ h = h_1$. Nous pouvons alors prendre pour $\varphi$ la composée
$$
L(h')^*E_2
\xrightarrow{L(h')^*\varphi_2^{-1}}
L(h')^*L(h_2)^*E = L(h_1)^*E
\xrightarrow{\varphi_1}
E_1
$$
et obtenir ainsi le morphisme
$(h, \varphi) : (T_1, g_1, E_1) \to (T_2, g_2, E_2)$
qui atteste la condition (2).
\end{proof}

\begin{lemma}
\label{lemma-complexes-diagonal}
Dans la situation \ref{situation-complexes}. Notons
$\mathcal{X} = \Complexesstack_{X/B}$. Alors
$\Delta : \mathcal{X} \to \mathcal{X} \times \mathcal{X}$ est représentable
par des espaces algébriques.
\end{lemma}

\begin{proof}
Considérons deux objets $x = (T, g, E)$ et $y = (T, g', E')$ de
$\mathcal{X}$ au-dessus d'un schéma $T$. Nous devons montrer que
$\mathit{Isom}_\mathcal{X}(x, y)$ est un espace algébrique sur $T$, voir
Champs algébriques, lemme \ref{algebraic-lemma-representable-diagonal}.
Si, pour $h : T' \to T$, les restrictions $x|_{T'}$ et $y|_{T'}$ sont
isomorphes dans la catégorie fibre $\mathcal{X}_{T'}$, alors
$g \circ h = g' \circ h$.
Il existe donc une transformation de préfaisceaux
$$
\mathit{Isom}_\mathcal{X}(x, y) \longrightarrow \text{Égaliseur}(g, g')
$$
Puisque la diagonale de $B$ est représentable (par des schémas), cet égaliseur
est un schéma. Nous pouvons donc remplacer $T$ par cet égaliseur et $E$, $E'$
par leurs images inverses. Nous pouvons ainsi supposer $g = g'$.

\medskip\noindent
Supposons $g = g'$. Après avoir remplacé $B$ par $T$ et $X$ par $X_T$, nous
arrivons au problème suivant. Étant donnés $E, E' \in D(\mathcal{O}_X)$
satisfaisant aux conditions (1), (2) du lemme \ref{lemma-complexes}, nous devons
montrer que $\mathit{Isom}(E, E')$ est un espace algébrique.
Ici, $\mathit{Isom}(E, E')$ est le foncteur
$$
(\Sch/B)^{opp} \to \textit{Sets},\quad
T \mapsto \{\varphi : E_T \to E'_T
\text{ isomorphisme dans }D(\mathcal{O}_{X_T})\}
$$
où $E_T$ et $E'_T$ sont les images inverses dérivées de $E$ et $E'$ sur
$X_T$. Soient maintenant $W \subset B$, resp.\ $W' \subset B$, les sous-espaces
ouverts de $B$ associés à $E, E'$, resp.\ à $E', E$, par le
lemme \ref{lemma-complexes-open-neg-exts-vanishing}.
Il est clair que, s'il existe un isomorphisme $E_T \to E'_T$ comme dans la
définition de $\mathit{Isom}(E, E')$, alors $T \to B$ se factorise à la fois
par $W$ et par $W'$ (parce que nous avons la condition (1) pour $E$ et $E'$, et
que nous aurons évidemment $E_t \cong E'_t$, donc aucune application non nulle
$E_t[i] \to E_t$ ou $E'_t[i] \to E_t$ sur la fibre $X_t$ pour $i > 0$).
Nous pouvons donc remplacer $B$ par l'ouvert $W \cap W'$.
Dans ce cas, le foncteur $H = \SheafHom(E, E')$
$$
(\Sch/B)^{opp} \to \textit{Sets},\quad T \mapsto
\Hom_{\mathcal{O}_{X_T}}(E_T, E'_T)
$$
est un espace algébrique affine et de présentation finie sur $B$ d'après le
lemme \ref{lemma-complexes-open-neg-exts-vanishing}.
Il en est de même de
$H' = \SheafHom(E', E)$,
$I = \SheafHom(E, E)$ et
$I' = \SheafHom(E', E')$.
Nous pouvons donc reprendre l'argument de la démonstration de la
proposition \ref{proposition-isom} pour voir que
$$
\mathit{Isom}(E, E') = (H' \times_B H) \times_{c, I \times_B I', \sigma} B
$$
pour certains morphismes $c$ et $\sigma$. Ainsi,
$\mathit{Isom}(E, E')$ est un espace algébrique.
\end{proof}

\begin{lemma}
\label{lemma-complexes-stack}
Dans la situation \ref{situation-complexes}, le foncteur
$p : \Complexesstack_{X/B} \longrightarrow (\Sch/S)_{fppf}$
est un champ en groupoïdes.
\end{lemma}

\begin{proof}
Pour prouver que $\Complexesstack_{X/B}$ est un champ en groupoïdes, nous
devons montrer que les préfaisceaux $\mathit{Isom}$ sont des faisceaux et que
les données de descente sont effectives. L'assertion sur $\mathit{Isom}$
résulte du lemme \ref{lemma-complexes-diagonal}, voir
Champs algébriques, lemme \ref{algebraic-lemma-representable-diagonal}.
Démontrons l'assertion sur les données de descente.

\medskip\noindent
Supposons que $\{a_i : T_i \to T\}$ soit un recouvrement fppf de schémas sur
$S$. Soit $(\xi_i, \varphi_{ij})$ une donnée de descente pour
$\{T_i \to T\}$ à valeurs dans $\Complexesstack_{X/B}$.
Pour chaque $i$, nous pouvons écrire $\xi_i = (T_i, g_i, E_i)$.
Notons $\text{pr}_0 : T_i \times_T T_j \to T_i$ et
$\text{pr}_1 : T_i \times_T T_j \to T_j$ les projections.
La condition
$\xi_i|_{T_i \times_T T_j} \cong \xi_j|_{T_i \times_T T_j}$ implique en
particulier que $g_i \circ \text{pr}_0 = g_j \circ \text{pr}_1$.
Il existe donc un unique morphisme $g : T \to B$ tel que
$g_i = g \circ a_i$, voir
Descente sur les espaces, lemme
\ref{spaces-descent-lemma-fpqc-universal-effective-epimorphisms}.
Notons $X_T = T \times_{g, B} X$. Posons
$X_i = X_{T_i} = T_i \times_{g_i, B} X = T_i \times_{a_i, T} X_T$ et
$$
X_{ij} = X_{T_i} \times_{X_T} X_{T_j} = X_i \times_{X_T} X_j
$$
avec les projections $\text{pr}_i$ et $\text{pr}_j$ vers $X_i$ et $X_j$.
Observons que l'image inverse de $(T_i, g_i, E_i)$ par
$\text{pr}_0 : T_i \times_T T_j \to T_i$ est donnée par
$(T_i \times_T T_j, g_i \circ \text{pr}_0, L\text{pr}_i^*E_i)$.
Une donnée de descente pour $\{T_i \to T\}$ dans $\Complexesstack_{X/B}$ est
donc donnée par les objets $(T_i, g \circ a_i, E_i)$ et, pour chaque couple
$i, j$, par un isomorphisme dans $D\mathcal{O}_{X_{ij}})$
$$
\varphi_{ij} :
L\text{pr}_i^*E_i \longrightarrow L\text{pr}_j^*E_j
$$
satisfaisant à la condition de cocycle sur l'image inverse de $X$ sur
$T_i \times_T T_j \times_T T_k$.
En utilisant l'annulation des Ext négatifs fournie par (b) du
lemme \ref{lemma-complexes}, nous pouvons appliquer
Espaces simpliciaux, lemme
\ref{spaces-simplicial-lemma-fppf-glue-neg-ext-zero}, pour obtenir la
descente\footnote{Pour le vérifier, considérons d'abord l'hyperrecouvrement
fppf de {\v C}ech $a : T_\bullet \to T$, où
$T_n = \coprod_{i_0 \ldots i_n} T_{i_0} \times_T \ldots \times_T T_{i_n}$,
puis posons $U_\bullet = T_\bullet \times_{a, T} X_T$. Alors
$U_\bullet \to X_T$ est un hyperrecouvrement fppf auquel nous pouvons
appliquer Espaces simpliciaux, lemme
\ref{spaces-simplicial-lemma-fppf-glue-neg-ext-zero}.}
pour ces complexes. Autrement dit, nous trouvons qu'il existe un objet $E$ de
$D_\QCoh(\mathcal{O}_{X_T})$ dont la restriction à $X_{T_i}$ est $E_i$,
compatiblement aux $\varphi_{ij}$. Rappelons qu'être $T$-parfait signifie être
pseudo-cohérent et avoir une dimension de Tor localement finie sur
$f^{-1}\mathcal{O}_T$.
Ainsi, $E$ est $T$-parfait par application de
Compléments sur les morphismes d'espaces, lemmes
\ref{spaces-more-morphisms-lemma-pseudo-coherent-descends-fpqc} et
\ref{spaces-more-morphisms-lemma-tor-amplitude-descends-fppf}.
Enfin, nous devons vérifier pour $E$ la condition (2) du
lemme \ref{lemma-complexes}.
Cela résulte immédiatement de la description de l'ouvert $W$ dans le
lemme \ref{lemma-complexes-open-neg-exts-vanishing} et du fait que (2) vaut
pour $E_i$ sur $X_{T_i}/T_i$.
\end{proof}

\begin{remark}
\label{remark-complexes-base-change}
Dans la situation \ref{situation-complexes}, la règle
$(T, g, E) \mapsto (T, g)$ définit un $1$-morphisme
$$
\Complexesstack_{X/B} \longrightarrow \mathcal{S}_B
$$
de champs en groupoïdes
(voir le lemme \ref{lemma-complexes-stack},
Champs algébriques, section \ref{algebraic-section-split}, et
Exemples de champs, section
\ref{examples-stacks-section-stack-associated-to-sheaf}).
Soit $B' \to B$ un morphisme d'espaces algébriques sur $S$.
Soit $\mathcal{S}_{B'} \to \mathcal{S}_B$ le $1$-morphisme associé de champs
fibrés en ensembles. Posons $X' = X \times_B B'$.
Nous obtenons un champ en groupoïdes
$\Complexesstack_{X'/B'} \to (\Sch/S)_{fppf}$
associé au changement de base $f' : X' \to B'$. Dans cette situation, le
diagramme
$$
\vcenter{
\xymatrix{
\Complexesstack_{X'/B'} \ar[r] \ar[d] &
\Complexesstack_{X/B} \ar[d] \\
\mathcal{S}_{B'} \ar[r] & \mathcal{S}_B
}
}
\quad
\begin{matrix}
\text{ou dans} \\
\text{une autre} \\
\text{notation}
\end{matrix}
\quad
\vcenter{
\xymatrix{
\Complexesstack_{X'/B'} \ar[r] \ar[d] &
\Complexesstack_{X/B} \ar[d] \\
\Sch/B' \ar[r] & \Sch/B
}
}
$$
est un carré $2$-produit fibré. Cette remarque triviale sera parfois utile
pour changer l'espace algébrique de base.
\end{remark}

\begin{lemma}
\label{lemma-complexes-limits}
Dans la situation \ref{situation-complexes}, supposons que $B \to S$ soit
localement de présentation finie. Alors
$p : \Complexesstack_{X/B} \to (\Sch/S)_{fppf}$ commute aux limites
(Axiomes d'Artin, définition \ref{artin-definition-limit-preserving}).
\end{lemma}

\begin{proof}
Notons $B(T)$ la catégorie discrète dont les objets sont les $S$-morphismes
$T \to B$. Soit $T = \lim T_i$ une limite filtrante de schémas affines sur
$S$. En associant à un objet $(T, h, E)$ de $\Complexesstack_{X/B, T}$ l'objet
$h$ de $B(T)$, nous obtenons un diagramme commutatif de catégories fibres
$$
\xymatrix{
\colim \Complexesstack_{X/B, T_i} \ar[r] \ar[d] &
\Complexesstack_{X/B, T} \ar[d] \\
\colim B(T_i) \ar[r] & B(T)
}
$$
Nous devons montrer que la flèche horizontale supérieure est une équivalence.
Puisque nous avons supposé que $B$ est localement de présentation finie sur
$S$, il résulte de
Limites d'espaces, remarque \ref{spaces-limits-remark-limit-preserving},
que la flèche horizontale inférieure est une équivalence. Cela signifie que
nous pouvons supposer que $T = \lim T_i$ est une limite filtrante de schémas
affines sur $B$. Notons $g_i : T_i \to B$ et $g : T \to B$ les morphismes
correspondants. Posons $X_i = T_i \times_{g_i, B} X$ et
$X_T = T \times_{g, B} X$.
Observons que $X_T = \colim X_i$.
D'après Compléments sur les morphismes d'espaces, lemme
\ref{spaces-more-morphisms-lemma-descend-relatively-perfect},
la catégorie des objets $T$-parfaits de $D(\mathcal{O}_{X_T})$ est la limite
inductive des catégories des objets $T_i$-parfaits de
$D(\mathcal{O}_{X_{T_i}})$.
Il nous reste donc seulement à démontrer qu'étant donné un objet $T_i$-parfait
$E_i$ de $D(\mathcal{O}_{X_{T_i}})$ tel que l'image inverse dérivée $E$ de
$E_i$ sur $X_T$ satisfasse à la condition (2) du lemme
\ref{lemma-complexes}, alors, quitte à augmenter $i$, $E_i$ satisfait à la
condition (2) du lemme \ref{lemma-complexes}.
Soit $W \subset |T_i|$ l'ouvert construit dans le
lemme \ref{lemma-complexes-open-neg-exts-vanishing} pour $E_i$ et $E_i$.
Par l'hypothèse sur $E$, nous trouvons que $T \to T_i$ se factorise par $T$.
Il existe donc un $i' \geq i$ tel que $T_{i'} \to T_i$ se factorise par $W$,
voir Limites, lemme
\ref{limits-lemma-limit-contained-in-constructible}
Alors $i'$ convient par construction de $W$.
\end{proof}

\begin{lemma}
\label{lemma-complexes-RS-star}
Dans la situation \ref{situation-complexes}. Soit
$$
\xymatrix{
Z \ar[r] \ar[d] & Z' \ar[d] \\
Y \ar[r] & Y'
}
$$
une somme amalgamée dans la catégorie des schémas sur $S$, où
$Z \to Z'$ est un épaississement d'ordre fini et $Z \to Y$ est affine, voir
Compléments sur les morphismes, lemme
\ref{more-morphisms-lemma-pushout-along-thickening}.
Alors le foncteur sur les catégories fibres
$$
\Complexesstack_{X/B, Y'}
\longrightarrow
\Complexesstack_{X/B, Y}
\times_{\Complexesstack_{X/B, Z}}
\Complexesstack_{X/B, Z'}
$$
est une équivalence.
\end{lemma}

\begin{proof}
Observons que l'application correspondante
$$
B(Y') \longrightarrow B(Y) \times_{B(Z)} B(Z')
$$
est une bijection, voir Sommes amalgamées d'espaces, lemme
\ref{spaces-pushouts-lemma-pushout-along-thickening-schemes}.
Ainsi, en utilisant le diagramme commutatif
$$
\xymatrix{
\Complexesstack_{X/B, Y'} \ar[r] \ar[d] &
\Complexesstack_{X/B, Y}
\times_{\Complexesstack_{X/B, Z}}
\Complexesstack_{X/B, Z'}
\ar[d] \\
B(Y') \ar[r] & B(Y) \times_{B(Z)} B(Z')
}
$$
nous voyons que nous pouvons supposer que $Y'$ est un schéma sur $B'$.
D'après la remarque \ref{remark-complexes-base-change}, nous pouvons remplacer
$B$ par $Y'$ et $X$ par $X \times_B Y'$.
Nous pouvons donc supposer $B = Y'$.

\medskip\noindent
Supposons $B = Y'$. Nous prouvons d'abord la pleine fidélité de notre foncteur.
Pour cela, soient $\xi_1, \xi_2$ deux objets de $\Complexesstack_{X/B}$
au-dessus de $Y'$. Nous devons alors montrer que
$$
\mathit{Isom}(\xi_1, \xi_2)(Y') \longrightarrow
\mathit{Isom}(\xi_1, \xi_2)(Y)
\times_{\mathit{Isom}(\xi_1, \xi_2)(Z)}
\mathit{Isom}(\xi_1, \xi_2)(Z')
$$
est bijective. Or nous savons déjà que $\mathit{Isom}(\xi_1, \xi_2)$ est un
espace algébrique sur $B = Y'$. Cette bijectivité résulte donc de
Axiomes d'Artin, lemme \ref{artin-lemma-pushout} (ou du lemme déjà mentionné de
Sommes amalgamées d'espaces,
\ref{spaces-pushouts-lemma-pushout-along-thickening-schemes}).

\medskip\noindent
Surjectivité essentielle. Soit $(E_Y, E_{Z'}, \alpha)$ un triplet, où
$E_Y \in D(\mathcal{O}_Y)$ et $E_{Z'} \in D(\mathcal{O}_{X_{Z'}})$ sont des
objets tels que $(Y, Y \to B, E_Y)$ soit un objet de
$\Complexesstack_{X/B}$ au-dessus de $Y$, que
$(Z', Z' \to B, E_{Z'})$ soit un objet de $\Complexesstack_{X/B}$ au-dessus de
$Z'$, et que
$\alpha : L(X_Z \to X_Y)^*E_Y \to L(X_Z \to X_{Z'})^*E_{Z'}$
soit un isomorphisme dans $D(\mathcal{O}_{Z'})$.
C'est-à-dire que
$$
((Y, Y \to B, E_Y), (Z', Z' \to B, E_{Z'}), \alpha)
$$
est un objet du but de la flèche de notre lemme.
Observons que le diagramme
$$
\xymatrix{
X_Z \ar[r] \ar[d] & X_{Z'} \ar[d] \\
X_Y \ar[r] & X_{Y'}
}
$$
est une somme amalgamée, avec $X_Z \to X_Y$ affine et
$X_Z \to X_{Z'}$ un épaississement
(voir Sommes amalgamées d'espaces, lemme
\ref{spaces-pushouts-lemma-equivalence-categories-spaces-pushout-flat}).
Par conséquent, d'après Sommes amalgamées d'espaces, lemme
\ref{spaces-pushouts-lemma-pushout-along-thickening-derived},
nous trouvons un objet $E_{Y'} \in D(\mathcal{O}_{X_{Y'}})$ ainsi que des
isomorphismes
$L(X_Y \to X_{Y'})^*E_{Y'} \to E_Y$ et
$L(X_{Z'} \to X_{Y'})^*E_{Y'} \to E_Z$
compatibles à $\alpha$. Il est clair que, si nous montrons que $E_{Y'}$ est
$Y'$-parfait, alors nous aurons terminé, car la propriété (2) du
lemme \ref{lemma-complexes} est une propriété portant sur les points (et $Y$
et $Y'$ ont les mêmes points).
Cela résulte de Compléments sur les morphismes d'espaces, lemme
\ref{spaces-more-morphisms-lemma-thickening-relatively-perfect}.
\end{proof}

\begin{lemma}
\label{lemma-complexes-tangent-space}
Dans la situation \ref{situation-complexes}, supposons que $S$ soit un schéma
localement noethérien et que $B \to S$ soit localement de présentation finie.
Soit $k$ un corps de type fini sur $S$ et soit
$x_0 = (\Spec(k), g_0, E_0)$ un objet de
$\mathcal{X} = \Complexesstack_{X/B}$ au-dessus de $k$.
Alors les espaces $T\mathcal{F}_{\mathcal{X}, k, x_0}$ et
$\text{Inf}(\mathcal{F}_{\mathcal{X}, k, x_0})$
(Axiomes d'Artin, section \ref{artin-section-tangent-spaces})
sont de dimension finie.
\end{lemma}

\begin{proof}
Observons que, d'après le lemme \ref{lemma-complexes-RS-star}, notre champ en
groupoïdes $\mathcal{X}$ satisfait à la propriété (RS*) définie dans
Axiomes d'Artin, section \ref{artin-section-RS-star}.
En particulier, $\mathcal{X}$ satisfait à (RS).
Par conséquent, toutes les catégories de prédéformations associées sont des
catégories de déformations
(Axiomes d'Artin, lemme \ref{artin-lemma-deformation-category}),
et l'énoncé a un sens.

\medskip\noindent
Dans ce paragraphe, nous montrons que nous pouvons nous ramener au cas
$B = \Spec(k)$. Posons $X_0 = \Spec(k) \times_{g_0, B} X$ et notons
$\mathcal{X}_0 = \Complexesstack_{X_0/k}$. Dans la remarque
\ref{remark-complexes-base-change}, nous avons vu que $\mathcal{X}_0$ est le
$2$-produit fibré de $\mathcal{X}$ et $\Spec(k)$ sur $B$ comme catégories
fibrées en groupoïdes sur $(\Sch/S)_{fppf}$. Ainsi, d'après
Axiomes d'Artin, lemme \ref{artin-lemma-fibre-product-tangent-spaces},
nous nous ramenons à démontrer que $B$, $\Spec(k)$ et $\mathcal{X}_0$ ont des
espaces tangents et des espaces d'automorphismes infinitésimaux de dimension
finie. Les espaces tangents de $B$ et de $\Spec(k)$ sont de dimension finie
d'après Axiomes d'Artin, lemme \ref{artin-lemma-finite-dimension}, et leurs
$\text{Inf}$ sont bien sûr nuls.
Il suffit donc de traiter $\mathcal{X}_0$.

\medskip\noindent
Soit $k[\epsilon]$ l'algèbre des nombres duaux sur $k$.
Soit $\Spec(k[\epsilon]) \to B$ la composée de $g_0 : \Spec(k) \to B$ et du
morphisme $\Spec(k[\epsilon]) \to \Spec(k)$ provenant de l'inclusion
$k \to k[\epsilon]$.
Posons $X_0 = \Spec(k) \times_B X$ et
$X_\epsilon = \Spec(k[\epsilon]) \times_B X$.
Observons que $X_\epsilon$ est un épaississement du premier ordre de $X_0$,
plat sur l'épaississement du premier ordre
$\Spec(k) \to \Spec(k[\epsilon])$.
Observons que $X_0$ et $X_\epsilon$ donnent lieu à des petits topos étales
canoniquement équivalents, voir
Compléments sur les morphismes d'espaces, section
\ref{spaces-more-morphisms-section-thickenings}.
D'après Compléments sur les morphismes d'espaces, lemme
\ref{spaces-more-morphisms-lemma-thickening-relatively-perfect},
$T\mathcal{F}_{\mathcal{X}_0, k, x_0}$ est l'ensemble des classes
d'isomorphisme de relèvements de $E_0$ à $X_\epsilon$ au sens de
Théorie des déformations, lemme
\ref{defos-lemma-first-order-defos-complex}.
Nous concluons que
$$
T\mathcal{F}_{\mathcal{X}_0, k, x_0} =
\Ext^1_{\mathcal{O}_{X_0}}(E_0, E_0)
$$
Nous avons utilisé ici l'identification $\epsilon k[\epsilon] \cong k$ de
$k[\epsilon]$-modules. En utilisant encore une fois
Théorie des déformations, lemme
\ref{defos-lemma-first-order-defos-complex}, nous voyons qu'il existe une
surjection
$$
\text{Inf}(\mathcal{F}_{\mathcal{X}, k, x_0})
\leftarrow
\Ext^0_{\mathcal{O}_{X_0}}(E_0, E_0)
$$
d'espaces vectoriels sur $k$. Comme $E_0$ est pseudo-cohérent, il appartient à
$D^-_{\textit{Coh}}(\mathcal{O}_{X_0})$ d'après
Catégories dérivées d'espaces, lemme
\ref{spaces-perfect-lemma-identify-pseudo-coherent-noetherian}.
Puisque $E_0$ est localement de dimension de Tor finie et que $X_0$ est
quasi-compact, nous avons
$E_0 \in D^b_{\textit{Coh}}(\mathcal{O}_{X_0})$.
Les $\Ext$ ci-dessus sont donc des espaces vectoriels de dimension finie sur
$k$ d'après Catégories dérivées d'espaces, lemme
\ref{spaces-perfect-lemma-ext-finite}.
\end{proof}

\begin{lemma}
\label{lemma-complexes-strong-effectiveness}
Dans la situation \ref{situation-complexes}, supposons que $B = S$ soit
localement noethérien. Alors l'effectivité formelle forte au sens de
Axiomes d'Artin, remarque \ref{artin-remark-strong-effectiveness},
vaut pour $p : \Complexesstack_{X/S} \to (\Sch/S)_{fppf}$.
\end{lemma}

\begin{proof}
Soit $(R_n)$ un système projectif de $S$-algèbres à applications de transition
surjectives dont les noyaux sont localement nilpotents. Posons
$R = \lim R_n$. Soit $(\xi_n)$ un système d'objets de
$\Complexesstack_{X/B}$ au-dessus des $(\Spec(R_n))$. Nous devons montrer que
$(\xi_n)$ est effectif, c'est-à-dire qu'il existe un objet $\xi$ de
$\Complexesstack_{X/B}$ au-dessus de $\Spec(R)$.

\medskip\noindent
Écrivons $X_R = \Spec(R) \times_S X$ et
$X_n = \Spec(R_n) \times_S X$.
Bien entendu, $X_n$ est le changement de base de $X_R$ par $R \to R_n$.
Puisque $S = B$, nous voyons que $\xi_n$ correspond simplement à un objet
$R_n$-parfait $E_n \in D(\mathcal{O}_{X_n})$ satisfaisant à la condition (2)
du lemme \ref{lemma-complexes}. En particulier, $E_n$ est pseudo-cohérent.
Les isomorphismes
$\xi_{n + 1}|_{\Spec(R_n)} \cong \xi_n$ correspondent à des isomorphismes
$L(X_n \to X_{n + 1})^*E_{n + 1} \to E_n$.
Par conséquent, d'après Platitude sur les espaces, théorème
\ref{spaces-flat-theorem-existence-derived}, nous trouvons un objet
pseudo-cohérent $E$ de $D(\mathcal{O}_{X_R})$ dont $E_n$ est l'image inverse
dérivée de $E$ pour tout $n$, compatiblement aux isomorphismes de transition.

\medskip\noindent
Observons que $(R, \Ker(R \to R_1))$ est une paire hensélienne, voir
Compléments d'algèbre, lemme \ref{more-algebra-lemma-limit-henselian}.
En particulier, $\Ker(R \to R_1)$ est contenu dans le radical de Jacobson de
$R$. Nous pouvons alors appliquer
Compléments sur les morphismes d'espaces, lemme
\ref{spaces-more-morphisms-lemma-henselian-relatively-perfect},
pour voir que $E$ est $R$-parfait.

\medskip\noindent
Enfin, nous devons vérifier la condition (2) du lemme \ref{lemma-complexes}.
D'après le lemme \ref{lemma-complexes-open-neg-exts-vanishing}, l'ensemble des
points $t$ de $\Spec(R)$ où les auto-Ext négatifs de $E_t$ s'annulent est un
ouvert. Comme cette condition est vraie dans $V(\Ker(R \to R_1))$ et que
$\Ker(R \to R_1)$ est contenu dans le radical de Jacobson de $R$, nous
concluons qu'elle vaut en tout point.
\end{proof}

\begin{theorem}[Algébricité des espaces de modules de complexes sur un morphisme propre]
\label{theorem-complexes-algebraic}
\begin{reference}
\cite{lieblich-complexes}
\end{reference}
Soit $S$ un schéma. Soit $f : X \to B$ un morphisme d'espaces algébriques
sur $S$. Supposons que $f$ soit propre, plat et de présentation finie.
Alors $\Complexesstack_{X/B}$ est un champ algébrique sur $S$.
\end{theorem}

\begin{proof}
Posons $\mathcal{X} = \Complexesstack_{X/B}$. Nous avons vu que $\mathcal{X}$
est un champ en groupoïdes sur $(\Sch/S)_{fppf}$ dont la diagonale est
représentable par des espaces algébriques
(lemmes \ref{lemma-complexes-stack} et \ref{lemma-complexes-diagonal}).
Il suffit donc de trouver un schéma $W$ et un morphisme surjectif et lisse
$W \to \mathcal{X}$.

\medskip\noindent
Soit $B'$ un schéma et soit $B' \to B$ un morphisme surjectif \'etale.
Posons $X' = B' \times_B X$ et notons $f' : X' \to B'$ la projection.
Alors $\mathcal{X}' = \Complexesstack_{X'/B'}$ est égal au $2$-produit fibré
de $\mathcal{X}$ et de la catégorie fibrée en ensembles associée à $B'$ sur la
catégorie fibrée en ensembles associée à $B$
(remarque \ref{remark-complexes-base-change}). D'après le contenu de
Champs algébriques, section
\ref{algebraic-section-representable-properties}, le morphisme
$\mathcal{X}' \to \mathcal{X}$ est surjectif et \'etale.
Il suffit donc de prouver le résultat pour $\mathcal{X}'$.
Autrement dit, nous pouvons supposer que $B$ est un schéma.

\medskip\noindent
Supposons que $B$ soit un schéma. Dans ce cas, nous pouvons remplacer $S$ par
$B$, voir Champs algébriques, section
\ref{algebraic-section-change-base-scheme}.
Nous pouvons donc supposer $S = B$.

\medskip\noindent
Supposons $S = B$. Choisissons un recouvrement ouvert affine
$S = \bigcup U_i$. Notons $\mathcal{X}_i$ la restriction de $\mathcal{X}$ à
$(\Sch/U_i)_{fppf}$. Si nous pouvons trouver des schémas $W_i$ sur $U_i$ et
des morphismes lisses surjectifs $W_i \to \mathcal{X}_i$, alors nous posons
$W = \coprod W_i$ et obtenons un morphisme lisse surjectif
$W \to \mathcal{X}$. Nous pouvons donc supposer que $S = B$ est affine.

\medskip\noindent
Supposons $S = B$ affine, disons $S = \Spec(\Lambda)$.
Écrivons $\Lambda = \colim \Lambda_i$ comme une limite inductive filtrante où
chaque $\Lambda_i$ est de type fini sur $\mathbf{Z}$. Pour un certain $i$,
nous pouvons trouver un morphisme d'espaces algébriques
$X_i \to \Spec(\Lambda_i)$ propre, plat et de présentation finie, dont le
changement de base à $\Lambda$ est $X$. Voir Limites d'espaces, lemmes
\ref{spaces-limits-lemma-descend-finite-presentation},
\ref{spaces-limits-lemma-descend-flat} et
\ref{spaces-limits-lemma-eventually-proper}.
Si nous montrons que $\Complexesstack_{X_i/\Spec(\Lambda_i)}$ est un champ
algébrique, alors il résulte du changement de base
(remarque \ref{remark-complexes-base-change} et
Champs algébriques, section \ref{algebraic-section-change-base-scheme})
que $\mathcal{X}$ est un champ algébrique.
Nous pouvons donc supposer que $\Lambda$ est une $\mathbf{Z}$-algèbre de type
fini.

\medskip\noindent
Supposons que $S = B = \Spec(\Lambda)$ soit affine de type fini sur
$\mathbf{Z}$. Dans ce cas, nous allons vérifier les conditions (1), (2), (3),
(4) et (5) de
Axiomes d'Artin, lemme \ref{artin-lemma-diagonal-representable}, pour conclure
que $\mathcal{X}$ est un champ algébrique.
Notons que $\Lambda$ est un G-anneau, voir
Compléments d'algèbre, proposition
\ref{more-algebra-proposition-ubiquity-G-ring}.
Ainsi, tous les anneaux locaux de $S$ sont des G-anneaux. Donc (5) vaut.
Pour vérifier (2), nous devons vérifier les axiomes [-1], [0], [1], [2] et [3]
de Axiomes d'Artin, section \ref{artin-section-axioms}.
Nous omettons la vérification de [-1], et les axiomes
[0], [1], [2], [3] correspondent respectivement aux
lemmes \ref{lemma-complexes-stack},
\ref{lemma-complexes-limits},
\ref{lemma-complexes-RS-star},
\ref{lemma-complexes-tangent-space}.
La condition (3) résulte du lemme
\ref{lemma-complexes-strong-effectiveness}.
La condition (1) est le lemme \ref{lemma-complexes-diagonal}.

\medskip\noindent
Il reste à montrer la condition (4), qui est l'ouverture de la versalité.
Pour cela, nous allons utiliser
Axiomes d'Artin, lemme \ref{artin-lemma-SGE-implies-openness-versality}.
Nous avons déjà vu que la diagonale de $\mathcal{X}$ est représentable par des
espaces algébriques, qu'il satisfait à (RS*) et qu'il commute aux
limites (voir les lemmes utilisés ci-dessus).
Il nous suffit donc de voir que $\mathcal{X}$ satisfait à l'effectivité
formelle forte formulée dans
Axiomes d'Artin, lemme \ref{artin-lemma-SGE-implies-openness-versality}.
Cela résulte du lemme \ref{lemma-complexes-strong-effectiveness}, et la
démonstration est achevée.
\end{proof}











\begin{multicols}{2}[\section{Autres chapitres}]
\noindent
Préliminaires
\begin{enumerate}
\item \hyperref[introduction-section-phantom]{Introduction}
\item \hyperref[conventions-section-phantom]{Conventions}
\item \hyperref[sets-section-phantom]{Théorie des ensembles}
\item \hyperref[categories-section-phantom]{Catégories}
\item \hyperref[topology-section-phantom]{Topologie}
\item \hyperref[sheaves-section-phantom]{Faisceaux sur les espaces}
\item \hyperref[sites-section-phantom]{Sites et faisceaux}
\item \hyperref[stacks-section-phantom]{Champs}
\item \hyperref[fields-section-phantom]{Corps}
\item \hyperref[algebra-section-phantom]{Algèbre commutative}
\item \hyperref[brauer-section-phantom]{Groupes de Brauer}
\item \hyperref[homology-section-phantom]{Algèbre homologique}
\item \hyperref[derived-section-phantom]{Catégories dérivées}
\item \hyperref[simplicial-section-phantom]{Méthodes simpliciales}
\item \hyperref[more-algebra-section-phantom]{Compléments d'algèbre}
\item \hyperref[smoothing-section-phantom]{Lissage des morphismes d'anneaux}
\item \hyperref[modules-section-phantom]{Faisceaux de modules}
\item \hyperref[sites-modules-section-phantom]{Modules sur les sites}
\item \hyperref[injectives-section-phantom]{Objets injectifs}
\item \hyperref[cohomology-section-phantom]{Cohomologie des faisceaux}
\item \hyperref[sites-cohomology-section-phantom]{Cohomologie sur les sites}
\item \hyperref[dga-section-phantom]{Algèbre différentielle graduée}
\item \hyperref[dpa-section-phantom]{Algèbre à puissances divisées}
\item \hyperref[sdga-section-phantom]{Faisceaux différentiels gradués}
\item \hyperref[hypercovering-section-phantom]{Hyperrecouvrements}
\end{enumerate}
Schémas
\begin{enumerate}
\setcounter{enumi}{25}
\item \hyperref[schemes-section-phantom]{Schémas}
\item \hyperref[constructions-section-phantom]{Constructions de schémas}
\item \hyperref[properties-section-phantom]{Propriétés des schémas}
\item \hyperref[morphisms-section-phantom]{Morphismes de schémas}
\item \hyperref[coherent-section-phantom]{Cohomologie des schémas}
\item \hyperref[divisors-section-phantom]{Diviseurs}
\item \hyperref[limits-section-phantom]{Limites de schémas}
\item \hyperref[varieties-section-phantom]{Variétés}
\item \hyperref[topologies-section-phantom]{Topologies sur les schémas}
\item \hyperref[descent-section-phantom]{Descente}
\item \hyperref[perfect-section-phantom]{Catégories dérivées de schémas}
\item \hyperref[more-morphisms-section-phantom]{Compléments sur les morphismes}
\item \hyperref[flat-section-phantom]{Compléments sur la platitude}
\item \hyperref[groupoids-section-phantom]{Schémas en groupoïdes}
\item \hyperref[more-groupoids-section-phantom]{Compléments sur les schémas en groupoïdes}
\item \hyperref[etale-section-phantom]{Morphismes étales de schémas}
\end{enumerate}
Thèmes de théorie des schémas
\begin{enumerate}
\setcounter{enumi}{41}
\item \hyperref[chow-section-phantom]{Homologie de Chow}
\item \hyperref[intersection-section-phantom]{Théorie de l'intersection}
\item \hyperref[pic-section-phantom]{Schémas de Picard des courbes}
\item \hyperref[weil-section-phantom]{Théories cohomologiques de Weil}
\item \hyperref[adequate-section-phantom]{Modules adéquats}
\item \hyperref[dualizing-section-phantom]{Complexes dualisants}
\item \hyperref[duality-section-phantom]{Dualité pour les schémas}
\item \hyperref[discriminant-section-phantom]{Discriminants et différentes}
\item \hyperref[derham-section-phantom]{Cohomologie de de Rham}
\item \hyperref[local-cohomology-section-phantom]{Cohomologie locale}
\item \hyperref[algebraization-section-phantom]{Géométrie algébrique et formelle}
\item \hyperref[curves-section-phantom]{Courbes algébriques}
\item \hyperref[resolve-section-phantom]{Résolution des surfaces}
\item \hyperref[models-section-phantom]{Réduction semi-stable}
\item \hyperref[functors-section-phantom]{Foncteurs et morphismes}
\item \hyperref[equiv-section-phantom]{Catégories dérivées de variétés}
\item \hyperref[pione-section-phantom]{Groupes fondamentaux de schémas}
\item \hyperref[etale-cohomology-section-phantom]{Cohomologie étale}
\item \hyperref[crystalline-section-phantom]{Cohomologie cristalline}
\item \hyperref[proetale-section-phantom]{Cohomologie pro-étale}
\item \hyperref[relative-cycles-section-phantom]{Cycles relatifs}
\item \hyperref[more-etale-section-phantom]{Compléments de cohomologie étale}
\item \hyperref[trace-section-phantom]{La formule des traces}
\end{enumerate}
Espaces algébriques
\begin{enumerate}
\setcounter{enumi}{64}
\item \hyperref[spaces-section-phantom]{Espaces algébriques}
\item \hyperref[spaces-properties-section-phantom]{Propriétés des espaces algébriques}
\item \hyperref[spaces-morphisms-section-phantom]{Morphismes d'espaces algébriques}
\item \hyperref[decent-spaces-section-phantom]{Espaces algébriques décents}
\item \hyperref[spaces-cohomology-section-phantom]{Cohomologie des espaces algébriques}
\item \hyperref[spaces-limits-section-phantom]{Limites d'espaces algébriques}
\item \hyperref[spaces-divisors-section-phantom]{Diviseurs sur les espaces algébriques}
\item \hyperref[spaces-over-fields-section-phantom]{Espaces algébriques sur des corps}
\item \hyperref[spaces-topologies-section-phantom]{Topologies sur les espaces algébriques}
\item \hyperref[spaces-descent-section-phantom]{Descente et espaces algébriques}
\item \hyperref[spaces-perfect-section-phantom]{Catégories dérivées d'espaces}
\item \hyperref[spaces-more-morphisms-section-phantom]{Compléments sur les morphismes d'espaces}
\item \hyperref[spaces-flat-section-phantom]{Platitude sur les espaces algébriques}
\item \hyperref[spaces-groupoids-section-phantom]{Groupoïdes dans les espaces algébriques}
\item \hyperref[spaces-more-groupoids-section-phantom]{Compléments sur les groupoïdes dans les espaces}
\item \hyperref[bootstrap-section-phantom]{Amorçage}
\item \hyperref[spaces-pushouts-section-phantom]{Sommes amalgamées d'espaces algébriques}
\end{enumerate}
Thèmes de géométrie
\begin{enumerate}
\setcounter{enumi}{81}
\item \hyperref[spaces-chow-section-phantom]{Groupes de Chow des espaces}
\item \hyperref[groupoids-quotients-section-phantom]{Quotients de groupoïdes}
\item \hyperref[spaces-more-cohomology-section-phantom]{Compléments sur la cohomologie des espaces}
\item \hyperref[spaces-simplicial-section-phantom]{Espaces simpliciaux}
\item \hyperref[spaces-duality-section-phantom]{Dualité pour les espaces}
\item \hyperref[formal-spaces-section-phantom]{Espaces algébriques formels}
\item \hyperref[restricted-section-phantom]{Algébrisation des espaces formels}
\item \hyperref[spaces-resolve-section-phantom]{Résolution des surfaces revisitée}
\end{enumerate}
Théorie des déformations
\begin{enumerate}
\setcounter{enumi}{89}
\item \hyperref[formal-defos-section-phantom]{Théorie formelle des déformations}
\item \hyperref[defos-section-phantom]{Théorie des déformations}
\item \hyperref[cotangent-section-phantom]{Le complexe cotangent}
\item \hyperref[examples-defos-section-phantom]{Problèmes de déformation}
\end{enumerate}
Champs algébriques
\begin{enumerate}
\setcounter{enumi}{93}
\item \hyperref[algebraic-section-phantom]{Champs algébriques}
\item \hyperref[examples-stacks-section-phantom]{Exemples de champs}
\item \hyperref[stacks-sheaves-section-phantom]{Faisceaux sur les champs algébriques}
\item \hyperref[criteria-section-phantom]{Critères de représentabilité}
\item \hyperref[artin-section-phantom]{Axiomes d'Artin}
\item \hyperref[quot-section-phantom]{Espaces de Quot et de Hilbert}
\item \hyperref[stacks-properties-section-phantom]{Propriétés des champs algébriques}
\item \hyperref[stacks-morphisms-section-phantom]{Morphismes de champs algébriques}
\item \hyperref[stacks-limits-section-phantom]{Limites de champs algébriques}
\item \hyperref[stacks-cohomology-section-phantom]{Cohomologie des champs algébriques}
\item \hyperref[stacks-perfect-section-phantom]{Catégories dérivées de champs}
\item \hyperref[stacks-introduction-section-phantom]{Introduction aux champs algébriques}
\item \hyperref[stacks-more-morphisms-section-phantom]{Compléments sur les morphismes de champs}
\item \hyperref[stacks-geometry-section-phantom]{La géométrie des champs}
\end{enumerate}
Thèmes de théorie des espaces de modules
\begin{enumerate}
\setcounter{enumi}{107}
\item \hyperref[moduli-section-phantom]{Champs de modules}
\item \hyperref[moduli-curves-section-phantom]{Espaces de modules de courbes}
\end{enumerate}
Divers
\begin{enumerate}
\setcounter{enumi}{109}
\item \hyperref[examples-section-phantom]{Exemples}
\item \hyperref[exercises-section-phantom]{Exercices}
\item \hyperref[guide-section-phantom]{Guide bibliographique}
\item \hyperref[desirables-section-phantom]{Desiderata}
\item \hyperref[coding-section-phantom]{Style de programmation}
\item \hyperref[obsolete-section-phantom]{Obsolète}
\item \hyperref[fdl-section-phantom]{Licence de documentation libre GNU}
\item \hyperref[index-section-phantom]{Index généré automatiquement}
\end{enumerate}
\end{multicols}


\bibliography{my}
\bibliographystyle{amsalpha}

\end{document}
